% El consumo responsable y la ética profesional — Espagnol 1ère A — Cours signé OnBuch+
% Programme officiel MINESEC. Compiler avec : tectonic EA15-consumo-responsable-etica-profesional.tex
\newcommand{\DOCMATIERE}{Espagnol}
\newcommand{\DOCNIVEAU}{1ère A}
\newcommand{\DOCMODULE}{Módulo 4 : Vie économique}
\newcommand{\DOCLECON}{EA15}
\newcommand{\DOCTITRE}{El consumo responsable y la ética profesional}
% OnBuch+ — préambule « cours signés » ENRICHI (compilé avec tectonic / XeLaTeX)
% Système visuel « Pop » d'OnBuch+ : fond crème (jamais blanc), bordures encre
% épaisses, ombres « dures » décalées, titres Archivo Black en capitales.
% Version MATHÉMATIQUES : graphes (pgfplots/tikz), boîtes pédagogiques
% (définition, propriété, méthode, attention, le savais-tu, exercice résolu,
% exercice, TP, bilan), page de garde et signature OnBuch+ ancrée en bas.
%
% Macros attendues dans main.tex AVANT \input{preamble} :
%   \DOCMATIERE  \DOCNIVEAU  \DOCMODULE  \DOCLECON (numéro)  \DOCTITRE
\documentclass[11pt]{article}

\usepackage{fontspec}
\usepackage[spanish,french]{babel} % français = langue principale ; l'espagnol via \esp{..} / espagnol
\usepackage[a4paper,margin=2cm,top=3.4cm,bottom=2.8cm,headheight=1.4cm,footskip=1.3cm]{geometry}
\usepackage{amsmath,amssymb,amsfonts}
\usepackage{mathtools} % \xmapsto, \coloneqq... souvent utilisés par les modèles
\usepackage{cancel} % simplifications barrées (bilans de forces)
% \abs{z} (module, valeur absolue) et \norm{u} (norme), utilisés par les modèles
\providecommand{\abs}{}\let\abs\relax\DeclarePairedDelimiter{\abs}{\lvert}{\rvert}
\providecommand{\norm}{}\let\norm\relax\DeclarePairedDelimiter{\norm}{\lVert}{\rVert}
\usepackage[table]{xcolor} % [table] : fournit aussi \rowcolors (alternance de lignes)
\usepackage{pagecolor}
\usepackage{tikz}
\usetikzlibrary{arrows.meta,positioning,calc,shapes.geometric,shapes.misc,shapes.arrows,decorations.pathmorphing,decorations.pathreplacing,decorations.markings,patterns,shadows,mindmap,trees,fit,backgrounds,matrix,intersections,angles,quotes}
\usepackage{pgfplots}
\usepackage[version=4]{mhchem} % équations nucléaires \ce{^{235}_{92}U}
\usepackage[european,siunitx]{circuitikz} % schémas électriques
\pgfplotsset{compat=1.18}
\usepgfplotslibrary{fillbetween,groupplots} % aires entre courbes (calcul intégral)
\usepackage{enumitem}
\usepackage{fancyhdr}
% [most] charge toutes les bibliothèques tcolorbox (listings, minted, raster,
% documentation...) et gonfle inutilement la mémoire TeX sur un document long
% (~300 boîtes) : on ne charge que ce qui est réellement utilisé.
\usepackage[skins,breakable]{tcolorbox}
\usepackage{graphicx}
\usepackage{colortbl}
\usepackage{booktabs}
\usepackage{tabularx}
\usepackage{diagbox} % case d'en-tête diagonale des tableaux à double entrée
% Colonnes extensibles centrée / gauche / droite (C, L, R), souvent utilisées par les modèles
\newcolumntype{C}{>{\centering\arraybackslash}X}
\newcolumntype{L}{>{\raggedright\arraybackslash}X}
\newcolumntype{R}{>{\raggedleft\arraybackslash}X}
\usepackage{array}
\usepackage{multicol}
\usepackage{siunitx}
% parse-numbers=false : accepte aussi \SI{2,5 \times 10^{-3}}{...}
\sisetup{locale=FR,output-decimal-marker={,},group-minimum-digits=4,parse-numbers=false}
\DeclareSIUnit\ohm{\ensuremath{\symup{\Omega}}} % U+2126 absent de la police
\DeclareSIUnit\division{div} % réglages d'oscilloscope (ms/div, V/div)
\DeclareSIUnit\tr{tr} % tours (vitesse de rotation, ex. \SI{1500}{\tr\per\minute})
\DeclareSIUnit\molar{M} % concentration molaire (ex. \SI{3}{\molar}, \SI{1}{\micro\molar})
\DeclareSIUnit\an{an} % année (le modèle confond parfois avec \year, un compteur TeX) — ex. \SI{5}{\kilo\meter\per\an}
\DeclareSIUnit\FCFA{FCFA} % franc CFA (ex. \SI{500}{\FCFA\per\kilo\gram})
\DeclareSIUnit\degreeBrix{\textdegree Brix} % teneur en sucre d'un sirop/jus (réfractométrie)
\DeclareSIUnit\month{mois} % le modèle utilise parfois \month, un compteur TeX, comme unité de durée
\usepackage{qrcode}
% \degree : commande fréquemment utilisée par les modèles pour un angle ou une
% température en dehors de \SI{}{} (non fournie par les paquets chargés).
\providecommand{\degree}{^{\circ}}
% \qed (fin de démonstration), utilisé par les modèles sans amsthm
\providecommand{\qed}{\hfill$\square$}
% \barre : double barre des valeurs interdites dans un tableau de variations (tkz-tab)
\providecommand{\barre}{\ensuremath{\|}}
% environnement proof (amsthm non chargé) utilisé par les modèles
\ifdefined\proof\else\newenvironment{proof}[1][Démonstration]{\par\noindent\textit{#1.}\ }{\qed\par}\fi
\usepackage{lastpage}
\usepackage{hyperref}
\hypersetup{hidelinks}

% ============================================================
% Palette « Pop » OnBuch+ — mêmes hex que la classe P (plus_ui.dart)
% ============================================================
\definecolor{poporange}{HTML}{F59321}
\definecolor{popdark}{HTML}{E07A0C}
\definecolor{popcream}{HTML}{FFF6E8}
\definecolor{popink}{HTML}{1C1714}
\definecolor{poppurple}{HTML}{7A5AE0}
\definecolor{popgreen}{HTML}{2E9E6A}
\definecolor{poppink}{HTML}{E0457B}
\definecolor{popblue}{HTML}{2F7DE1}
\definecolor{popgold}{HTML}{C98A12}
\definecolor{popmuted}{HTML}{8A7F76}
\definecolor{popline}{HTML}{EADFD2}
\definecolor{popgray}{HTML}{8A7F76}
\colorlet{popgrey}{popgray}
% Alias tolérants (le modèle nomme parfois les couleurs différemment)
\colorlet{popwhite}{white}
\colorlet{popblack}{popink}
\colorlet{popred}{poppink}
\colorlet{popyellow}{popgold}
% Teintes claires (fonds de tableaux/encadrés) — le modèle nomme parfois
% les couleurs avec un suffixe L (ex. popblueL) sans les avoir définies.
\colorlet{poporangeL}{poporange!20}
\colorlet{popdarkL}{popdark!20}
\colorlet{poppurpleL}{poppurple!20}
\colorlet{popgreenL}{popgreen!20}
\colorlet{poppinkL}{poppink!20}
\colorlet{popblueL}{popblue!20}
\colorlet{popgoldL}{popgold!20}
\colorlet{popredL}{popred!20}
\colorlet{popgrayL}{popgray!20}
% Teintes claires pour fonds de tableaux / zones
\colorlet{popblueL}{popblue!10!white}
\colorlet{poporangeL}{poporange!14!white}
\colorlet{popgreenL}{popgreen!12!white}
\colorlet{poppurpleL}{poppurple!12!white}
\colorlet{poppinkL}{poppink!12!white}
\colorlet{popgoldL}{popgold!14!white}

\pagecolor{popcream}

% ============================================================
% Typographie
% ============================================================
\defaultfontfeatures{Ligatures=TeX}
\setmainfont{PlusJakartaSans-Medium.ttf}[Path=../../../fonts/,BoldFont=PlusJakartaSans-Bold.ttf]
% Maths en sans-serif (Fira Math) avec chiffres et lettres droites de Plus
% Jakarta, pour que les formules se fondent dans le texte.
\usepackage[math-style=ISO,bold-style=ISO]{unicode-math}
\setmathfont{FiraMath-Regular.otf}[Path=../../../fonts/]
\setmathfont{PlusJakartaSans-Medium.ttf}[Path=../../../fonts/,range={up/num,up/{Latin,latin},\mathup}]
\usepackage{icomma} % virgule décimale sans espace en mode math
% Caractères mathématiques Unicode absents des polices de texte (Plus Jakarta,
% Archivo Black) : rendus par leur équivalent mathématique, en texte comme en
% formule — sinon ils s'affichent en carré vide (ex. « Arithmétique dans ℤ »).
\usepackage{newunicodechar}
\newunicodechar{ℝ}{\ensuremath{\mathbb{R}}}
\newunicodechar{ℤ}{\ensuremath{\mathbb{Z}}}
\newunicodechar{ℂ}{\ensuremath{\mathbb{C}}}
\newunicodechar{ℕ}{\ensuremath{\mathbb{N}}}
\newunicodechar{ℚ}{\ensuremath{\mathbb{Q}}}
\newunicodechar{𝔻}{\ensuremath{\mathbb{D}}}
\newunicodechar{α}{\ensuremath{\alpha}}
\newunicodechar{β}{\ensuremath{\beta}}
\newunicodechar{γ}{\ensuremath{\gamma}}
\newunicodechar{δ}{\ensuremath{\delta}}
\newunicodechar{ε}{\ensuremath{\varepsilon}}
\newunicodechar{θ}{\ensuremath{\theta}}
\newunicodechar{λ}{\ensuremath{\lambda}}
\newunicodechar{μ}{\ensuremath{\mu}}
\newunicodechar{σ}{\ensuremath{\sigma}}
\newunicodechar{τ}{\ensuremath{\tau}}
\newunicodechar{φ}{\ensuremath{\varphi}}
\newunicodechar{ψ}{\ensuremath{\psi}}
\newunicodechar{ω}{\ensuremath{\omega}}
\newunicodechar{Σ}{\ensuremath{\Sigma}}
% majuscules produites par \MakeUppercase dans les titres (α-aminés -> Α)
\newunicodechar{Α}{\ensuremath{\alpha}}
\newunicodechar{Β}{\ensuremath{\beta}}
\newunicodechar{Γ}{\ensuremath{\Gamma}}
\newunicodechar{Δ}{\ensuremath{\Delta}}
\newunicodechar{Ε}{\ensuremath{\varepsilon}}
\newunicodechar{Θ}{\ensuremath{\Theta}}
\newunicodechar{Λ}{\ensuremath{\Lambda}}
\newunicodechar{Μ}{\ensuremath{\mu}}
\newunicodechar{Τ}{\ensuremath{\tau}}
\newunicodechar{Φ}{\ensuremath{\Phi}}
\newunicodechar{Ψ}{\ensuremath{\Psi}}
\newunicodechar{Ω}{\ensuremath{\Omega}}
\newunicodechar{↦}{\ensuremath{\mapsto}}
\newunicodechar{∈}{\ensuremath{\in}}
\newunicodechar{∉}{\ensuremath{\notin}}
\newunicodechar{∧}{\ensuremath{\wedge}}
\newunicodechar{⇔}{\ensuremath{\Leftrightarrow}}
\newunicodechar{⟺}{\ensuremath{\Longleftrightarrow}}
\newunicodechar{⇒}{\ensuremath{\Rightarrow}}
\newunicodechar{⟹}{\ensuremath{\Longrightarrow}}
\newunicodechar{∘}{\ensuremath{\circ}}
\newunicodechar{∪}{\ensuremath{\cup}}
\newunicodechar{∩}{\ensuremath{\cap}}
\newunicodechar{≡}{\ensuremath{\equiv}}
\newunicodechar{⊂}{\ensuremath{\subset}}
\newunicodechar{⊥}{\ensuremath{\perp}}
\newunicodechar{∃}{\ensuremath{\exists}}
\newunicodechar{∀}{\ensuremath{\forall}}
\newunicodechar{∛}{\ensuremath{\sqrt[3]{\,}}}
\newunicodechar{∜}{\ensuremath{\sqrt[4]{\,}}}
\newunicodechar{⃗}{\ensuremath{{}^{\to}}}
\newunicodechar{ⁿ}{\ifmmode^{n}\else\textsuperscript{\ensuremath{n}}\fi}
\newunicodechar{ˣ}{\ifmmode^{x}\else\textsuperscript{\ensuremath{x}}\fi}
\newunicodechar{ᵇ}{\ifmmode^{b}\else\textsuperscript{\ensuremath{b}}\fi}
\newunicodechar{ʳ}{\ifmmode^{r}\else\textsuperscript{\ensuremath{r}}\fi}
\newunicodechar{ᵘ}{\ifmmode^{u}\else\textsuperscript{\ensuremath{u}}\fi}
\newunicodechar{ʸ}{\ifmmode^{y}\else\textsuperscript{\ensuremath{y}}\fi}
\newunicodechar{ᵃ}{\ifmmode^{a}\else\textsuperscript{\ensuremath{a}}\fi}
\newunicodechar{ᵉ}{\ifmmode^{e}\else\textsuperscript{\ensuremath{e}}\fi}
\newunicodechar{ᵏ}{\ifmmode^{k}\else\textsuperscript{\ensuremath{k}}\fi}
\newunicodechar{ᵖ}{\ifmmode^{p}\else\textsuperscript{\ensuremath{p}}\fi}
\newunicodechar{ᴺ}{\ifmmode^{N}\else\textsuperscript{\ensuremath{N}}\fi}
\newunicodechar{ᶜ}{\ifmmode^{c}\else\textsuperscript{\ensuremath{c}}\fi}
\newunicodechar{ʰ}{\ifmmode^{h}\else\textsuperscript{\ensuremath{h}}\fi}
\newunicodechar{ᵠ}{\ifmmode^{q}\else\textsuperscript{\ensuremath{q}}\fi}
\newunicodechar{ₙ}{\ifmmode_{n}\else\textsubscript{\ensuremath{n}}\fi}
\newunicodechar{ₐ}{\ifmmode_{a}\else\textsubscript{\ensuremath{a}}\fi}
\newunicodechar{ᵢ}{\ifmmode_{i}\else\textsubscript{\ensuremath{i}}\fi}
\newunicodechar{ᵣ}{\ifmmode_{r}\else\textsubscript{\ensuremath{r}}\fi}
\newunicodechar{ₖ}{\ifmmode_{k}\else\textsubscript{\ensuremath{k}}\fi}
\newunicodechar{ᵦ}{\ifmmode_{\beta}\else\textsubscript{\ensuremath{\beta}}\fi}
\newunicodechar{ₓ}{\ifmmode_{x}\else\textsubscript{\ensuremath{x}}\fi}
\newfontfamily\popdisplay{ArchivoBlack-Regular.ttf}[Path=../../../fonts/]
\newfontfamily\poplogo{SpaceGrotesk-Bold.ttf}[Path=../../../fonts/]
\newfontfamily\popsemi{PlusJakartaSans-SemiBold.ttf}[Path=../../../fonts/]
\newfontfamily\popxbold{PlusJakartaSans-ExtraBold.ttf}[Path=../../../fonts/]
\newfontfamily\popxboldls{PlusJakartaSans-ExtraBold.ttf}[Path=../../../fonts/,LetterSpace=3.0]

\setlist[enumerate]{leftmargin=1.7em,itemsep=5pt,topsep=4pt}
\setlist[itemize]{leftmargin=1.7em,itemsep=4pt,topsep=4pt,label={\color{poporange}\textbullet}}
\setlength{\parindent}{0pt}
\setlength{\parskip}{5pt}
\renewcommand{\arraystretch}{1.25}

% pgfplots : style Pop par défaut
\pgfplotsset{
  every axis/.append style={axis line style={popink,line width=1pt},tick style={popink},
    grid style={popline,line width=0.5pt},label style={font=\small},tick label style={font=\footnotesize}},
  popcurve/.style={poporange,line width=1.8pt,no markers,smooth},
}
% Même style disponible dans un \draw TikZ simple (hors pgfplots)
\tikzset{popcurve/.style={poporange,line width=1.8pt}}
\tikzset{popgrid/.style={popline,very thin}}
\colorlet{popcurve}{poporange} % le nom du style est parfois utilisé comme couleur

% ============================================================
% Pill / squiggle
% ============================================================
\newcommand{\poppill}[3]{%
  \tikz[baseline=-0.55ex]\node[draw=popink,line width=1.2pt,rounded corners=2.6pt,%
    fill=#2,inner xsep=6.5pt,inner ysep=3pt]{%
    \popxboldls\fontsize{7}{7}\selectfont\color{#3}\MakeUppercase{#1}};%
}
\newcommand{\popsquiggle}[1]{%
  \tikz[baseline=-0.4ex]{
    \draw[#1,line width=2.6pt,line cap=round]
      plot [smooth] coordinates {(0,0.09) (0.34,0.01) (0.68,0.21) (1.02,0.01) (1.36,0.10)};
  }%
}
% Mot-clé surligné (marqueur orange sous le texte)
\newcommand{\cle}[1]{{\popxbold\color{popdark}#1}}

% ============================================================
% En-tête / pied
% ============================================================
\pagestyle{fancy}
\fancyhf{}
\renewcommand{\headrulewidth}{0pt}
\renewcommand{\footrulewidth}{0pt}
\lhead{\raisebox{-0.15ex}{\poplogo\fontsize{13}{13}\selectfont\color{popink}OnBuch\color{poporange}+}}
\rhead{\poppill{\DOCMATIERE\ \textbullet\ \DOCNIVEAU\ \textbullet\ Leçon \DOCLECON}{white}{popink}}
\lfoot{\popxbold\fontsize{7.6}{7.6}\selectfont\color{popmuted}\MakeUppercase{Cours signé OnBuch+}\ \textbullet\ {\popsemi\DOCTITRE}}
\rfoot{\tikz[baseline=-0.6ex]\node[rounded corners=8.5pt,fill=popink,minimum height=17pt,minimum width=17pt,inner xsep=5pt,inner sep=0]{\popxbold\fontsize{8}{8}\selectfont\color{popcream}\thepage\,/\,\pageref*{LastPage}};}

% ============================================================
% Titres
% ============================================================
\newcommand{\chaptitre}[2]{%
  \begin{tcolorbox}[enhanced,colback=poporange,colframe=popink,boxrule=2.6pt,arc=11pt,
      drop shadow=popink,left=15pt,right=15pt,top=12pt,bottom=11pt,before skip=3pt,after skip=18pt]
    \poppill{#2}{popink}{popcream}\par\vspace{9pt}
    {\raggedright\hyphenpenalty=10000\exhyphenpenalty=10000\popdisplay\fontsize{22}{24}\selectfont\color{popcream}\MakeUppercase{#1}\par}
  \end{tcolorbox}
}
% Section numérotée : pastille orange avec le numéro + titre Archivo
\newcounter{popsec}
\newcommand{\coursec}[1]{%
  \stepcounter{popsec}\setcounter{popsub}{0}%
  \par\vspace{16pt}\noindent
  \begin{minipage}{\linewidth}
  \tikz[baseline=-0.7ex]\node[draw=popink,line width=1.6pt,fill=poporange,rounded corners=4pt,minimum size=22pt,inner sep=0]{\popdisplay\fontsize{12}{12}\selectfont\color{popcream}\thepopsec};%
  \hspace{8pt}{\popdisplay\fontsize{15}{17}\selectfont\color{popink}\MakeUppercase{#1}}\par
  \vspace{4pt}\noindent{\color{poporange}\rule{36pt}{3.6pt}}
  \end{minipage}\par\nopagebreak\vspace{8pt}
}
\newcounter{popsub}
\newcommand{\courssub}[1]{%
  \stepcounter{popsub}%
  \par\vspace{9pt}\noindent{\popxbold\fontsize{11.5}{13}\selectfont\color{popink}{\color{poporange}\thepopsec.\thepopsub}\hspace{6pt}#1}\par\nopagebreak\vspace{4pt}
}

% ============================================================
% Boîtes pédagogiques — même recette : fond blanc, bordure épaisse colorée,
% ombre dure encre, badge plein. Titre optionnel [#1].
% ============================================================
\tcbset{popbase/.style={breakable,enhanced,colback=white,boxrule=2.3pt,arc=9pt,
  shadow={3pt}{-3pt}{0pt}{popink},left=14pt,right=14pt,top=11pt,bottom=11pt,before skip=11pt,after skip=15pt,
  fontupper=\popsemi\fontsize{10.6}{14.5}\selectfont\color{popink},
  fontlower=\popsemi\fontsize{10.6}{14.5}\selectfont\color{popink}}}
% Coche dessinée (le glyphe ✓ n'existe pas dans les polices du document)
\newcommand{\popcheck}[1]{\tikz[baseline=-0.5ex]\draw[#1,line width=1.6pt,line cap=round,line join=round](-0.13,0.0)--(-0.04,-0.09)--(0.13,0.1);}
\providecommand{\coche}{\popcheck{popgreen}}
\providecommand{\resultat}[1]{\par\smallskip\noindent\fcolorbox{popgreen}{popgreenL}{\parbox{\dimexpr\linewidth-2\fboxsep-2\fboxrule\relax}{\textbf{Résultat.} #1}}\par\smallskip}
\newcommand{\popboxhead}[3]{\poppill{#1}{#2}{white}\ifx\relax#3\relax\else\hspace{6pt}{\popxbold\fontsize{10.6}{12}\selectfont\color{popink}#3}\fi\par\vspace{7pt}}

\newtcolorbox{aretenir}[1][]{popbase,colframe=popgreen,before upper={\popboxhead{À retenir}{popgreen}{#1}}}
% Texte en espagnol (document de lecture, dialogue, poème, extrait) : cadre
% d'étude. Un titre optionnel (source, auteur) s'affiche dans l'en-tête.
\newtcolorbox{textebox}[1][]{popbase,colframe=popblue,before upper={\popboxhead{Texto}{popblue}{#1}\begin{otherlanguage}{spanish}},after upper={\end{otherlanguage}}}
% Marques ✓ / ✗ (forme correcte / fautive) souvent utilisées par les modèles
\providecommand{\cross}{\textcolor{poppink}{\ensuremath{\times}}}
\providecommand{\xmark}{\cross}\providecommand{\crossmark}{\cross}\providecommand{\cmark}{\ensuremath{\checkmark}}
% Ligne de réponse / justification à compléter (inventée par les modèles)
\providecommand{\justification}[1]{\par\noindent\textit{Justification~:} \dotfill\par}
% \textipa (package tipa non chargé) : la police ne couvre pas l'API -> simple passage
\providecommand{\textipa}[1]{#1}
% Soulignement (ulem non chargé) : \uline, \ul
\providecommand{\uline}[1]{\underline{#1}}\providecommand{\ul}[1]{\underline{#1}}
% Ordinaux français écrits en macro par les modèles : 1\ière, 3\ième
\newcommand{\ière}{\textsuperscript{re}}\newcommand{\ième}{\textsuperscript{e}}
% \exemple (inventée par les modèles) : étiquette « Exemple : »
\providecommand{\exemple}{\textbf{Exemple~:}\ }
% Mot ou phrase en espagnol à l'intérieur d'un paragraphe français
\newcommand{\esp}[1]{\ifmmode\text{\textit{#1}}\else\begingroup\selectlanguage{spanish}\itshape #1\endgroup\fi}
\newtcolorbox{exemplebox}[1][]{popbase,colframe=poporange,before upper={\popboxhead{Exemple}{poporange}{#1}}}
\newtcolorbox{definition}[1][]{popbase,colframe=popblue,before upper={\popboxhead{Définition}{popblue}{#1}}}
\newtcolorbox{propriete}[1][]{popbase,colframe=poppurple,before upper={\popboxhead{Propriété}{poppurple}{#1}}}
\newtcolorbox{methode}[1][]{popbase,colframe=popgold,before upper={\popboxhead{Méthode}{popgold}{#1}}}
\newtcolorbox{attention}[1][]{popbase,colframe=poppink,before upper={\popboxhead{Attention}{poppink}{#1}}}
\newtcolorbox{experience}[1][]{popbase,colframe=popblue,colback=popblueL,before upper={\popboxhead{Expérience}{popblue}{#1}}}
% « Le savais-tu ? » : carte violette pleine (contexte, culture, Cameroun)
\newtcolorbox{savaistu}[1][]{popbase,colback=poppurple,colframe=popink,
  fontupper=\popsemi\fontsize{10.4}{14.2}\selectfont\color{white},
  before upper={\poppill{Le savais-tu\,?}{white}{poppurple}\ifx\relax#1\relax\else\hspace{6pt}{\popxbold\color{white}#1}\fi\par\vspace{7pt}}}
% Exercice résolu : énoncé en haut, \tcblower puis la solution sur fond crème
\newcounter{popexo}\newcounter{popexr}
\newtcolorbox{exoresolu}[1][]{popbase,colframe=poporange,colbacklower=popcream,
  segmentation style={popink,line width=1.2pt,dashed},
  before upper={\stepcounter{popexr}\popboxhead{Exercice résolu \thepopexr}{poporange}{#1}},
  before lower={\poppill{Solution}{popink}{popcream}\par\vspace{6pt}}}
% Exercice (non résolu en place) — la correction est dans la partie Corrigés
\newtcolorbox{exercice}[1][]{popbase,colframe=popink,
  before upper={\stepcounter{popexo}\popboxhead{Exercice \thepopexo}{popink}{#1}}}
% Corrigé
\newtcolorbox{corrige}[1][]{popbase,colframe=popgreen,colback=popgreenL,
  before upper={\popboxhead{Corrigé}{popgreen}{#1}}}
% Objectifs / prérequis / bilan
\newtcolorbox{objectifs}{popbase,colframe=popink,colback=poporangeL,
  before upper={\popboxhead{Objectifs de la leçon}{popink}{}}}
\newtcolorbox{prerequis}{popbase,colframe=popink,colback=popblueL,
  before upper={\popboxhead{Prérequis}{popblue}{}}}
\newtcolorbox{bilan}{popbase,colframe=popink,colback=white,boxrule=2.8pt,
  before upper={\popboxhead{Fiche bilan}{poporange}{L'essentiel à connaître pour l'examen}}}
% Figure encadrée avec légende
\newtcolorbox{popfigure}[1][]{enhanced,breakable=false,colback=white,colframe=popline,boxrule=1.4pt,arc=8pt,
  left=10pt,right=10pt,top=10pt,bottom=8pt,before skip=10pt,after skip=12pt,halign=center,
  lower separated=false,halign lower=center,
  fontlower=\popsemi\fontsize{9}{11.5}\selectfont\color{popmuted},
  #1}
\newcommand{\legende}[1]{\par\vspace{6pt}{\popsemi\fontsize{9}{11.5}\selectfont\color{popmuted}#1}}

% ============================================================
% Signature OnBuch+
% ============================================================
\newcommand{\poplogoinline}[1]{{\poplogo\fontsize{#1}{#1}\selectfont\color{popink}OnBuch\color{poporange}+}}
\newcommand{\signatureonbuch}{%
  \begin{tcolorbox}[enhanced,colback=popink,colframe=popink,boxrule=2.2pt,arc=10pt,
      drop shadow=poporange,left=14pt,right=14pt,top=10pt,bottom=10pt,before skip=0pt,after skip=0pt]
    \tikz[baseline=-0.7ex]\node[circle,fill=popgreen,draw=popcream,line width=1.3pt,minimum size=22pt,inner sep=0]{\popcheck{white}};%
    \hspace{9pt}\begin{minipage}[c]{0.62\linewidth}
      {\popxbold\fontsize{12}{13}\selectfont\color{popcream}Signé\ \textbullet\ {\poplogo OnBuch\color{poporange}+}}\par\vspace{2pt}
      {\raggedright\popsemi\fontsize{8}{10.5}\selectfont\color{popline}\MakeUppercase{Équipe pédagogique OnBuch+}\\\MakeUppercase{Conforme au programme officiel MINESEC}\par}
    \end{minipage}\hfill
    {\poplogo\fontsize{22}{22}\selectfont\color{popcream}OnBuch\color{poporange}+}
  \end{tcolorbox}%
}

% ============================================================
% Page de garde
% #1 titre · #2 matière · #3 niveau · #4 module · #5 n° leçon · #6 durée ·
% #7 description / objectifs (texte ou itemize)
% ============================================================
\newcommand{\pagegarde}[7]{%
  \thispagestyle{empty}
  \newgeometry{a4paper,margin=1.7cm,top=1.6cm,bottom=1.4cm}%
  \begin{tikzpicture}[remember picture,overlay]
    \fill[poporange!18] ([xshift=-3.2cm,yshift=-3.6cm]current page.north east) circle (4.6cm);
    \fill[poppurple!14] ([xshift=2.4cm,yshift=7.6cm]current page.south west) circle (3.8cm);
  \end{tikzpicture}%
  {\poplogo\fontsize{30}{30}\selectfont\color{popink}OnBuch\color{poporange}+}\hfill
  \raisebox{0.6ex}{\poppill{Cours signé \textbullet\ Premium}{popink}{popcream}}\par
  \vspace{1pt}\popsquiggle{poppurple}\par
  \vspace{8pt}
  \begin{tcolorbox}[enhanced,colback=poporange,colframe=popink,boxrule=2.8pt,arc=13pt,
      drop shadow=popink,left=18pt,right=18pt,top=12pt,bottom=13pt,after skip=10pt]
    \poppill{#2 \textbullet\ #3}{popink}{popcream}\hspace{5pt}\poppill{#4}{white}{popink}\par\vspace{8pt}
    {\popdisplay\fontsize{14}{14}\selectfont\color{popink}LEÇON #5}\par\vspace{4pt}
    {\raggedright\hyphenpenalty=10000\exhyphenpenalty=10000\popdisplay\fontsize{25}{27}\selectfont\color{popcream}\MakeUppercase{#1}\par}
    \vspace{8pt}
    {\popxbold\fontsize{9.5}{9.5}\selectfont\color{popink}\MakeUppercase{Durée conseillée : #6}}
  \end{tcolorbox}
  \begin{tcolorbox}[enhanced,colback=white,colframe=popink,boxrule=2.2pt,arc=10pt,
      drop shadow=popink,left=16pt,right=16pt,top=10pt,bottom=10pt,after skip=8pt]
    \poppill{Dans cette leçon}{poppurple}{white}\par\vspace{5pt}
    \popsemi\fontsize{9.3}{12.6}\selectfont\color{popink}#7
  \end{tcolorbox}
  \noindent\begin{minipage}[t]{0.487\linewidth}
    \begin{tcolorbox}[enhanced,colback=popgreen,colframe=popink,boxrule=2.2pt,arc=10pt,
        drop shadow=popink,left=11pt,right=11pt,top=10pt,bottom=10pt,height=3.1cm,valign=center]
      \tcbox[colback=white,colframe=popink,boxrule=1.3pt,arc=5pt,boxsep=0pt,nobeforeafter,
        left=3pt,right=3pt,top=3pt,bottom=3pt,valign=center]{\qrcode[height=1.9cm]{https://play.google.com/store/apps/details?id=com.luvvix.onbuch&hl=fr}}%
      \hspace{8pt}\begin{minipage}[c]{0.5\linewidth}
        {\popdisplay\fontsize{11}{12.5}\selectfont\color{white}\MakeUppercase{Télécharge OnBuch}}\par\vspace{4pt}
        {\raggedright\popsemi\fontsize{8.4}{11}\selectfont\color{white}Gratuit \textbullet\ Google Play\\Tuteur IA, annales, concours, résultats\par}
      \end{minipage}
    \end{tcolorbox}
  \end{minipage}\hfill
  \begin{minipage}[t]{0.487\linewidth}
    \begin{tcolorbox}[enhanced,colback=poppurple,colframe=popink,boxrule=2.2pt,arc=10pt,
        drop shadow=popink,left=11pt,right=11pt,top=10pt,bottom=10pt,height=3.1cm,valign=center]
      \tikz[baseline=-0.7ex]\node[circle,fill=white,draw=popink,line width=1.3pt,minimum size=1.6cm,inner sep=0]{\popdisplay\fontsize{24}{24}\selectfont\color{poppurple}+};%
      \hspace{8pt}\begin{minipage}[c]{0.6\linewidth}
        {\popdisplay\fontsize{11}{12.5}\selectfont\color{white}\MakeUppercase{Passe à OnBuch+}}\par\vspace{4pt}
        {\raggedright\popsemi\fontsize{8.4}{11}\selectfont\color{white}Tous les cours signés, Tuteur IA illimité, fiches PDF — onglet OnBuch+ de l'app\par}
      \end{minipage}
    \end{tcolorbox}
  \end{minipage}
  \vspace{8pt}
  \popcontenu
  \vfill
  \signatureonbuch
  \restoregeometry
  \newpage
}

% Rangée « Ce cours contient » (page de garde)
\newcommand{\poptile}[3]{% #1 couleur #2 gros texte #3 légende
  \begin{tcolorbox}[enhanced,colback=#1,colframe=popink,boxrule=2pt,arc=9pt,shadow={2.5pt}{-2.5pt}{0pt}{popink},
      left=6pt,right=6pt,top=9pt,bottom=9pt,halign=center,valign=center,height=1.75cm,width=\linewidth,nobeforeafter]
    {\popdisplay\fontsize{12.5}{13}\selectfont\color{white}\MakeUppercase{#2}}\par\vspace{3pt}
    {\centering\popsemi\fontsize{7.8}{9.5}\selectfont\color{white}#3\par}
  \end{tcolorbox}}
\newcommand{\popcontenu}{%
  \par\noindent{\popxboldls\fontsize{7.5}{7.5}\selectfont\color{popmuted}CE COURS SIGNÉ CONTIENT}\par\vspace{7pt}
  \noindent\begin{minipage}[t]{0.235\linewidth}\poptile{popblue}{Cours}{complet, illustré, pas à pas}\end{minipage}\hfill
  \begin{minipage}[t]{0.235\linewidth}\poptile{poporange}{Méthodes}{et exercices résolus}\end{minipage}\hfill
  \begin{minipage}[t]{0.235\linewidth}\poptile{poppink}{Exercices}{type examen + corrigés}\end{minipage}\hfill
  \begin{minipage}[t]{0.235\linewidth}\poptile{popgreen}{Fiche bilan}{l'essentiel à retenir}\end{minipage}\par}

% Sommaire
\newtcolorbox{sommaire}{enhanced,colback=white,colframe=popink,boxrule=2.2pt,arc=10pt,
  drop shadow=popink,left=18pt,right=18pt,top=13pt,bottom=13pt,before skip=6pt,after skip=20pt,
  before upper={\poppill{Sommaire}{poppurple}{white}\par\vspace{9pt}\popsemi\fontsize{10.6}{14.5}\selectfont\color{popink}}}

% Signature de fin de cours — poussée tout en bas de la dernière page
\newcommand{\findecours}{%
  \par\vfill\nopagebreak
  \begin{center}\popsquiggle{poporange}\end{center}\vspace{4pt}
  \signatureonbuch
}
\AtBeginDocument{\def\llbracket{\mathopen{[\mkern-3.5mu[}}\def\rrbracket{\mathclose{]\mkern-3.5mu]}}} % intervalles d'entiers (glyphe ⟦ absent de la police)
% Glyphes absents de Fira Math : redéfinis à partir de glyphes disponibles
\AtBeginDocument{%
  \def\setminus{\mathbin{\backslash}}\let\smallsetminus\setminus
  \def\vdots{\mathord{\vbox{\baselineskip3.2pt\lineskiplimit0pt\kern4pt\hbox{.}\hbox{.}\hbox{.}}}}%
  \def\ddots{\mathinner{\mkern1mu\raise6.5pt\hbox{.}\mkern2mu\raise3.6pt\hbox{.}\mkern2mu\raise0.7pt\hbox{.}\mkern1mu}}%
  \def\triangle{\mathord{\scalebox{0.9}{$\symup{\Delta}$}}}%
  \def\nabla{\mathord{\raisebox{\depth}{\scalebox{1}[-1]{$\symup{\Delta}$}}}}%
  \def\checkmark{\popcheck{popdark}}%
  \def\star{\mathbin{\ast}}\def\bigstar{\mathord{\ast}}%
  \def\diamond{\mathbin{\scalebox{0.6}{$\Diamond$}}}%
  \def\bigcup{\mathop{\mathchoice{\raisebox{-0.3em}{\scalebox{1.7}{$\cup$}}}{\scalebox{1.25}{$\cup$}}{\cup}{\cup}}\displaylimits}%
  \def\bigcap{\mathop{\mathchoice{\raisebox{-0.3em}{\scalebox{1.7}{$\cap$}}}{\scalebox{1.25}{$\cap$}}{\cap}{\cap}}\displaylimits}%
  \def\lesseqqgtr{\mathrel{\lessgtr}}\def\bigoplus{\mathop{\oplus}\displaylimits}%
}
\providecommand{\ssi}{si et seulement si} % abréviation fréquente
\AtBeginDocument{\providecommand{\wideparen}{\overparen}} % arc de cercle

%% FIGFIX-BEGIN v5 — correctifs globaux des figures (docs/onbuch-generation/tools/figfix)
\usepackage{adjustbox}\usepackage{etoolbox}\tcbuselibrary{raster}
% toute figure TikZ/pgfplots plus large que la ligne est réduite à la largeur disponible
\BeforeBeginEnvironment{tikzpicture}{\begin{adjustbox}{max width=\linewidth}}
\AfterEndEnvironment{tikzpicture}{\end{adjustbox}}
% légendes pgfplots lisibles par-dessus les courbes
\pgfplotsset{every axis/.append style={legend style={fill=white,fill opacity=0.92,text opacity=1,draw=black!15,inner sep=3pt}}}
% étiquettes d'arêtes (syntaxe edge["..."]) avec fond blanc
\tikzset{every edge quotes/.append style={fill=white,inner sep=1.2pt}}
%% FIGFIX-END
\begin{document}

\pagegarde{El consumo responsable y la ética profesional}{Espagnol}{1ère A}{Módulo 4 : Vie économique}{EA15}{2 h}{Cette leçon propose la lecture et le commentaire d'un texte sur la consommation responsable, enrichi d'un lexique thématique complet (consommation, commerce équitable, éthique professionnelle). Elle permet aussi d'exprimer l'obligation et l'interdiction en espagnol (deber, tener que, no se debe) pour formuler des règles de conduite professionnelle.
\vspace{4pt}
\begin{multicols}{2}\small
\begin{itemize}
\item À la fin de cette leçon, je sais lire et comprendre un texte espagnol sur la consommation responsable.
\item À la fin de cette leçon, je sais distinguer la consommation responsable du gaspillage (derroche).
\item À la fin de cette leçon, je sais utiliser le lexique de la consommation : consumo, consumidor, publicidad, ahorro, producto local, comercio justo.
\item À la fin de cette leçon, je sais exprimer l'obligation avec deber + infinitif et tener que + infinitif.
\item À la fin de cette leçon, je sais exprimer l'interdiction et le conseil général avec no se debe + infinitif.
\item À la fin de cette leçon, je sais énoncer des règles d'éthique professionnelle : honnêteté, responsabilité, ponctualité, respect du client.
\end{itemize}\end{multicols}}

\begin{sommaire}
\begin{multicols}{2}
\begin{enumerate}
\item Texte support : « Comprar menos, comprar mejor »
\item Compréhension du texte
\item Lexique thématique par champs
\item Grammaire outillée : obligation et interdiction
\item Méthode du commentaire de texte
\item Production guidée
\item Activité d'intégration
\item Méthodes et exercices résolus
\item Exercices
\item Corrigés des exercices
\item Fiche bilan
\end{enumerate}
\end{multicols}
\end{sommaire}

\begin{objectifs}
\begin{itemize}
\item À la fin de cette leçon, je sais lire et comprendre un texte espagnol sur la consommation responsable.
\item À la fin de cette leçon, je sais distinguer la consommation responsable du gaspillage (derroche).
\item À la fin de cette leçon, je sais utiliser le lexique de la consommation : consumo, consumidor, publicidad, ahorro, producto local, comercio justo.
\item À la fin de cette leçon, je sais exprimer l'obligation avec deber + infinitif et tener que + infinitif.
\item À la fin de cette leçon, je sais exprimer l'interdiction et le conseil général avec no se debe + infinitif.
\item À la fin de cette leçon, je sais énoncer des règles d'éthique professionnelle : honnêteté, responsabilité, ponctualité, respect du client.
\item À la fin de cette leçon, je sais commenter un texte argumentatif simple en suivant une méthode claire.
\item À la fin de cette leçon, je sais rédiger un court texte ou une affiche de sensibilisation à la consommation responsable.
\end{itemize}
\end{objectifs}

\begin{prerequis}
\begin{itemize}
\item Le présent de l'indicatif des verbes réguliers et des verbes usuels (deber, tener, hacer).
\item Les articles définis et indéfinis, l'accord de l'adjectif (Seconde).
\item Le vocabulaire de base de l'argent et des achats : comprar, vender, precio, dinero.
\item La construction impersonnelle simple avec se (se vende, se habla).
\item Les connecteurs logiques de base : pero, porque, además, sin embargo.
\end{itemize}
\end{prerequis}

\newpage

\coursec{Texte support : « Comprar menos, comprar mejor »}

\courssub{Situation d'accroche et problématique}

Au marché Mokolo de Yaoundé, Amina hésite : acheter des tomates importées, moins chères, ou celles des producteurs de la région, un peu plus chères mais fraîches ? Son frère, lui, vient d'ouvrir un atelier de couture et se demande comment gagner la confiance de ses clients. Ces deux questions quotidiennes — \textbf{comment consommer} et \textbf{comment travailler} — sont au cœur de deux notions essentielles du monde hispanophone : \esp{el consumo responsable} (la consommation responsable) et \esp{la ética profesional} (l'éthique professionnelle).

Voici notre problématique : \emph{Comment dire en espagnol qu'il ne faut pas gaspiller ? Comment formuler les règles d'un bon professionnel ?} Pour y répondre, nous allons d'abord lire un texte court qui raconte la vie d'une famille camerounaise au marché, puis en extraire le vocabulaire et les structures utiles.

\courssub{Lecture du texte}

\begin{methode}[Lire un texte en espagnol : les cinq étapes]
\begin{enumerate}
\item \textbf{Lire le titre} et deviner le thème : que signifie \esp{Comprar menos, comprar mejor} ?
\item \textbf{Première lecture globale}, silencieuse, \emph{sans dictionnaire} : il s'agit seulement de comprendre l'idée générale.
\item \textbf{Souligner les mots connus} et les \cle{mots transparents} (ceux qui ressemblent au français : \esp{productos}, \esp{economía}, \esp{publicidad}...).
\item \textbf{Deuxième lecture} avec le glossaire : on précise le sens des phrases difficiles.
\item \textbf{Résumer le texte en une phrase} : qui ? où ? quoi ? pourquoi ?
\end{enumerate}
\end{methode}

\begin{textebox}[Comprar menos, comprar mejor]
En el mercado de mi barrio, mi madre siempre compra productos locales: tomates, plátanos y frijoles de los campesinos de la región. Dice que así ayudamos a la economía del país y comemos más sano. Mi hermano, en cambio, se deja llevar por la publicidad y compra cosas innecesarias: es un derroche de dinero. Nosotros debemos consumir de manera responsable: tenemos que ahorrar, reutilizar los envases y respetar el medio ambiente. No se debe tirar la comida mientras otras personas tienen hambre. El comercio justo también es importante: el consumidor responsable paga un precio digno al productor. Comprar menos, pero comprar mejor: esa es la clave.
\end{textebox}

\textbf{Glossaire :}
\begin{itemize}
\item \esp{el mercado} : le marché ; \esp{el barrio} : le quartier
\item \esp{los campesinos} : les paysans, les cultivateurs
\item \esp{comer sano} : manger sainement
\item \esp{en cambio} : en revanche, par contre
\item \esp{dejarse llevar por} : se laisser entraîner par
\item \esp{el derroche} : le gaspillage
\item \esp{ahorrar} : économiser ; \esp{el ahorro} : l'économie (d'argent)
\item \esp{reutilizar} : réutiliser ; \esp{el envase} : l'emballage
\item \esp{tirar la comida} : jeter la nourriture
\item \esp{un precio digno} : un prix juste, digne
\item \esp{la clave} : la clé, le secret
\end{itemize}

\textbf{Traduction du texte :} « Au marché de mon quartier, ma mère achète toujours des produits locaux : des tomates, des bananes plantain et des haricots des paysans de la région. Elle dit qu'ainsi nous aidons l'économie du pays et que nous mangeons plus sainement. Mon frère, en revanche, se laisse entraîner par la publicité et achète des choses inutiles : c'est un gaspillage d'argent. Nous devons consommer de manière responsable : nous devons économiser, réutiliser les emballages et respecter l'environnement. On ne doit pas jeter la nourriture pendant que d'autres personnes ont faim. Le commerce équitable est important aussi : le consommateur responsable paie un prix juste au producteur. Acheter moins, mais acheter mieux : voilà la clé. »

\courssub{Compréhension du texte}

\textbf{Questions de compréhension (réponses attendues en espagnol) :}
\begin{enumerate}
\item ¿Dónde compra la madre? ¿Qué compra? \emph{(Où la mère achète-t-elle ? Qu'achète-t-elle ?)}
\item ¿Por qué compra productos locales? \emph{(Pourquoi achète-t-elle des produits locaux ?)}
\item ¿Qué hace el hermano? ¿Es responsable? \emph{(Que fait le frère ? Est-il responsable ?)}
\item ¿Qué debemos hacer, según el texto? \emph{(Que devons-nous faire, selon le texte ?)}
\item ¿Qué es el comercio justo? \emph{(Qu'est-ce que le commerce équitable ?)}
\end{enumerate}

\textbf{Repérage de la situation :} le narrateur est un enfant ou un adolescent d'une famille camerounaise. Les personnages sont la mère (consommatrice responsable) et le frère (victime de la publicité). La scène se passe au marché du quartier. Le thème est \cle{la consommation responsable} dans la vie quotidienne.

\begin{definition}[El consumo responsable]
Le \esp{consumo responsable} est une \textbf{manière de consommer qui respecte l'environnement, l'économie locale et les producteurs}. Le consommateur responsable réfléchit avant d'acheter : il préfère les produits locaux, évite le gaspillage, réutilise et recycle.
\end{definition}

\courssub{Lexique thématique : deux champs opposés}

Le texte s'organise autour de deux champs lexicaux opposés. Relevons les mots de chaque champ :

\begin{center}
\begin{tabularx}{\linewidth}{|X|X|X|}
\hline
\rowcolor{popblueL} \textbf{Consumo responsable} & \textbf{Derroche / Surconsommation} & \textbf{Traduction} \\
\hline
producto local & publicidad & produit local / publicité \\
\hline
comercio justo & compras innecesarias & commerce équitable / achats inutiles \\
\hline
ahorrar & derroche de dinero & économiser / gaspillage d'argent \\
\hline
reutilizar los envases & tirar la comida & réutiliser les emballages / jeter la nourriture \\
\hline
respetar el medio ambiente & dejarse llevar & respecter l'environnement / se laisser entraîner \\
\hline
precio digno & cosas innecesarias & prix juste / choses inutiles \\
\hline
reciclar & consumismo & recycler / consumérisme \\
\hline
\end{tabularx}
\end{center}

\begin{popfigure}
\begin{tikzpicture}[mindmap, grow cyclic, every node/.style={concept, font=\small},
  root concept/.append style={concept color=popblue, font=\small\bfseries, text=white},
  level 1/.append style={level distance=3.2cm, sibling angle=90},
  level 1 concept/.append style={font=\small}]
\node[root concept]{consumo responsable}
  child[concept color=popgreen] { node {producto local} }
  child[concept color=poporange] { node {comercio justo} }
  child[concept color=poppurple] { node {ahorro} }
  child[concept color=poppink] { node {no derroche} };
\end{tikzpicture}
\legende{Carte mentale : les quatre idées du texte autour du consumo responsable.}
\end{popfigure}

\begin{popfigure}
\begin{tikzpicture}
\node[draw=popgreen, fill=popgreenL, rounded corners, minimum width=5.6cm, minimum height=0.9cm, font=\small\bfseries] (a) at (0,0) {consumo responsable};
\node[draw=poppink, fill=poppinkL, rounded corners, minimum width=5.6cm, minimum height=0.9cm, font=\small\bfseries] (b) at (6.4,0) {derroche};
\node[draw=popline, rounded corners, align=left, font=\small, text width=5.2cm, anchor=north west] at (-2.8,-0.7)
  {producto local\\ comercio justo\\ ahorrar\\ reutilizar\\ reciclar\\ precio digno};
\node[draw=popline, rounded corners, align=left, font=\small, text width=5.2cm, anchor=north west] at (3.6,-0.7)
  {publicidad\\ derroche\\ compras innecesarias\\ tirar la comida\\ dejarse llevar\\ consumismo};
\end{tikzpicture}
\legende{Deux colonnes : les mots du texte rangés selon leur champ lexical (vocabulaire étendu).}
\end{popfigure}

\begin{attention}[Faux amis et confusions fréquentes]
\begin{itemize}
\item Ne confonds pas \esp{consumidor} (le \textbf{consommateur}, la personne) et \esp{consumo} (la \textbf{consommation}, l'action). \esp{El consumidor compra; el consumo aumenta.}
\item \esp{Derroche} est un faux ami partiel : il signifie \textbf{gaspillage}, et non « débauche ». \esp{Es un derroche de dinero} = « c'est un gaspillage d'argent ».
\item \esp{Publicidad} signifie « publicité » (activité commerciale), pas « notoriété ». On dit \esp{la publicidad engañosa} (la publicité mensongère).
\item \esp{Realizar} ne signifie pas « réaliser » (comprendre) mais \textbf{effectuer, réaliser} (une action). \esp{Realizar una compra} = effectuer un achat.
\end{itemize}
\end{attention}

\begin{savaistu}[Le commerce équitable, du café au cacao]
Le \esp{comercio justo} (commerce équitable) est né en partie en Amérique latine pour garantir un \esp{precio digno} aux producteurs de café et de cacao — deux produits que le Cameroun exporte aussi ! Quand un consommateur paie un prix juste, le paysan peut vivre dignement de son travail. En Guinée équatoriale, voisin hispanophone du Cameroun, le cacao est également une richesse importante. Des labels comme \esp{Fairtrade} ou \esp{WFTO} certifient ce commerce.
\end{savaistu}

\coursec{Grammaire outillée : Obligation et interdiction}

\courssub{Observation et règle}

Le texte contient trois structures essentielles pour exprimer l'obligation et l'interdiction. Observons-les dans leur contexte avant d'en fixer la règle.

\begin{center}
\begin{tabularx}{\linewidth}{|X|X|X|}
\hline
\rowcolor{popblueL} \textbf{Phrase du texte} & \textbf{Structure} & \textbf{Traduction} \\
\hline
\esp{Debemos consumir de manera responsable.} & \esp{deber} + infinitif & Nous devons consommer de manière responsable. \\
\hline
\esp{Tenemos que ahorrar.} & \esp{tener que} + infinitif & Nous devons économiser. \\
\hline
\esp{No se debe tirar la comida.} & \esp{no se debe} + infinitif & On ne doit pas jeter la nourriture. \\
\hline
\end{tabularx}
\end{center}

\begin{propriete}[Obligation : \esp{deber} vs \esp{tener que} vs \esp{hay que}]
\textbf{1. \esp{Deber} + infinitif} (obligation morale, conseil, devoir intérieur) \\
\esp{Debes estudiar más.} (Tu dois / Tu devrais étudier plus.) \\
\emph{Nuance :} au conditionnel (\esp{debería}), c'est un conseil : \esp{Deberías comer más verduras.}

\textbf{2. \esp{Tener que} + infinitif} (obligation extérieure, contrainte, nécessité factuelle) \\
\esp{Tengo que ir al mercado.} (Je dois aller au marché — obligation concrète.) \\
\emph{Attention :} \esp{tener} est un verbe à diphtongue (\esp{tengo, tienes, tiene, tenemos, tenéis, tienen}).

\textbf{3. \esp{Hay que} + infinitif} (obligation impersonnelle, règle générale, « il faut ») \\
\esp{Hay que respetar las normas.} (Il faut respecter les règles / On doit respecter les règles.) \\
\emph{Invariable :} toujours \esp{hay que} (pas de conjugaison de \esp{haber} au présent pour cette structure).

\textbf{4. \esp{No se debe} / \esp{No hay que} + infinitif} (interdiction générale, impersonnelle) \\
\esp{No se debe tirar basura.} = \esp{No hay que tirar basura.} (On ne doit pas / Il ne faut pas jeter des déchets.) \\
\emph{Différence :} \esp{No se debe} insiste sur le devoir moral ; \esp{No hay que} sur la règle pratique.
\end{propriete}

\begin{center}
\begin{tabular}{|l|l|l|l|}
\hline
\rowcolor{popblueL} \textbf{Personne} & \textbf{Deber} & \textbf{Tener que} & \textbf{Hay que} \\
\hline
yo & debo & tengo que & \multicolumn{1}{c|}{---} \\
\hline
tú & debes & tienes que & \multicolumn{1}{c|}{---} \\
\hline
él/ella/usted & debe & tiene que & \multicolumn{1}{c|}{\esp{hay que}} \\
\hline
nosotros & debemos & tenemos que & \multicolumn{1}{c|}{---} \\
\hline
vosotros & debéis & tenéis que & \multicolumn{1}{c|}{---} \\
\hline
ellos/ustedes & deben & tienen que & \multicolumn{1}{c|}{---} \\
\hline
\end{tabular}
\end{center}

\begin{exemplebox}[Exemples contrastés]
\begin{itemize}
\item \esp{Debo ayudar a mi madre.} (Obligation morale, je le sens.) \emph{« Je dois aider ma mère. »}
\item \esp{Tengo que ayudar a mi madre porque está enferma.} (Contrainte extérieure, circonstance.) \emph{« Je dois aider ma mère car elle est malade. »}
\item \esp{Hay que ayudar a los mayores.} (Règle générale, valable pour tous.) \emph{« Il faut aider les personnes âgées. »}
\item \esp{No se debe mentir.} (Interdiction morale générale.) \emph{« On ne doit pas mentir. »}
\item \esp{No hay que llegar tarde.} (Règle pratique.) \emph{« Il ne faut pas arriver en retard. »}
\end{itemize}
\end{exemplebox}

\begin{attention}[Pièges pour francophones]
\begin{itemize}
\item \textbf{Ne pas confondre \esp{deber} (verbe) et \esp{deber} (nom masculin : le devoir, la dette).} \esp{Tengo muchos deberes} = J'ai beaucoup de devoirs (devoirs scolaires).
\item \textbf{\esp{Tener que} : ne pas oublier \esp{que}.} \esp{*Tengo ir} est faux. \rightarrow \esp{Tengo que ir}.
\item \textbf{\esp{Hay que} : invariable !} On n'écrit jamais \esp{*hays que}, \esp{*había que} (imparfait : \esp{había que}), \esp{*habrá que} (futur). Au passé composé : \esp{ha habido que}.
\item \textbf{\esp{No se debe} : le \esp{se} est impersonnel.} Il ne s'accorde pas. \esp{No se deben tirar papeles} (pluriel du sujet patient \esp{papeles}) mais \esp{No se debe tirar papel} (singulier). À l'oral, on entend souvent \esp{no se deben} pour du pluriel, mais la norme cultive l'accord avec le sujet patient.
\end{itemize}
\end{attention}

\courssub{Application immédiate}

\begin{exoresolu}[Choix de la structure d'obligation]
\textbf{Énoncé :} Complète avec \esp{debo}, \esp{tengo que}, \esp{hay que} ou \esp{no se debe} (présent de l'indicatif) selon le contexte.
\begin{enumerate}
\item \esp{\_\_\_\_\_\_\_\_\_\_ reciclar las botellas de plástico.} (Règle générale pour tous)
\item \esp{Yo \_\_\_\_\_\_\_\_\_\_ estudiar para el examen de mañana.} (Obligation personnelle ressentie)
\item \esp{Mañana \_\_\_\_\_\_\_\_\_\_ levantarme temprano porque el autobús sale a las 6.} (Contrainte horaire extérieure)
\item \esp{\_\_\_\_\_\_\_\_\_\_ desperdiciar agua.} (Interdiction morale générale)
\end{enumerate}
\tcblower
\textbf{Solution :}
\begin{enumerate}
\item \esp{Hay que reciclar las botellas de plástico.} (Impersonnel, règle générale)
\item \esp{Yo \textbf{debo} estudiar para el examen de mañana.} (Devoir moral / conseil à soi-même)
\item \esp{Mañana \textbf{tengo que} levantarme temprano porque el autobús sale a las 6.} (Contrainte extérieure, horaire)
\item \esp{\textbf{No se debe} desperdiciar agua.} (Interdiction générale, impersonnelle)
\end{enumerate}
\end{exoresolu}

\coursec{Méthode : Le commentaire de texte en espagnol}

\courssub{Structure type (Bac / Évaluation)}

\begin{methode}[Commenter un texte : 4 étapes clés]
\begin{enumerate}
\item \textbf{Introduction (3-4 lignes)} :
      \begin{itemize}
      \item \textbf{Présenter} : titre (entre guillemets \esp{« ... »}), auteur (si connu, sinon \esp{texto anónimo}), thème général (\esp{el consumo responsable}), type de texte (article d'opinion, témoignage, dialogue...).
      \item \textbf{Problématique} : une question ouverte qui guide l'analyse. Ex. : \esp{¿Cómo promueve el texto un consumo más ético y sostenible?}
      \item \textbf{Annonce du plan} : \esp{En primer lugar...; en segundo lugar...; finalmente...}
      \end{itemize}
\item \textbf{Développement (2 ou 3 paragraphes)} : Un paragraphe = une idée force.
      \begin{itemize}
      \item \textbf{Phrase d'introduction} (idée directrice).
      \item \textbf{Citation courte} intégrée dans la phrase (\esp{El texto afirma que « ... »}).
      \item \textbf{Analyse / Explication} : vocabulaire, structures (ex. \esp{deber + inf.}), opposition (champs lexicaux), ton.
      \item \textbf{Mini-transition} vers le paragraphe suivant.
      \end{itemize}
\item \textbf{Conclusion (2-3 lignes)} :
      \begin{itemize}
      \item Réponse synthétique à la problématique.
      \item Ouverture : lien avec l'actualité (Cameroun, Guinée équatoriale), expérience personnelle, autre texte.
      \end{itemize}
\item \textbf{Relecture} : accords, accents, \esp{¿ ¡}, connecteurs logiques (\esp{por tanto, sin embargo, además, en conclusión}).
\end{enumerate}
\end{methode}

\begin{exemplebox}[Introduction modèle pour « Comprar menos, comprar mejor »]
\esp{El texto « Comprar menos, comprar mejor » es un testimonio en primera persona que aborda el tema del consumo responsable en un contexto camerunés. El narrador contrasta la actitud de su madre, que compra productos locales, con la de su hermano, influenciado por la publicidad. La pregunta que guía nuestro análisis es: ¿De qué manera el texto opone el consumo responsable al derroche para defender una ética de consumo? Veremos primero la oposición entre dos modelos de consumo, luego las propuestas concretas de acción responsable, y finalmente el valor ético del comercio justo.}
\end{exemplebox}

\begin{aretenir}[Connecteurs logiques indispensables]
\begin{center}
\begin{tabularx}{\linewidth}{|X|X|}
\hline
\rowcolor{popblueL} \textbf{Français} & \textbf{Español} \\
\hline
En premier lieu / D'abord & \esp{En primer lugar / Primero} \\
\hline
Ensuite / De plus & \esp{En segundo lugar / Además / Asimismo} \\
\hline
Cependant / Mais & \esp{Sin embargo / No obstante / Pero} \\
\hline
Par exemple & \esp{Por ejemplo / Como se ve en el texto} \\
\hline
Par conséquent / Donc & \esp{Por tanto / Por consiguiente / Así que} \\
\hline
En conclusion / Pour finir & \esp{En conclusión / Para terminar / En definitiva} \\
\hline
\end{tabularx}
\end{center}
\end{aretenir}

\coursec{Production guidée : Éthique professionnelle et consommation}

\courssub{Expression écrite : Charte de l'artisan responsable}

\begin{experience}[Jeu de rôle \& Rédaction : « La carta de compromiso »]
\textbf{Contexte :} Ton frère a ouvert son atelier de couture à Douala. Il veut afficher une \esp{Carta de compromiso profesional} (Charte d'engagement professionnel) sur le mur pour rassurer ses clients.

\textbf{Consigne :} Rédige cette charte en espagnol (80-100 mots). Tu dois :
\begin{enumerate}
\item Utiliser \textbf{au moins 3 obligations} (\esp{deber / tener que / hay que} + infinitif).
\item Utiliser \textbf{au moins 2 interdictions} (\esp{no se debe / no hay que} + infinitif).
\item Intégrer \textbf{5 mots du lexique} de la leçon (\esp{honestidad, puntualidad, respeto, precio digno, calidad, cliente, producto local, comercio justo}).
\item Soigner la présentation (puces, impératif ou infinitif de règle).
\end{enumerate}

\textbf{Aide au démarrage (phares de début de phrase) :}
\begin{itemize}
\item \esp{En este taller, nos comprometemos a...}
\item \esp{Es fundamental...}
\item \esp{El cliente tiene derecho a...}
\item \esp{No se permite...}
\item \esp{La base de nuestro trabajo es la...}
\end{itemize}
\end{experience}

\begin{exemplebox}[Modèle de production attendue (Corrigé type)]
\esp{
\textbf{CARTA DE COMPROMISO PROFESIONAL}
\begin{itemize}
\item \textbf{Honestidad:} \esp{Debemos ofrecer un precio digno y transparente. No se debe ocultar ningún costo.}
\item \textbf{Calidad:} \esp{Tenemos que usar telas de buena calidad, preferiblemente producto local.}
\item \textbf{Puntualidad:} \esp{Hay que entregar los pedidos en la fecha acordada.}
\item \textbf{Respeto:} \esp{El cliente merece un trato amable. No hay que ignorar sus sugerencias.}
\item \textbf{Ética:} \esp{Nos comprometemos con el comercio justo y el consumo responsable.}
\end{itemize}
}
\emph{Traduction : « CHARTE D'ENGAGEMENT PROFESSIONNEL — Honnêteté : Nous devons offrir un prix juste et transparent. On ne doit pas cacher de coûts. — Qualité : Nous devons utiliser de bons tissus, de préférence produit local. — Ponctualité : Il faut livrer les commandes à la date convenue. — Respect : Le client mérite un accueil aimable. Il ne faut pas ignorer ses suggestions. — Éthique : Nous nous engageons pour le commerce équitable et la consommation responsable. »}
\end{exemplebox}

\courssub{Expression orale : Débat « Consommer local »}

\begin{experience}[Débat guidé : \esp{¿Comprar importado o local?}]
\textbf{Organisation :} Deux groupes. Groupe A : \esp{Defienden el producto local}. Groupe B : \esp{Defienden la libertad de elegir (precio bajo)}.
\textbf{Durée :} 5 min préparation (mots-clés au tableau), 10 min débat.
\textbf{Contraintes linguistiques :} Chaque élève doit parler au moins une fois en utilisant :
\begin{itemize}
\item Une structure d'obligation (\esp{deber / tener que / hay que}).
\item Un connecteur (\esp{sin embargo, además, por tanto}).
\item Un mot du champ lexical (\esp{derroche, ahorro, comercio justo, publicidad}).
\end{itemize}
\textbf{Rôle du prof :} Noter les erreurs récurrentes (accords, \esp{ser/estar} si apparu, \esp{por/para}) pour correction collective finale.
\end{experience}

\courssub{Entraînement global (Évaluation formative)}

\begin{exercice}[Exercices de synthèse]
\begin{enumerate}
\item \textbf{Traduction (Thème) :} Traduis en espagnol.
      \begin{enumerate}
      \item Nous devons acheter des produits locaux pour aider l'économie.
      \item Il ne faut pas gaspiller l'eau ni l'électricité.
      \item Mon père doit se lever tôt demain pour aller au marché.
      \item En tant que consommateurs, nous avons le devoir de recycler.
      \item Le commerce équitable garantit un prix juste au producteur.
      \end{enumerate}
\item \textbf{Grammaire :} Transforme les phrases avec \esp{deber} en \esp{tener que} et vice-versa, en notant le changement de nuance.
      \begin{enumerate}
      \item \esp{Debes respetar a tus mayores.} $\rightarrow$ \esp{...}
      \item \esp{Tengo que trabajar el domingo.} $\rightarrow$ \esp{...}
      \end{enumerate}
\item \textbf{Version :} Traduis en français.
      \esp{« No se debe tirar la comida mientras otros pasan hambre. El consumidor responsable sabe que ahorrar no es guardar, es invertir en el futuro. »}
\item \textbf{Rédaction courte (40 mots) :} Tu es journaliste pour le journal du lycée. Écris un court article (titre + 3 phrases) pour inciter les élèves à apporter leur gourde au lieu d'acheter des bouteilles en plastique. Utilise \esp{hay que}, \esp{no se debe}, \esp{reutilizar}.
\end{enumerate}
\end{exercice}

\begin{corrige}[Exercices de synthèse]
\begin{enumerate}
\item \textbf{Thème :}
      \begin{enumerate}
      \item \esp{Debemos / Tenemos que comprar productos locales para ayudar a la economía.}
      \item \esp{No se debe / No hay que derrochar el agua ni la electricidad.}
      \item \esp{Mi padre tiene que levantarse temprano mañana para ir al mercado.} (Contrainte extérieure)
      \item \esp{Como consumidores, debemos reciclar.} (Devoir moral)
      \item \esp{El comercio justo garantiza un precio digno al productor.}
      \end{enumerate}
\item \textbf{Transformation :}
      \begin{enumerate}
      \item \esp{Tienes que respetar a tus mayores.} (Nuance : contrainte sociale/extérieure vs devoir moral).
      \item \esp{Debo trabajar el domingo.} (Nuance : obligation sentie comme un devoir vs contrainte subie).
      \end{enumerate}
\item \textbf{Version :} « On ne doit pas jeter la nourriture pendant que d'autres ont faim. Le consommateur responsable sait qu'économiser ce n'est pas mettre de côté, c'est investir dans l'avenir. »
\item \textbf{Rédaction (exemple) :}
      \esp{\textbf{¡Trae tu botella!} Hay que reducir el plástico en el colegio. No se debe comprar botellas de un solo uso. Tenemos que reutilizar nuestras gourdes cada día. ¡El planeta lo agradece!}
\end{enumerate}
\end{corrige}

\begin{aretenir}[Bilan de la leçon EA15]
\begin{itemize}
\item \textbf{Lexique :} \esp{consumo responsable, derroche, producto local, comercio justo, ahorrar, reutilizar, precio digno, consumidor, publicidad, honestidad, puntualidad, respeto.}
\item \textbf{Grammaire :} Obligation : \esp{deber} (moral), \esp{tener que} (extérieur), \esp{hay que} (général). Interdiction : \esp{no se debe}, \esp{no hay que} + infinitif.
\item \textbf{Compétences :} Lire/comprendre un texte d'opinion, opposer deux champs lexicaux, exprimer des règles de vie (éthique pro / conso), commenter un texte structuré, rédiger une charte / un article court.
\end{itemize}
\end{aretenir}

\coursec{Compréhension du texte}

Dans la section précédente, nous avons lu et découvert le texte \esp{« Comprar menos, comprar mejor »}, un récit simple où une mère camerounaise fait ses courses au marché de Mokolo à Yaoundé avec ses deux enfants. Nous allons maintenant vérifier et approfondir notre compréhension de ce texte : c'est l'étape où le lecteur passe de la \emph{découverte} à la \emph{maîtrise} du contenu. Au baccalauréat, les questions de compréhension représentent une part essentielle de l'épreuve : savoir y répondre correctement, en espagnol et avec méthode, est donc indispensable.

Pour travailler sereinement, relisons d'abord le texte support dans sa version intégrale.

\begin{textebox}[Comprar menos, comprar mejor (texte support)]
Todos los sábados, mi madre va al mercado de Mokolo, en Yaoundé, con mi hermana y conmigo. Ella compra productos locales: tomates, plátanos, ñame y pescado seco. « Los productos locales son frescos y ayudan a los productores de nuestro país », dice siempre. Mi madre compara los precios y no compra cosas innecesarias. « Debemos comprar lo que necesitamos, no lo que la publicidad quiere vender », explica.

Mi hermano mayor, en cambio, es diferente. Él compra ropa nueva cada mes, tira la comida que no le gusta y gasta todo su dinero en cosas inútiles. Mi madre le dice: « No se debe derrochar el dinero ni tirar la comida. El derroche es un problema para la familia y para el planeta ».

En casa, tenemos reglas claras: debemos apagar la luz cuando salimos de una habitación, tenemos que ahorrar agua y no se debe comprar lo que no necesitamos. Mi madre también es vendedora en el mercado y siempre es honesta con sus clientes: da el precio justo y nunca engaña a nadie. « La honestidad es la base del trabajo », afirma.

Yo pienso que mi madre tiene razón: comprar menos y comprar mejor es bueno para la familia, para los productores locales y para el medio ambiente. El consumo responsable empieza con pequeñas acciones de cada día.
\end{textebox}

\textbf{Glossaire :} \esp{plátanos} : bananes plantains ; \esp{ñame} : igname ; \esp{pescado seco} : poisson séché ; \esp{innecesarias} : inutiles, pas nécessaires ; \esp{tira} : jette (verbe \esp{tirar}) ; \esp{gasta} : dépense (verbe \esp{gastar}) ; \esp{derrochar} : gaspiller ; \esp{el derroche} : le gaspillage ; \esp{apagar la luz} : éteindre la lumière ; \esp{ahorrar} : économiser ; \esp{engañar} : tromper ; \esp{el precio justo} : le prix juste ; \esp{tiene razón} : a raison ; \esp{el medio ambiente} : l'environnement.

\courssub{La méthode pour répondre aux questions de compréhension}

Avant de répondre aux questions, il faut connaître la bonne méthode. Beaucoup d'élèves perdent des points non pas parce qu'ils n'ont pas compris le texte, mais parce qu'ils répondent mal.

\begin{methode}[Comment répondre à une question de compréhension]
\begin{enumerate}
\item \textbf{Lire la question attentivement} et souligner le mot interrogatif : \esp{¿Qué?} (que), \esp{¿Quién?} (qui), \esp{¿Por qué?} (pourquoi), \esp{¿Dónde?} (où), \esp{¿Cuándo?} (quand).
\item \textbf{Repérer dans le texte} la phrase ou le passage qui contient la réponse.
\item \textbf{Rédiger une phrase complète en espagnol}, en reprenant les mots de la question. Exemple : à la question \esp{¿Qué compra la madre?}, on ne répond pas seulement \esp{« productos locales »}, mais : \esp{« La madre compra productos locales. »}
\item \textbf{Citer le texte} entre guillemets quand on demande une justification : \esp{El texto dice: « ... »}.
\item \textbf{Relire sa réponse} pour vérifier l'orthographe, les accents et l'accord des adjectifs.
\end{enumerate}
\end{methode}

\begin{attention}[Le piège de la réponse à « ¿Por qué... ? »]
À la question \esp{¿Por qué... ?} (pourquoi), la réponse commence \textbf{toujours} par \esp{porque} (parce que), sans accent et en un seul mot. Mais attention : au baccalauréat, il est \textbf{interdit} de répondre par un simple \esp{porque} suivi de trois mots arrachés au texte. Il faut une \textbf{phrase complète} avec sujet et verbe.

\begin{itemize}
\item Mauvaise réponse : \esp{Porque son frescos.} (phrase sans sujet clair, trop courte)
\item Bonne réponse : \esp{La madre compra productos locales porque son frescos y ayudan a los productores del país.}
\end{itemize}

Ne confonds pas les quatre formes : \esp{porque} (parce que), \esp{¿por qué?} (pourquoi), \esp{por qué} (pour quelle raison, dans les interrogatives indirectes), \esp{porqué} (le motif, nom masculin). À ton niveau, retiens surtout les deux premières.
\end{attention}

\courssub{Comprendre sans dictionnaire : les mots transparents}

Avant même de répondre aux questions, remarque que le texte contient beaucoup de mots que tu comprends déjà sans les avoir appris : ce sont les \cle{mots transparents}, c'est-à-dire les mots espagnols qui ressemblent au français parce qu'ils ont la même origine latine.

\begin{propriete}[Les mots transparents, alliés de la compréhension]
Un mot transparent se reconnaît par sa forme proche du français. Dans notre texte : \esp{responsable} (responsable), \esp{publicidad} (publicité), \esp{productor} (producteur), \esp{necesario} (nécessaire), \esp{familia} (famille), \esp{problema} (problème), \esp{planeta} (planète), \esp{honestidad} (honnêteté), \esp{acciones} (actions), \esp{importados} (importés), \esp{diferente} (différent), \esp{inútiles} (inutiles).

Stratégie : quand tu rencontres un mot inconnu, demande-toi d'abord s'il ressemble à un mot français. Cela suffit souvent à comprendre le sens général de la phrase.
\end{propriete}

\begin{attention}[Attention aux faux amis]
Tous les mots qui se ressemblent ne se traduisent pas de la même façon ! \esp{Publicidad} signifie « publicité » (et non « publicité » au sens de « notoriété publique ») ; \esp{gastar} signifie « dépenser » (et non « gâter ») ; \esp{tirar} signifie « jeter » (et non « tirer ») ; \esp{realizar} signifie « réaliser » au sens d'« accomplir » (et non « réaliser » au sens de « comprendre »). Vérifie toujours le contexte avant de traduire un mot transparent.
\end{attention}

\courssub{Questions de compréhension globale}

Les premières questions portent sur le sens général du texte. On y répond après une première lecture, sans entrer dans les détails.

\begin{exemplebox}[Questions globales et réponses modèles]
\textbf{1. ¿De qué trata el texto?} (De quoi parle le texte ?)

\esp{El texto trata del consumo responsable y de la ética profesional: una madre compra productos locales y enseña a sus hijos a no derrochar.}

« Le texte parle de la consommation responsable et de l'éthique professionnelle : une mère achète des produits locaux et apprend à ses enfants à ne pas gaspiller. »

\textbf{2. ¿Quién habla en el texto?} (Qui parle dans le texte ?)

\esp{Habla un narrador, un hijo de la familia, que cuenta la historia de su madre.}

« C'est un narrateur, un enfant de la famille, qui raconte l'histoire de sa mère. »

\textbf{3. ¿Dónde y cuándo ocurre la historia?}

\esp{La historia ocurre en el mercado de Mokolo, en Yaoundé, todos los sábados.}
\end{exemplebox}

Observe la formule magique pour la question « de quoi s'agit-il » : \esp{El texto trata de...} + nom ou \esp{El texto trata de...} + infinitif. C'est la réponse type attendue au baccalauréat.

\courssub{Questions de détail}

Les questions de détail exigent de retrouver une information précise dans le texte. La réponse doit reprendre les mots de la question et citer le texte.

\begin{exemplebox}[Questions de détail et réponses modèles]
\textbf{1. ¿Qué compra la madre?} (Qu'achète la mère ?)

\esp{La madre compra productos locales: tomates, plátanos, ñame y pescado seco.}

« La mère achète des produits locaux : des tomates, des bananes plantains, de l'igname et du poisson séché. »

\textbf{2. ¿Por qué compra productos locales?} (Pourquoi achète-t-elle des produits locaux ?)

\esp{Compra productos locales porque son frescos y ayudan a los productores de su país. El texto dice: « Los productos locales son frescos y ayudan a los productores de nuestro país ».}

\textbf{3. ¿Qué reglas hay en la casa?}

\esp{En la casa hay reglas claras: deben apagar la luz, tienen que ahorrar agua y no se debe comprar lo que no necesitan.}
\end{exemplebox}

Remarque la structure de la réponse justifiée : d'abord la réponse complète, ensuite la citation introduite par \esp{El texto dice: « ... »} (« Le texte dit : ... »). C'est exactement la démarche attendue dans l'épreuve de commentaire.

\courssub{Question d'opposition : le contraste entre les personnages}

Le texte oppose deux attitudes : celle de la mère (responsable) et celle du frère (gaspilleur). Ce contraste est signalé par la locution \esp{en cambio} (« en revanche, par contre »), un connecteur d'opposition très utile.

\begin{exemplebox}[Question d'opposition et réponse modèle]
\textbf{1. ¿Qué hace el hermano?} (Que fait le frère ?)

\esp{El hermano compra ropa nueva cada mes, tira la comida que no le gusta y gasta todo su dinero en cosas inútiles.}

« Le frère achète des vêtements neufs chaque mois, jette la nourriture qu'il n'aime pas et dépense tout son argent en choses inutiles. »

\textbf{2. ¿Qué problema representa su actitud?} (Quel problème représente son attitude ?)

\esp{Su actitud representa el derroche, que es un problema para la familia y para el planeta. El texto dice: « El derroche es un problema para la familia y para el planeta ».}
\end{exemplebox}

\begin{center}
\begin{tabularx}{\linewidth}{|X|X|X|}
\hline
\rowcolor{popblueL} La madre (responsable) & El hermano (derrochador) & Conector \\
\hline
\esp{compra productos locales} & \esp{compra ropa nueva cada mes} & \esp{en cambio} \\
\hline
\esp{compara los precios} & \esp{gasta todo su dinero} & \esp{en cambio} \\
\hline
\esp{no compra cosas innecesarias} & \esp{tira la comida} & \esp{en cambio} \\
\hline
\end{tabularx}
\end{center}

\courssub{Questions sur les règles énoncées dans le texte}

Le texte formule des règles de vie avec les verbes d'obligation (\esp{deber}, \esp{tener que}) et d'interdiction (\esp{no se debe}). Ces questions préparent directement la section de grammaire.

\begin{exemplebox}[Questions sur les règles et réponses modèles]
\textbf{1. ¿Qué debemos hacer según el texto?} (Que devons-nous faire selon le texte ?)

\esp{Según el texto, debemos apagar la luz cuando salimos de una habitación, ahorrar agua y comprar lo que necesitamos.}

« Selon le texte, nous devons éteindre la lumière en sortant d'une pièce, économiser l'eau et acheter ce dont nous avons besoin. »

\textbf{2. ¿Qué no se debe hacer?} (Que ne doit-on pas faire ?)

\esp{No se debe derrochar el dinero ni tirar la comida.}

« On ne doit pas gaspiller l'argent ni jeter la nourriture. »
\end{exemplebox}

Note la construction impersonnelle \esp{no se debe} + infinitif : elle exprime une interdiction générale, comme le « on ne doit pas » français ou les panneaux « défense de... ». Exemple : \esp{No se debe fumar.} « Il est interdit de fumer. »

\courssub{Question d'interprétation : comprendre le titre}

La question d'interprétation est la plus difficile : la réponse ne se trouve pas mot à mot dans le texte, il faut la \emph{déduire} de l'ensemble.

\begin{exoresolu}[Interpréter la phrase « comprar menos, comprar mejor »]
Énoncé : \esp{¿Qué significa la frase « comprar menos, comprar mejor »?} Explique le sens du titre avec tes propres mots.
\tcblower
Solution détaillée : le titre résume toute la philosophie du texte. \esp{Comprar menos} signifie ne pas acheter de choses inutiles, éviter le gaspillage d'argent et d'objets (comme le fait le frère). \esp{Comprar mejor} signifie choisir la qualité : des produits locaux, frais, qui aident les producteurs du pays et respectent l'environnement (comme le fait la mère).

Réponse modèle complète en espagnol :

\esp{La frase « comprar menos, comprar mejor » significa que no debemos comprar muchas cosas innecesarias, sino elegir productos de calidad, locales y útiles. Es mejor para la familia, para los productores y para el medio ambiente.}

« La phrase "acheter moins, acheter mieux" signifie que nous ne devons pas acheter beaucoup de choses inutiles, mais choisir des produits de qualité, locaux et utiles. C'est meilleur pour la famille, pour les producteurs et pour l'environnement. »

Méthode : pour une question d'interprétation, 1) repère les passages liés à l'idée, 2) dégage l'idée générale, 3) reformule avec \emph{tes} mots, sans recopier une phrase entière du texte.
\end{exoresolu}

\courssub{Vue d'ensemble : l'organigramme de compréhension}

Pour retenir la structure logique du texte, voici un organigramme qui va du thème général jusqu'à la morale. Sache refaire ce schéma pour n'importe quel texte narratif ou argumentatif : c'est un excellent outil de révision avant le commentaire.

\begin{popfigure}
\begin{center}
\begin{tikzpicture}[
  node distance=0.55cm and 1.1cm,
  box/.style={draw=popblue, thick, rounded corners, fill=popblueL, align=center, inner sep=5pt, font=\small},
  arr/.style={-{Stealth[length=2.5mm]}, thick, popdark}
]
\node[box, fill=popgoldL, draw=popgold, align=center] (tema){\textbf{Tema}\\ consumo responsable\\ y ética profesional};
\node[box, below left=0.7cm and -0.4cm of tema, align=center] (pers){\textbf{Personajes}\\ la madre (responsable)\\ el hermano (derrochador)};
\node[box, below right=0.7cm and -0.4cm of tema, align=center] (prob){\textbf{Problema}\\ el derroche : gastar\\ y tirar la comida};
\node[box, below=0.7cm of $(pers)!0.5!(prob)$, yshift=-0.6cm, align=center] (sol){\textbf{Solución}\\ reglas claras : ahorrar,\\ comprar local, ser honesto};
\node[box, below=0.7cm of sol, fill=poppurpleL, draw=poppurple, align=center] (mor){\textbf{Morale}\\ « comprar menos, comprar mejor »};
\draw[arr] (tema) -- (pers);
\draw[arr] (tema) -- (prob);
\draw[arr] (pers) -- (sol);
\draw[arr] (prob) -- (sol);
\draw[arr] (sol) -- (mor);
\end{tikzpicture}
\end{center}
\legende{Organigramme de compréhension : du thème à la morale du texte.}
\end{popfigure}

\courssub{Vrai ou faux justifié}

L'exercice du vrai ou faux (\esp{verdadero o falso}) est un classique du baccalauréat. Règle d'or : chaque réponse doit être justifiée par une \textbf{citation exacte} du texte, introduite par \esp{El texto dice: « ... »} ou \esp{Según el texto, ...}.

\begin{exercice}[¿Verdadero o falso? Justifica con una cita del texto]
Indique si chaque phrase est vraie (\esp{verdadero}) ou fausse (\esp{falso}) et justifie par une citation du texte.

\begin{enumerate}
\item \esp{La madre compra productos importados de Europa.}
\item \esp{El hermano tira la comida que no le gusta.}
\item \esp{En la casa, no hay reglas de consumo.}
\item \esp{La madre es vendedora honesta con sus clientes.}
\item \esp{El narrador piensa que su madre no tiene razón.}
\end{enumerate}
\end{exercice}

\begin{corrige}[Exercice : verdadero o falso]
\begin{enumerate}
\item \textbf{Falso.} \esp{El texto dice: « Ella compra productos locales: tomates, plátanos, ñame y pescado seco ».} La mère achète des produits \emph{locaux}, pas importés.
\item \textbf{Verdadero.} \esp{El texto dice: « tira la comida que no le gusta ».}
\item \textbf{Falso.} \esp{El texto dice: « En casa, tenemos reglas claras ».} Il y a bien des règles : éteindre la lumière, économiser l'eau.
\item \textbf{Verdadero.} \esp{El texto dice: « siempre es honesta con sus clientes: da el precio justo y nunca engaña a nadie ».}
\item \textbf{Falso.} \esp{El texto dice: « Yo pienso que mi madre tiene razón ».} Le narrateur est d'accord avec sa mère.
\end{enumerate}
\end{corrige}

\begin{attention}[Erreurs fréquentes dans le vrai ou faux]
\begin{itemize}
\item Justifier avec une phrase \emph{inventée} au lieu d'une citation exacte : la justification doit reprendre les mots du texte.
\item Répondre \esp{verdadero} ou \esp{falso} sans justification : au Bac, une réponse non justifiée ne vaut aucun point.
\item Mal recopier la citation (oublier les accents : \esp{habitación}, \esp{ñame}, \esp{está}). Une citation fautive est une justification fausse.
\end{itemize}
\end{attention}

\courssub{Reformulation : résumer le texte en deux phrases}

Dernière étape de la compréhension : le résumé (\esp{el resumen}). Résumer, c'est dire l'essentiel avec \emph{ses propres mots}, sans détails ni citations. En deux phrases, on doit trouver : 1) le thème et les personnages, 2) le problème et la leçon du texte.

\begin{exemplebox}[Modèle de résumé en deux phrases]
\esp{El texto presenta a una madre de Yaoundé que practica el consumo responsable: compra productos locales y es honesta con sus clientes. Frente al derroche de su hijo, ella enseña que debemos comprar menos y comprar mejor para ayudar a la familia, a los productores y al planeta.}

« Le texte présente une mère de Yaoundé qui pratique la consommation responsable : elle achète des produits locaux et est honnête avec ses clients. Face au gaspillage de son fils, elle enseigne que nous devons acheter moins et acheter mieux pour aider la famille, les producteurs et la planète. »
\end{exemplebox}

\begin{methode}[Réussir un résumé en deux phrases]
\begin{enumerate}
\item \textbf{Phrase 1} : thème + personnage principal + situation (\esp{El texto presenta a... que...}).
\item \textbf{Phrase 2} : problème ou opposition + leçon du texte (\esp{Frente a..., enseña que...}).
\item Utilise des connecteurs : \esp{frente a} (face à), \esp{sin embargo} (cependant), \esp{por eso} (c'est pourquoi).
\item Interdiction : ne recopie aucune phrase entière du texte et ne donne pas ton opinion personnelle dans un résumé.
\end{enumerate}
\end{methode}

\begin{exercice}[À ton tour]
Rédige ton propre résumé du texte en deux phrases complètes en espagnol, sans recopier le modèle ci-dessus. Utilise au moins un connecteur (\esp{sin embargo}, \esp{por eso}, \esp{frente a}). Puis réponds à cette question d'opinion : \esp{¿Y tú, compras productos locales? ¿Por qué?} (réponse en deux ou trois phrases commençant par \esp{Sí, compro... porque...} ou \esp{No compro... porque...}).
\end{exercice}

Tu sais désormais lire finement le texte : identifier son thème, ses personnages, son opposition centrale, ses règles et sa morale, puis le résumer. Ces compétences de compréhension vont maintenant nous servir à construire le \cle{lexique thématique} de la consommation responsable, objet de la section suivante.

\coursec{El consumo responsable y la ética profesional}

\courssub{Texte support : « Comprar menos, comprar mejor »}

\begin{textebox}[Comprar menos, comprar mejor]
Vivimos en una sociedad de consumo. La publicidad nos invade en la televisión, en las redes sociales y en las calles de Yaoundé o Douala. Nos dice: « ¡Compra! ¡Necesitas esto! ¡Sé feliz! ». Pero, ¿es cierto? A menudo compramos cosas innecesarias: ropa que no usamos, gadgets que se rompen pronto, comida que tiramos a la basura. Este derroche tiene consecuencias graves: contaminación, agotamiento de los recursos naturales y deudas para las familias.

La alternativa es el consumo responsable. Consumir responsable significa comprar menos, pero comprar mejor. Significa elegir productos locales, como el cacao, el plátano o el café de nuestros campesinos, y apoyar el comercio justo. Cuando pagamos un precio justo al productor, él vive dignamente y cuida su tierra. También significa ahorrar: reparar en lugar de tirar, reutilizar bolsas y cajas, reciclar papel y plástico. El ahorro no es solo guardar dinero; es respetar el medio ambiente.

Detrás de todo buen comercio hay una ética profesional. Un vendedor honrado no engaña al cliente sobre la calidad o el precio. Llega puntual a su puesto, trata bien a la gente y cumple su palabra. La honestidad, la puntualidad y el respeto al cliente crean confianza. Y la confianza es la base de cualquier negocio duradero. Como dice el proverbio: \esp{El cliente siempre tiene razón}. No para que el comerciante pierda, sino para que gane el respeto de todos.
\end{textebox}

\textbf{Glossaire :} \esp{la sociedad de consumo} : la société de consommation ; \esp{invadir} : envahir ; \esp{las redes sociales} : les réseaux sociaux ; \esp{el gadget} : le gadget ; \esp{tirar a la basura} : jeter à la poubelle ; \esp{el agotamiento} : l'épuisement ; \esp{los recursos naturales} : les ressources naturelles ; \esp{las deudas} : les dettes ; \esp{la alternativa} : l'alternative ; \esp{el campesino} : le paysan ; \esp{el precio justo} : le prix juste ; \esp{vivir dignamente} : vivre dignement ; \esp{reparar} : réparer ; \esp{cumplir su palabra} : tenir sa parole ; \esp{duradero/a} : durable.

\courssub{Compréhension du texte}

\begin{exercice}[Questions de compréhension]
\begin{enumerate}
\item \textbf{Question globale :} Quel est le message principal de ce texte ? Résumez-le en deux phrases en français.
\item \textbf{Vocabulaire en contexte :} Expliquez en français le sens de \esp{derroche}, \esp{comercio justo} et \esp{ética profesional} d'après le texte.
\item \textbf{Informations explicites :} Citez trois exemples de « cosas innecesarias » mentionnés dans le texte.
\item \textbf{Conséquences :} Selon le texte, quelles sont les trois conséquences graves du gaspillage (derroche) ?
\item \textbf{Actions concrètes :} Le texte propose quatre actions pour consommer responsable. Lesquelles ? (Produits locaux, commerce juste, réparer/réutiliser/recycler, économiser).
\item \textbf{Éthique professionnelle :} Quelles sont les trois qualités du « vendedor honrado » citées dans le dernier paragraphe ?
\item \textbf{Inférence :} Pourquoi l'auteur dit-il que « el ahorro no es solo guardar dinero » ? Justifiez par le texte.
\end{enumerate}
\end{exercice}

\begin{corrige}[Exercice « Compréhension du texte »]
\begin{enumerate}
\item Le texte dénonce la surconsommation poussée par la publicité et propose le consommation responsable (acheter moins mais mieux : local, équitable, durable) fondée sur une éthique professionnelle (honnêteté, respect).
\item \esp{Derroche} = gaspillage (jeter de la nourriture, acheter inutilement). \esp{Comercio justo} = système qui garantit un prix juste aux producteurs (paysans). \esp{Ética profesional} = ensemble de valeurs (honnêteté, ponctualité, respect) qui guident le bon commerçant.
\item Vêtements qu'on n'utilise pas, gadgets qui cassent vite, nourriture qu'on jette.
\item Pollution (contaminación), épuisement des ressources naturelles (agotamiento de recursos), dettes familiales (deudas para las familias).
\item 1. Choisir des produits locaux. 2. Soutenir le commerce équitable. 3. Réparer, réutiliser, recycler. 4. Économiser (ahorrar).
\item Honnêteté (honestidad / ser honrado), ponctualité (llegar puntual), respect du client (tratar bien al cliente / respeto al cliente).
\item Parce que l'épargne inclut aussi le respect de l'environnement : réparer au lieu de jeter, réutiliser, recycler. C'est une gestion durable des ressources, pas seulement financière.
\end{enumerate}
\end{corrige}

\coursec{Lexique thématique par champs}

Après avoir lu et compris le texte \esp{Comprar menos, comprar mejor}, nous allons maintenant organiser son vocabulaire en \cle{champs lexicaux}. Un champ lexical est un ensemble de mots qui appartiennent à un même thème : ici, la consommation, la publicité, l'économie responsable et l'éthique professionnelle. Classer les mots par champs est la meilleure méthode pour les mémoriser durablement et pour pouvoir les réemployer dans vos propres productions écrites et orales.

\courssub{Pourquoi travailler par champs lexicaux ?}

\begin{methode}[Mémoriser le vocabulaire par champs]
\begin{enumerate}
\item \textbf{Repérer} dans le texte tous les mots qui se rapportent au même thème.
\item \textbf{Classer} ces mots dans un tableau : nom, verbe, adjectif, traduction.
\item \textbf{Observer} les familles de mots (même racine, suffixes différents).
\item \textbf{Réemployer} chaque mot dans une phrase personnelle, si possible liée à votre vie quotidienne (le marché, la boutique du quartier, l'école).
\item \textbf{Vérifier} les faux amis et les pièges orthographiques (accents, h muette).
\end{enumerate}
\end{methode}

\courssub{Champ 1 : La consommation}

Observons d'abord ces phrases tirées du texte support :

\esp{El consumidor compra muchos productos cada semana.} « Le consommateur achète beaucoup de produits chaque semaine. »

\esp{Debemos consumir menos y comprar mejor.} « Nous devons consommer moins et acheter mieux. »

On remarque que trois mots partagent la même racine \cle{consum-} : le nom \esp{el consumo}, la personne \esp{el consumidor}, le verbe \esp{consumir}. C'est une famille de mots, exactement comme en français « consommation / consommateur / consommer ».

\begin{definition}[Champ de la consommation]
\begin{itemize}
\item \esp{el consumo} : la consommation. \esp{El consumo de agua aumenta en verano.} « La consommation d'eau augmente en été. »
\item \esp{el consumidor / la consumidora} : le consommateur / la consommatrice. \esp{La consumidora elige productos locales.} « La consommatrice choisit des produits locaux. »
\item \esp{consumir} : consommer (verbe régulier en \esp{-ir}). \esp{Consumimos demasiada carne.} « Nous consommons trop de viande. »
\item \esp{la compra} : l'achat (acte d'acheter). \esp{Hago la compra en el mercado de Mvog-Mbi.} « Je fais les courses au marché de Mvog-Mbi. »
\item \esp{comprar} : acheter. \esp{Compro fruta fresca cada sábado.} « J'achète des fruits frais chaque samedi. »
\item \esp{el producto} : le produit. \esp{Este producto es de buena calidad.} « Ce produit est de bonne qualité. »
\end{itemize}
\end{definition}

\begin{attention}[Faire la compra ou hacer la compra ?]
En espagnol, « faire les courses » (les achats quotidiens) se dit \esp{hacer la compra} (et non \esp{hacer compras}, qui signifie plutôt « faire des achats » de façon générale ou ponctuelle). Ne calquez pas le français « faire un achat » : on dira \esp{hacer una compra} pour un achat précis et unique.
\end{attention}

\courssub{Champ 2 : La publicité et le gaspillage}

\esp{La publicidad nos invita a comprar cosas innecesarias.} « La publicité nous invite à acheter des choses inutiles. »

\esp{El derroche de comida es un problema grave.} « Le gaspillage de nourriture est un problème grave. »

\begin{definition}[Champ de la publicité et du gaspillage]
\begin{itemize}
\item \esp{la publicidad} : la publicité (au sens commercial : ensemble des annonces). \esp{Veo mucha publicidad en la televisión.} « Je vois beaucoup de publicité à la télévision. »
\item \esp{el anuncio} : l'annonce publicitaire, le spot, l'affiche. \esp{El anuncio del nuevo teléfono es muy largo.} « La pub du nouveau téléphone est très longue. »
\item \esp{el derroche} : le gaspillage (nom masculin). \esp{Tirar comida es un derroche.} « Jeter de la nourriture est un gaspillage. »
\item \esp{derrochar} : gaspiller (verbe régulier en \esp{-ar}). \esp{No derroches el agua, por favor.} « Ne gaspille pas l'eau, s'il te plaît. »
\item \esp{innecesario/a} : inutile, superflu. \esp{Comprar otro bolso es innecesario.} « Acheter un autre sac est inutile. »
\item \esp{la basura} : les ordures, la poubelle. \esp{Tira el papel a la basura.} « Jette le papier à la poubelle. »
\end{itemize}
\end{definition}

\begin{attention}[La publicidad : attention au sens]
\esp{La publicidad} signifie « la publicité » uniquement au sens \textbf{commercial} (annonces, spots, affiches). Si « publicité » signifie « notoriété, fait de rendre public », l'espagnol utilise d'autres mots comme \esp{la difusión} ou \esp{la divulgación}. De même, « faire de la publicité pour un produit » se dit \esp{hacer publicidad de un producto} ou \esp{anunciar un producto}.
\end{attention}

\courssub{Champ 3 : L'économie et la consommation responsable}

\esp{El ahorro es importante para la familia.} « L'épargne est importante pour la famille. »

\esp{Comprar productos locales ayuda a los campesinos y protege el medio ambiente.} « Acheter des produits locaux aide les paysans et protège l'environnement. »

\begin{definition}[Champ de la consommation responsable]
\begin{itemize}
\item \esp{el ahorro} : l'économie, l'épargne (nom masculin). \esp{Gracias al ahorro, mi padre compró una moto.} « Grâce à ses économies, mon père a acheté une moto. »
\item \esp{ahorrar} : économiser (verbe régulier \esp{-ar}). \esp{Ahorro mil francos cada mes.} « J'économise mille francs chaque mois. »
\item \esp{el ahorrador / la ahorradora} : l'économe (personne qui économise). \esp{Mi madre es muy ahorradora.} « Ma mère est très économe. »
\item \esp{el producto local} : le produit local. \esp{El cacao es un producto local del Camerún.} « Le cacao est un produit local du Cameroun. »
\item \esp{el comercio justo} : le commerce équitable. \esp{Apoyamos el comercio justo comprando café certificado.} « Nous soutenons le commerce équitable en achetant du café certifié. »
\item \esp{reutilizar} : réutiliser. \esp{Reutilizamos las bolsas de plástico.} « Nous réutilisons les sacs en plastique. »
\item \esp{reciclar} : recycler. \esp{En la escuela reciclamos el papel.} « À l'école, nous recyclons le papier. »
\item \esp{el medio ambiente} : l'environnement (toujours au masculin malgré sa forme en \esp{-e}). \esp{Debemos cuidar el medio ambiente.} « Nous devons protéger l'environnement. »
\item \esp{cuidar} : prendre soin de, protéger. \esp{Cuidar el planeta es tarea de todos.} « Protéger la planète est l'affaire de tous. »
\end{itemize}
\end{definition}

\begin{definition}[Le comercio justo]
\esp{El comercio justo} (le commerce équitable) est un système commercial qui garantit un prix équitable et de bonnes conditions de travail aux producteurs, par exemple aux cultivateurs de cacao ou de café. Grâce au commerce équitable, le paysan reçoit un prix stable et juste pour sa récolte, et non un prix imposé par les intermédiaires.
\end{definition}

\begin{exemplebox}[Exemples traduits]
\begin{itemize}
\item \esp{El ahorro es importante para la familia.} « L'épargne est importante pour la famille. »
\item \esp{La honestidad es la base del comercio.} « L'honnêteté est la base du commerce. »
\item \esp{Reciclar y reutilizar protegen el medio ambiente.} « Recycler et réutiliser protègent l'environnement. »
\item \esp{Comprar productos locales es consumir de forma responsable.} « Acheter des produits locaux, c'est consommer de façon responsable. »
\end{itemize}
\end{exemplebox}

\courssub{Champ 4 : L'éthique professionnelle}

\esp{Un vendedor honrado no engaña a sus clientes.} « Un vendeur honnête ne trompe pas ses clients. »

\esp{La puntualidad y el respeto al cliente crean confianza.} « La ponctualité et le respect du client créent la confiance. »

\begin{definition}[Champ de l'éthique professionnelle]
\begin{itemize}
\item \esp{la honestidad} : l'honnêteté (nom féminin, h muet). \esp{La honestidad del comerciante atrae clientes.} « L'honnêteté du commerçant attire des clients. »
\item \esp{honrado/a} : honnête (adjectif le plus courant). \esp{Mi tío es un hombre honrado.} « Mon oncle est un homme honnête. »
\item \esp{la responsabilidad} : la responsabilité. \esp{Cada trabajador tiene su responsabilidad.} « Chaque travailleur a sa responsabilité. »
\item \esp{la puntualidad} : la ponctualité. \esp{La puntualidad es una cualidad profesional esencial.} « La ponctualité est une qualité professionnelle essentielle. »
\item \esp{el respeto al cliente} : le respect du client. \esp{El respeto al cliente es fundamental.} « Le respect du client est fondamental. »
\item \esp{la calidad} : la qualité. \esp{Vendo productos de buena calidad.} « Je vends des produits de bonne qualité. »
\item \esp{la confianza} : la confiance. \esp{Los clientes tienen confianza en ella.} « Les clients ont confiance en elle. »
\item \esp{engañar} : tromper, duper. \esp{Un vendedor honrado no engaña nunca.} « Un vendeur honnête ne trompe jamais. »
\item \esp{cumplir} : tenir (sa parole, un engagement), respecter (une règle). \esp{Cumplo mi palabra.} « Je tiens ma parole. »
\end{itemize}
\end{definition}

\begin{attention}[Honestidad et honrado : l'orthographe et l'usage]
Le nom \esp{la honestidad} s'écrit avec un \textbf{h muet} (on prononce « o-nes-ti-dad »). L'adjectif courant, de registre courant à soutenu, est \esp{honrado/a} (« on-ra-do ») : \esp{un vendedor honrado}, \esp{una mujer honrada}. L'adjectif \esp{honesto/a} existe mais est plus rare, plus soutenu, souvent littéraire ; à votre niveau, employez toujours \esp{honrado/a}. N'oubliez pas l'accord en genre : \esp{honrado} (masc.), \esp{honrada} (fém.).
\end{attention}

\begin{savaistu}[El cliente siempre tiene razón]
Dans tout le monde hispanophone, on connaît la règle commerciale : \esp{El cliente siempre tiene razón} — « Le client a toujours raison ». Cette expression résume l'éthique du bon commerçant : écouter le client, le respecter et ne jamais le tromper, même quand on n'est pas d'accord avec lui. À Yaoundé comme à Madrid ou à Malabo, un commerçant honnête garde ses clients toute la vie. Notez la structure \esp{tener razón} (avoir raison) : \esp{Yo tengo razón, tú tienes razón, él tiene razón}.
\end{savaistu}

\courssub{Expressions utiles à réemployer}

\begin{exemplebox}[Quatre expressions clés pour l'oral et l'écrit]
\begin{itemize}
\item \esp{comprar productos locales} : acheter local. \esp{Prefiero comprar productos locales en el mercado.} « Je préfère acheter local au marché. »
\item \esp{pagar un precio justo} : payer un prix juste. \esp{En el comercio justo, pagamos un precio justo al campesino.} « Dans le commerce équitable, nous payons un prix juste au paysan. »
\item \esp{llegar puntual} : arriver à l'heure (ponctuel). \esp{El empleado llega puntual todos los días.} « L'employé arrive à l'heure tous les jours. »
\item \esp{tratar bien al cliente} : bien traiter le client. \esp{Un buen vendedor trata bien al cliente.} « Un bon vendeur traite bien le client. »
\end{itemize}
\end{exemplebox}

\courssub{Les familles de mots : observer pour multiplier}

En espagnol comme en français, beaucoup de mots se construisent autour d'une même racine. Observer ces familles permet d'apprendre trois ou quatre mots pour le prix d'un.

\begin{popfigure}
\begin{center}
\begin{tikzpicture}[node distance=1.6cm and 2.2cm,
raiz/.style={rectangle, rounded corners, draw=popblue, fill=popblueL, thick, align=center, font=\bfseries},
hijo/.style={rectangle, rounded corners, draw=poporange, fill=poporangeL, thick, align=center},
flecha/.style={-{Stealth[length=3mm]}, thick, draw=popdark}]
\node[raiz, align=center] (r){Racine\\ \esp{consum-}};
\node[hijo, above right=0.4cm and 2.6cm of r, align=center] (a){\esp{el consumo}\\ (nom : la consommation)};
\node[hijo, right=2.6cm of r, align=center] (b){\esp{el consumidor}\\ (personne : le consommateur)};
\node[hijo, below right=0.4cm and 2.6cm of r, align=center] (c){\esp{consumir}\\ (verbe : consommer)};
\draw[flecha] (r) -- (a);
\draw[flecha] (r) -- (b);
\draw[flecha] (r) -- (c);
\end{tikzpicture}
\end{center}
\legende{Famille de mots autour de la racine \esp{consum-}}
\end{popfigure}

\begin{propriete}[Suffixes productifs des familles de mots]
\begin{itemize}
\item Racine \cle{consum-} : \esp{consumo} (nom, suffixe \esp{-o}) ; \esp{consumidor} (personne qui agit, suffixe \esp{-dor}, comme \esp{trabajar} $\rightarrow$ \esp{trabajador}) ; \esp{consumir} (verbe infinitif en \esp{-ir}).
\item Racine \cle{ahorr-} : \esp{el ahorro} (nom d'action/résultat, suffixe \esp{-o}) ; \esp{ahorrar} (verbe) ; \esp{ahorrador/a} (agent, suffixe \esp{-dor/a}).
\item Racine \cle{honr-} : \esp{la honestidad} (nom abstrait, suffixe \esp{-dad}, correspondant au français « -té ») ; \esp{honrado/a} (adjectif, participe passé adjectivé de \esp{honrar}).
\end{itemize}
\cle{Règle mnémotechnique} : Le suffixe espagnol \esp{-dad} (féminin) correspond très souvent au français « -té » : \esp{honestidad} (honnêteté), \esp{responsabilidad} (responsabilité), \esp{puntualidad} (ponctualité), \esp{calidad} (qualité), \esp{verdad} (vérité), \esp{libertad} (liberté). Le suffixe \esp{-dor/-dora} (masc./fém.) désigne l'agent : \esp{trabajador}, \esp{consumidor}, \esp{vendedor}, \esp{ahorrador}.
\end{propriete}

\courssub{Tableau récapitulatif des quatre champs lexicaux}

\begin{center}
\begin{tabularx}{\linewidth}{|>{\bfseries}l|X|X|X|}
\hline
\rowcolor{popblueL} \textbf{Champ} & \textbf{Mot espagnol} & \textbf{Traduction} & \textbf{Exemple contextualisé} \\
\hline
Consommation & el consumidor / la consumidora & le consommateur / la consommatrice & \esp{El consumidor responsable elige bien.} \\
\hline
Consommation & comprar & acheter & \esp{Debemos comprar menos y mejor.} \\
\hline
Publicité & la publicidad & la publicité (commerciale) & \esp{La publicidad incita al derroche.} \\
\hline
Publicité & el anuncio & l'annonce, le spot & \esp{Vi un anuncio de productos locales.} \\
\hline
Gaspillage & el derroche & le gaspillage & \esp{Evitemos el derroche de agua.} \\
\hline
Gaspillage & derrochar & gaspiller & \esp{No derroches tu dinero.} \\
\hline
Économie resp. & el ahorro & l'épargne, l'économie & \esp{El ahorro permite invertir.} \\
\hline
Économie resp. & el comercio justo & le commerce équitable & \esp{El comercio justo ayuda al campesino.} \\
\hline
Économie resp. & reciclar / reutilizar & recycler / réutiliser & \esp{Reciclar es cuidar el medio ambiente.} \\
\hline
Éthique prof. & la honestidad & l'honnêteté & \esp{La honestidad genera confianza.} \\
\hline
Éthique prof. & honrado/a & honnête & \esp{Es un comerciante honrado.} \\
\hline
Éthique prof. & la puntualidad & la ponctualité & \esp{La puntualidad es respeto.} \\
\hline
Éthique prof. & el respeto al cliente & le respect du client & \esp{El respeto al cliente es ley.} \\
\hline
\end{tabularx}
\end{center}

\courssub{Texte de réinvestissement : vocabulaire en contexte}

\begin{textebox}[Una vendedora honrada en el mercado de Douala]
Doña Carmen tiene una pequeña tienda en el mercado de Douala. Ella es una vendedora honrada y responsable. Cada mañana llega puntual a las seis y abre su tienda con una sonrisa. Trata bien al cliente y nunca engaña con los precios. Por eso, los consumidores tienen confianza en ella y vuelven cada semana.

Carmen vende productos locales: plátanos, tomates, cacao y café de los campesinos de la región. Ella paga un precio justo a los productores porque cree en el comercio justo. «Si el campesino gana bien, yo también gano», dice siempre.

En su tienda no hay derroche: Carmen reutiliza las cajas de cartón y recicla las bolsas de plástico. «Debemos cuidar el medio ambiente y evitar el consumo innecesario», explica a sus hijos. Gracias al ahorro, ha comprado una nevera para conservar la fruta fresca.

Carmen no gasta dinero en publicidad: su mejor anuncio es la calidad de sus productos y su honestidad. Para ella, la ética profesional es sencilla: ser honrado, llegar puntual, respetar al cliente y ofrecer calidad. «El cliente siempre tiene razón», repite con una sonrisa. Sus hijos quieren aprender su oficio porque, gracias a su ejemplo, comprenden que comprar menos y comprar mejor es posible, y que la honestidad es la base del comercio.
\end{textebox}

\textbf{Glossaire :} \esp{la tienda} : la boutique ; \esp{engañar} : tromper ; \esp{el plátano} : la banane plantain ; \esp{el campesino} : le paysan ; \esp{la caja de cartón} : le carton ; \esp{la nevera} : le réfrigérateur ; \esp{conservar} : conserver ; \esp{el oficio} : le métier ; \esp{la sonrisa} : le sourire ; \esp{ganar} : gagner (de l'argent) ; \esp{evitar} : éviter.

\textbf{Questions de compréhension (réinvestissement) :}
\begin{enumerate}
\item ¿A qué hora llega Carmen al mercado? \esp{Llega a las seis.}
\item ¿Qué productos vende en su tienda? \esp{Vende productos locales: plátanos, tomates, cacao y café.}
\item ¿Por qué paga un precio justo a los campesinos? \esp{Porque cree en el comercio justo.}
\item ¿Qué hace Carmen para evitar el derroche? \esp{Reutiliza cajas y recicla bolsas de plástico.}
\item ¿Cuál es el mejor anuncio de Carmen? \esp{La calidad de sus productos y su honestidad.}
\end{enumerate}

\courssub{Exercice résolu et entraînement}

\begin{exoresolu}[Compléter avec le mot du champ lexical qui convient]
Énoncé — Complétez chaque phrase avec le mot espagnol approprié parmi : \esp{ahorro}, \esp{derroche}, \esp{honrado}, \esp{publicidad}, \esp{puntualidad}, \esp{productos locales}.
\begin{enumerate}
\item Un vendedor \ldots\ldots\ldots no engaña a sus clientes.
\item Tirar la comida es un \ldots\ldots\ldots
\item Gracias al \ldots\ldots\ldots, he comprado un teléfono.
\item La \ldots\ldots\ldots de la televisión es muy cara.
\item Comprar \ldots\ldots\ldots ayuda a los campesinos del Camerún.
\item La \ldots\ldots\ldots es importante en el trabajo: llega siempre a la hora.
\end{enumerate}
\tcblower
Solution détaillée :
\begin{enumerate}
\item \esp{honrado} : un vendeur \textbf{honnête} ne trompe pas ses clients (adjectif de la famille \esp{honestidad}, accord au masculin singulier).
\item \esp{derroche} : jeter la nourriture est un \textbf{gaspillage} (nom masculin, sujet de \esp{es}).
\item \esp{ahorro} : grâce à l'\textbf{épargne}, j'ai acheté un téléphone (nom masculin, après \esp{al} = \esp{a + el}).
\item \esp{publicidad} : la \textbf{publicité} de la télévision est très chère (sens commercial, féminin singulier).
\item \esp{productos locales} : acheter des \textbf{produits locaux} aide les paysans du Cameroun (nom masculin pluriel + adjectif accordé).
\item \esp{puntualidad} : la \textbf{ponctualité} est importante au travail (nom féminin en \esp{-dad}, sujet de \esp{es}).
\end{enumerate}
\end{exoresolu}

\begin{exercice}[À vous de jouer : traduction et manipulation]
\begin{enumerate}
\item Traduisez en espagnol :
\begin{enumerate}
\item Le consommateur achète des produits locaux.
\item L'honnêteté est la base du commerce.
\item Nous devons recycler et réutiliser pour protéger l'environnement.
\item Un bon employé arrive à l'heure et traite bien le client.
\end{enumerate}
\item Trouvez le mot intrus dans chaque série et justifiez en français :
\begin{enumerate}
\item \esp{consumo / consumidor / consumir / publicidad}
\item \esp{ahorrar / ahorro / derroche / ahorrador}
\item \esp{honrado / puntualidad / honestidad / responsabilidad}
\end{enumerate}
\item Construisez une phrase complète et correcte avec chacune de ces expressions :
\esp{pagar un precio justo}, \esp{llegar puntual}, \esp{tratar bien al cliente}.
\end{enumerate}
\end{exercice}

\begin{corrige}[Exercice « À vous de jouer »]
\begin{enumerate}
\item a) \esp{El consumidor compra productos locales.} b) \esp{La honestidad es la base del comercio.} c) \esp{Debemos reciclar y reutilizar para proteger el medio ambiente.} d) \esp{Un buen empleado llega puntual y trata bien al cliente.}
\item a) L'intrus est \esp{publicidad} (champ publicité/gaspillage) ; les trois autres appartiennent à la famille morphologique \esp{consum-}. b) L'intrus est \esp{derroche} (sémantique : gaspillage, antonyme d'épargne) ; les trois autres appartiennent à la famille \esp{ahorr-}. c) L'intrus est \esp{honrado} (adjectif) ; les trois autres sont des noms féminins en \esp{-dad} (qualités abstraites).
\item Exemples possibles : \esp{En el comercio justo, pagamos un precio justo al campesino.} ; \esp{La secretaria llega puntual a la oficina todos los días.} ; \esp{Doña Carmen trata bien al cliente y por eso tiene muchos clientes fieles.}
\end{enumerate}
\end{corrige}

\begin{aretenir}[L'essentiel du lexique — Fiche de révision]
\begin{itemize}
\item \textbf{Quatre champs lexicaux} à maîtriser :
\begin{itemize}
\item Consommation : \esp{el consumo, el consumidor/la consumidora, consumir, la compra, comprar, el producto}.
\item Publicité \& Gaspillage : \esp{la publicidad, el anuncio, el derroche, derrochar, innecesario, la basura}.
\item Économie responsable : \esp{el ahorro, ahorrar, el ahorrador, el producto local, el comercio justo, reutilizar, reciclar, el medio ambiente, cuidar}.
\item Éthique professionnelle : \esp{la honestidad, honrado/a, la responsabilidad, la puntualidad, el respeto al cliente, la calidad, la confianza, engañar, cumplir}.
\end{itemize}
\item \textbf{Morphologie} : suffixe \esp{-dad} = fr. « -té » ; suffixe \esp{-dor/-dora} = agent (celui qui fait).
\item \textbf{Pièges majeurs} :
\begin{itemize}
\item \esp{honestidad} : h muet, féminin ; adjectif = \esp{honrado/a} (pas \esp{honesto} à ce niveau).
\item \esp{publicidad} = publicité commerciale seulement.
\item \esp{el medio ambiente} : masculin (accord : \esp{el medio ambiente \textbf{es} importante}).
\item \esp{hacer la compra} = faire les courses (quotidien) ; \esp{hacer una compra} = faire un achat.
\end{itemize}
\item \textbf{Expression culturelle} : \esp{El cliente siempre tiene razón} (\esp{tener razón} = avoir raison).
\end{itemize}
\end{aretenir}

\coursec{Grammaire outillée : Obligation et interdiction}

Le texte support utilise \esp{Debemos cuidar...} (« Nous devons protéger... »), \esp{Debemos consumir menos...} (« Nous devons consommer moins... »). Pour exprimer des règles d'éthique professionnelle ou des conseils de consommation responsable, l'espagnol dispose de plusieurs structures. C'est un point clé pour l'expression écrite (règles de vie, chartes, conseils) et pour le thème.

\courssub{Deber + infinitif : le devoir moral, l'obligation interne}

\begin{propriete}[Deber + infinitif]
\esp{Deber} exprime un devoir moral, une obligation que l'on s'impose à soi-même ou qui relève de la conscience, de l'éthique. C'est le « doit » de la responsabilité.
\begin{center}
\begin{tabular}{|l|l|l|}
\hline
\rowcolor{popblueL} \textbf{Personne} & \textbf{Conjugaison (Présent)} & \textbf{Traduction} \\
\hline
yo & debo & je dois \\
\hline
tú & debes & tu dois \\
\hline
él/ella/usted & debe & il/elle doit / vous devez \\
\hline
nosotros/as & debemos & nous devons \\
\hline
vosotros/as & debéis & vous devez \\
\hline
ellos/ellas/ustedes & deben & ils/elles doivent / vous devez \\
\hline
\end{tabular}
\end{center}
\textbf{Exemples :}
\esp{Debemos comprar productos locales.} « Nous devons acheter des produits locaux. » (Devoir civique/éthique)
\esp{Un vendedor honrado no debe engañar al cliente.} « Un vendeur honnête ne doit pas tromper le client. » (Règle morale)
\esp{¿Debo reciclar esta botella?} « Dois-je recycler cette bouteille ? » (Demande de conseil moral)
\end{propriete}

\courssub{Tener que + infinitif : l'obligation externe, la contrainte}

\begin{propriete}[Tener que + infinitif]
\esp{Tener que} exprime une obligation imposée par les circonstances, la loi, le règlement, une tierce personne. C'est le « avoir à » / « être obligé de » du fait de contraintes extérieures. \cle{Attention} : ne pas oublier la préposition \esp{que} avant l'infinitif.
\begin{center}
\begin{tabular}{|l|l|l|}
\hline
\rowcolor{poporangeL} \textbf{Personne} & \textbf{Conjugaison (Présent)} & \textbf{Traduction} \\
\hline
yo & tengo que + inf. & j'ai à / je dois \\
\hline
tú & tienes que + inf. & tu as à / tu dois \\
\hline
él/ella/usted & tiene que + inf. & il/elle a à / il/elle doit \\
\hline
nosotros/as & tenemos que + inf. & nous avons à / nous devons \\
\hline
vosotros/as & tenéis que + inf. & vous avez à / vous devez \\
\hline
ellos/ellas/ustedes & tienen que + inf. & ils/elles ont à / ils/elles doivent \\
\hline
\end{tabular}
\end{center}
\textbf{Exemples :}
\esp{Tengo que llegar puntual a las seis.} « Je dois (je suis obligé d') arriver à l'heure à 6h. » (Règle du patron/horaire)
\esp{Los consumidores tienen que pagar el IVA.} « Les consommateurs doivent payer la TVA. » (Obligation légale)
\esp{¿Tenemos que pagar un precio justo?} « Sommes-nous obligés de payer un prix juste ? » (Contrainte contractuelle)
\end{propriete}

\begin{attention}[Deber vs Tener que : la nuance décisive]
\begin{itemize}
\item \esp{Debo estudiar} = « Je dois étudier » (c'est mon devoir d'étudiant, ma conscience me le dit).
\item \esp{Tengo que estudiar} = « Je dois étudier » (j'ai un examen demain, mes parents m'y obligent, je n'ai pas le choix).
\item Dans le texte : \esp{Debemos cuidar el medio ambiente} (devoir moral collectif) ; mais \esp{Carmen tiene que abrir la tienda a las seis} (contrainte horaire).
\end{itemize}
\end{attention}

\courssub{Hay que + infinitif : l'obligation impersonnelle, la règle générale}

\begin{propriete}[Hay que + infinitif]
\esp{Hay que} (forme impersonnelle de \esp{haver} + \esp{que}) s'utilise pour des règles générales, des consignes, des vérités générales, sans sujet précis. Cela correspond au français « il faut », « on doit ». Il n'y a \textbf{jamais} de sujet (pas de \esp{yo, tú, nosotros}...). Le verbe \esp{hay} est invariable.
\textbf{Exemples :}
\esp{Hay que reciclar el papel.} « Il faut recycler le papier. / On doit recycler le papier. »
\esp{En este mercado, hay que pagar en efectivo.} « Dans ce marché, il faut payer en espèces. »
\esp{Para ser honrado, hay que decir la verdad.} « Pour être honnête, il faut dire la vérité. »
\end{propriete}

\courssub{L'interdiction : No se debe / No hay que / Está prohibido}

\begin{propriete}[Exprimer l'interdiction]
Trois structures principales au niveau A2/B1 :
\begin{enumerate}
\item \esp{No se debe + infinitif} : « Il ne faut pas / On ne doit pas » (interdiction morale, conseil négatif, impersonnel \textbf{se} passif). \esp{No se debe derrochar agua.} « Il ne faut pas gaspiller l'eau. »
\item \esp{No hay que + infinitif} : « Il ne faut pas / Il n'est pas nécessaire de » (interdiction ou absence d'obligation selon contexte, impersonnel). \esp{No hay que engañar al cliente.} « Il ne faut pas tromper le client. »
\item \esp{Está prohibido + infinitif / nom} : « Il est interdit de... » (interdiction formelle, règlement, panneau). \esp{Está prohibido tirar basura en la calle.} « Il est interdit de jeter des ordures dans la rue. » \esp{Está prohibido el paso.} « Passage interdit. »
\end{enumerate}
\end{propriete}

\begin{attention}[No se debe vs No hay que]
\esp{No se debe} insiste sur le \textbf{devoir moral} (c'est mal). \esp{No hay que} insiste sur la \textbf{règle générale} (ce n'est pas la procédure). Souvent interchangeables pour l'interdiction, mais \esp{No se debe} est plus fort éthiquement.
\cle{Piège} : Ne dites pas \esp{*No hay que debe} ou \esp{*No se tiene que}. Structure fixe : \esp{No + se debe/hay que + infinitif}.
\end{attention}

\courssub{Tableau récapitulatif : Choix de la structure}

\begin{center}
\begin{tabularx}{\linewidth}{|X|X|X|}
\hline
\rowcolor{popblueL} \textbf{Nuance} & \textbf{Structure} & \textbf{Exemple contexte éthique / consommation} \\
\hline
Devoir moral, conseil éthique & \esp{Deber + inf.} & \esp{Debemos apoyar el comercio justo.} \\
\hline
Obligation externe, contrainte, règle & \esp{Tener que + inf.} & \esp{El vendedor tiene que dar factura.} \\
\hline
Règle générale, consigne impersonnelle & \esp{Hay que + inf.} & \esp{Hay que reciclar en casa.} \\
\hline
Interdiction morale / conseil négatif & \esp{No se debe + inf.} & \esp{No se debe comprar marfil.} \\
\hline
Interdiction formelle / panneau / loi & \esp{Está prohibido + inf.} & \esp{Está prohibido vender tabaco a menores.} \\
\hline
\end{tabularx}
\end{center}

\courssub{Exercices d'application grammaticale}

\begin{exoresolu}[Choisir la structure d'obligation/interdiction la plus adaptée]
Énoncé — Complétez avec \esp{debo}, \esp{tengo que}, \esp{hay que}, \esp{no se debe}, \esp{está prohibido} (conjuguez \esp{deber} et \esp{tener} si nécessaire).
\begin{enumerate}
\item \ldots\ldots\ldots (yo) reciclar las botellas de plástico para proteger el planeta. (Conseil moral personnel)
\item \ldots\ldots\ldots (nosotros) pagar la factura de la electricidad mañana. (Contrainte externe, échéance)
\item En la escuela \ldots\ldots\ldots llegar puntual a las 7:30. (Règle générale, règlement intérieur)
\item \ldots\ldots\ldots tirar basura en el río. (Interdiction formelle, panneau écologique)
\item Un comerciante honrado \ldots\ldots\ldots (él) mentir sobre los precios. (Interdiction morale forte)
\end{enumerate}
\tcblower
Solution détaillée :
\begin{enumerate}
\item \esp{Debo} (devoir moral personnel : \esp{yo debo}).
\item \esp{Tenemos que} (obligation externe, échéance : \esp{nosotros tenemos que}).
\item \esp{hay que} (règle générale impersonnelle, pas de sujet).
\item \esp{Está prohibido} (interdiction formelle, affichage public).
\item \esp{no debe} (interdiction morale, devoir professionnel : \esp{él no debe}).
\end{enumerate}
\end{exoresolu}

\begin{exercice}[Entraînement : Obligation et interdiction au marché]
\begin{enumerate}
\item Mettez au bon forme (présent) : a) Yo \ldots\ldots\ldots (deber) comprar menos ropa. b) Los vendedores \ldots\ldots\ldots (tener que) pesar bien la fruta. c) Para ahorrar, \ldots\ldots\ldots (hay que) apagar la luz. d) \ldots\ldots\ldots (estar prohibido) fumar en la tienda. e) Tú \ldots\ldots\ldots (no / deber) derrochar comida.
\item Traduisez en espagnol :
\begin{enumerate}
\item Il faut (règle générale) respecter le client.
\item Je dois (obligation morale) être honnête.
\item Nous sommes obligés (contrainte) de payer des impôts.
\item Il est interdit de gaspiller l'eau.
\item On ne doit pas (conseil) acheter des produits inutiles.
\end{enumerate}
\end{enumerate}
\end{exercice}

\begin{corrige}[Exercice « Entraînement »]
\begin{enumerate}
\item a) \esp{debo} b) \esp{tienen que} c) \esp{hay que} d) \esp{Está prohibido} e) \esp{no debes}
\item a) \esp{Hay que respetar al cliente.} b) \esp{Debo ser honrado.} (ou \esp{honesta}) c) \esp{Tenemos que pagar impuestos.} d) \esp{Está prohibido derrochar el agua.} e) \esp{No se debe comprar productos innecesarios.} (ou \esp{No hay que comprar...})
\end{enumerate}
\end{corrige}

\coursec{Méthode : Commenter un texte argumentatif}

Le commentaire de texte est une épreuve majeure du Baccalauréat (épreuve écrite de compréhension et expression). Le texte de cette leçon (\esp{Comprar menos, comprar mejor}) est un texte \textbf{argumentatif} : il défend une thèse (la consommation responsable) contre une antithèse (la société de consommation/gaspillage).

\begin{methode}[Étapes du commentaire de texte argumentatif (niveau Première)]
\begin{enumerate}
\item \textbf{Lecture active (10 min) :} Lisez 2 fois. 1ère : compréhension globale (sujet, thèse). 2nde : soulignez connecteurs, arguments, exemples, champ lexical, temps verbaux. Notez la structure (paragraphes).
\item \textbf{Introduction (rédaction soignée) :}
\begin{itemize}
\item \textit{Accroche} : Situez le thème (ex. \esp{El consumo responsable es un tema de actualidad...}).
\item \textit{Présentation du texte} : Titre (entre guillemets), auteur (si connu, sinon « texto anónimo »), source, date, thème principal.
\item \textit{Thèse} : Résumez l'idée principale en une phrase. \esp{El autor defiende la idea de que...}
\item \textit{Annonce du plan} : \esp{En primer lugar analizaremos...; luego...; finalmente...}
\end{itemize}
\item \textbf{Développement (2 ou 3 parties équilibrées) :} Une partie = un argument principal du texte.
\begin{itemize}
\item \textit{Idée directrice} : \esp{El primer argumento del autor es...}
\item \textit{Citation} : Intégrez un court extrait (max 1 ligne) entre guillemets espagnols \esp{« ... »}.
\item \textit{Analyse} : Expliquez la citation, le vocabulaire clé, la structure de la phrase, l'effet produit. Liez au thème.
\item \textit{Exemple personnel / Contexte camerounais} : Montrez votre culture (marché local, cacao, bendskin, etc.).
\end{itemize}
\item \textbf{Conclusion (synthèse + ouverture) :}
\begin{itemize}
\item \textit{Bilan} : \esp{En definitiva, el texto nos demuestra que...}
\item \textit{Ouverture} : Question, citation, lien avec un autre texte, actualité (COP, Objectifs Développement Durable, commerce équitable au Cameroun).
\end{itemize}
\item \textbf{Relecture (5 min) :} Vérifiez accords, accents, \esp{ser/estar}, \esp{por/para}, subjonctif après \esp{es importante que}, connecteurs (\esp{por tanto, sin embargo, además, en conclusión}).
\end{enumerate}
\end{methode}

\begin{exemplebox}[Modèle d'introduction pour « Comprar menos, comprar mejor »]
« \esp{El consumo responsable es un tema de actualidad urgente en un mundo dominado por la publicidad y el derroche.} » (Accroche). Le texte \esp{« Comprar menos, comprar mejor »}, de source pédagogique, traite de la nécessité de changer nos habitudes de consommation. \esp{La tesis principal es que debemos consumir menos pero mejor, apoyando el comercio local y justo, y adoptando una ética profesional honrada.} (Thèse). \esp{En primer lugar, analizaremos la crítica de la sociedad de consumo; luego, las propuestas de consumo responsable; finalmente, el papel de la ética profesional.} (Plan).
\end{exemplebox}

\begin{attention}[Erreurs fréquentes au commentaire]
\begin{itemize}
\item \textbf{Résumé au lieu d'analyse :} Ne racontez pas le texte, \textbf{analysez}-le (pourquoi ce mot ? quel effet ?).
\item \textbf{Citations trop longues :} Maximum une ligne, intégrée dans votre phrase.
\item \textbf{Oubli de l'ouverture :} La conclusion ne doit pas être une simple répétition.
\item \textbf{Langue :} Écrivez en espagnol. Utilisez les connecteurs logiques (\esp{por lo tanto, no obstante, además, en primer lugar}).
\end{itemize}
\end{attention}

\coursec{Production guidée : Rédiger une charte de consommation responsable}

\courssub{Consigne}

Vous êtes délégués de la classe de Première A du Lycée de Yaoundé. Vous devez rédiger une \textbf{Charte de la Consommation Responsable au Lycée} (150--180 mots) qui sera affichée dans les couloirs et à la cantine. Cette charte s'adresse aux élèves (\esp{tú} ou \esp{ustedes}) et au personnel.

\courssub{Plan imposé}

\begin{enumerate}
\item \textbf{Titre accrocheur} (ex. \esp{¡Compra menos, compra mejor!}).
\item \textbf{Préambule} (2--3 lignes) : Pourquoi cette charte ? (Constat : gaspillage, publicité, environnement). Utilisez \esp{Hay que...}, \esp{Es necesario...}.
\item \textbf{5 Règles d'or} (numérotées) : Utilisez \textbf{au moins 3 structures différentes} d'obligation/interdiction étudiées (\esp{Debemos, Tenemos que, Hay que, No se debe, Está prohibido}). Intégrez le vocabulaire des 4 champs lexicaux.
\item \textbf{Engagement éthique} (2 lignes) : Rappel des valeurs (honnêteté, respect, ponctualité) pour les délégués et les vendeurs de la cantine. Vocabulaire : \esp{honestidad, honrado, puntualidad, respeto}.
\item \textbf{Signature / Date}.
\end{enumerate}

\courssub{Banque de connecteurs et structures utiles}

\begin{center}
\begin{tabularx}{\linewidth}{|X|X|}
\hline
\rowcolor{popblueL} \textbf{Fonction} & \textbf{Outils linguistiques} \\
\hline
Ordonner & \esp{En primer lugar, En segundo lugar, Finalmente, Para terminar} \\
\hline
Ajouter un argument & \esp{Además, Asimismo, También, Por otra parte} \\
\hline
Opposer / Nuancer & \esp{Sin embargo, No obstante, Aunque...} \\
\hline
Conséquence & \esp{Por tanto, Por consiguiente, Así que} \\
\hline
Cause & \esp{Porque, Como, Debido a que} \\
\hline
Exemplifier & \esp{Por ejemplo, Como por ejemplo} \\
\hline
\end{tabularx}
\end{center}

\courssub{Modèle de rédaction (guide)}

\begin{textebox}[Modèle de Charte — ¡Consumamos con cabeza en el Lycée!]
\textbf{¡CONSUMO RESPONSABLE EN NUESTRO LICEO!}

Queridos compañeros: Cada día tiramos kilos de comida y papel. La publicidad nos empuja a comprar cosas innecesarias. \textbf{Hay que cambiar} esta situación ya.

\textbf{Nuestras 5 Reglas de Oro:}
1. \textbf{Debemos comprar productos locales} (plátanos, maíz, cacao) en el mercado escolar para apoyar a los campesinos cameruneses.
2. \textbf{Tenemos que llevar} nuestra propia botella de agua y tupper; \textbf{está prohibido} usar vasos de plástico de un solo uso.
3. \textbf{No se debe derrochar} la comida en la cantina: \textbf{hay que} servir solo lo que vamos a comer.
4. \textbf{Debemos reciclar} papel y plástico en los contenedores azules y amarillos del patio.
5. \textbf{Los vendedores de la cantina tienen que ser honrados}: \textbf{no deben engañar} con los precios ni el peso.

\textbf{Compromiso ético:} La \esp{honestidad}, la \esp{puntualidad} y el \esp{respeto} son la base de nuestra convivencia. Un delegado \esp{honrado} \esp{cumple su palabra}.

¡Firmado: Los Delegados de 1ère A --- Yaoundé, \today
\end{textebox}

\begin{experience}[Atelier oral : « Le marché équitable »]
\textbf{Objectif :} Réinvestir lexique + grammaire (obligation/interdiction) en interaction.
\textbf{Déroulement (15 min) :}
\begin{enumerate}
\item \textbf{Binômes} : Un élève = \esp{El campesino} (producteur de cacao/café), l'autre = \esp{El comprador} (représentant d'une coopérative de commerce équitable).
\item \textbf{Scénario} : Négociation d'un contrat annuel. Le paysan veut un prix juste, paiement à l'avance, respect de l'environnement. L'acheteur veut qualité, régularité, certification.
\item \textbf{Contraintes linguistiques} : Utiliser \textbf{au moins 4 fois} des structures d'obligation/interdiction (\esp{debo, tengo que, hay que, no se debe, está prohibido}) et \textbf{6 mots} du lexique de la leçon (\esp{comercio justo, precio justo, honrado, puntualidad, medio ambiente, derroche, ahorro, calidad}).
\item \textbf{Présentation} : 3 binômes jouent la scène devant la classe. La classe note les structures entendues.
\end{enumerate}
\end{experience}

\begin{aretenir}[Bilan de la leçon EA15 — Ce qu'il faut maîtriser pour le Bac]
\begin{itemize}
\item \textbf{Lexique} : 4 champs (consommation, publicité/gaspillage, économie responsable, éthique pro) + familles \esp{consum-}, \esp{ahorr-}, \esp{honr-} + suffixes \esp{-dad}, \esp{-dor}.
\item \textbf{Grammaire} : Distinguer et employer correctement \esp{deber} (devoir moral), \esp{tener que} (contrainte), \esp{hay que} (règle générale), \esp{no se debe / no hay que / está prohibido} (interdictions).
\item \textbf{Méthode} : Structure du commentaire argumentatif (Intro thèse plan, Développement citation+analyse, Conclusion synthèse+ouverture).
\item \textbf{Production} : Rédiger une charte / article d'opinion / lettre formelle en mobilisant lexique + grammaire + connecteurs.
\item \textbf{Culture} : \esp{El cliente siempre tiene razón}, commerce équitable (cacao/café Cameroun/Guinée Équatoriale), marchés locaux (Mvog-Mbi, Douala).
\end{itemize}
\end{aretenir}

\coursec{Grammaire outillée : obligation et interdiction}

Dans la section précédente, nous avons constitué le lexique thématique de la consommation responsable et de l'éthique professionnelle : \esp{consumo}, \esp{consumidor}, \esp{ahorro}, \esp{comercio justo}, \esp{honestidad}... Mais un texte argumentatif comme « Comprar menos, comprar mejor » ne se contente pas de nommer les choses : il dit ce qu'il \emph{faut} faire et ce qu'il ne faut \emph{pas} faire. C'est précisément le rôle des tournures d'obligation et d'interdiction que nous allons maintenant observer, comprendre et manipuler. Cette grammaire est indispensable pour exprimer des règles d'éthique professionnelle, compétence visée par notre leçon.

\courssub{Observation : que disent les phrases du texte ?}

Reprenons trois phrases tirées du texte support et observons-les attentivement.

\begin{exemplebox}[Phrases extraites du texte support]
\begin{enumerate}
\item \esp{Debemos consumir de manera responsable.} « Nous devons consommer de manière responsable. »
\item \esp{Tenemos que ahorrar.} « Nous devons économiser. »
\item \esp{No se debe tirar la comida.} « On ne doit pas jeter la nourriture. »
\end{enumerate}
\end{exemplebox}

Observons la structure de chaque phrase :

\begin{itemize}
\item Phrase 1 : verbe \esp{deber} conjugué + \textbf{infinitif} directement (\esp{consumir}).
\item Phrase 2 : verbe \esp{tener} conjugué + \textbf{que} + infinitif (\esp{ahorrar}).
\item Phrase 3 : \esp{no} + \esp{se} + \esp{debe} (3\textsuperscript{e} personne du singulier) + infinitif (\esp{tirar}).
\end{itemize}

Trois constructions différentes, mais une même famille de sens : l'obligation et l'interdiction. Chacune a sa nuance propre. Distinguons-les une à une.

\courssub{Règle 1 : \esp{deber} + infinitif — le devoir moral}

\begin{propriete}[Deber + infinitif]
\textbf{Forme :} \esp{deber} conjugué + infinitif (\textbf{sans} \esp{que}).

\textbf{Emploi :} \esp{deber} + infinitif exprime le \cle{devoir moral}, l'obligation intérieure, ce que la conscience, l'éthique ou la règle professionnelle commandent. C'est l'équivalent du français « devoir » au sens moral.

\textbf{Exemple :} \esp{Debo ser honesto con los clientes.} « Je dois être honnête avec les clients. »
\end{propriete}

Conjugaison complète de \esp{deber} au présent de l'indicatif (verbe régulier en \esp{-er}) :

\begin{center}
\begin{tabular}{|l|l|l|}
\hline
\rowcolor{popblueL} \textbf{Personne} & \textbf{Deber} & \textbf{Exemple} \\
\hline
yo & debo & Debo ser puntual. \\
\hline
tú & debes & Debes respetar al cliente. \\
\hline
él / ella / usted & debe & Un profesional debe cumplir sus promesas. \\
\hline
nosotros / nosotras & debemos & Debemos consumir de manera responsable. \\
\hline
vosotros / vosotras & debéis & Debéis ahorrar agua. \\
\hline
ellos / ellas / ustedes & deben & Los comerciantes deben ser honestos. \\
\hline
\end{tabular}
\end{center}

\begin{exemplebox}[Deber + infinitif : exemples traduits]
\begin{itemize}
\item \esp{Un profesional debe cumplir sus promesas.} « Un professionnel doit tenir ses promesses. »
\item \esp{Debes decir la verdad.} « Tu dois dire la vérité. »
\item \esp{Debemos respetar el medio ambiente.} « Nous devons respecter l'environnement. »
\item \esp{Los estudiantes deben trabajar con seriedad.} « Les élèves doivent travailler avec sérieux. »
\end{itemize}
\end{exemplebox}

\begin{attention}[Erreur fréquente des francophones]
En français on dit « je dois \textbf{que}... » dans certaines tournures familières, et surtout on est tenté de calquer « avoir à » ou « devoir que ». En espagnol, \esp{debo que trabajar} est \textbf{FAUX}. Le verbe \esp{deber} s'emploie \textbf{sans} \esp{que} devant l'infinitif : \esp{debo trabajar}. À l'inverse, \esp{tener} exige toujours \esp{que} : \esp{tengo que trabajar}. Retenez le contraste : \esp{debo trabajar} mais \esp{tengo que trabajar}.
\end{attention}

\courssub{Règle 2 : \esp{tener que} + infinitif — l'obligation pratique}

\begin{propriete}[Tener que + infinitif]
\textbf{Forme :} \esp{tener} conjugué + \esp{que} + infinitif.

\textbf{Emploi :} \esp{tener que} + infinitif exprime l'\cle{obligation pratique}, la nécessité concrète imposée par les circonstances extérieures (l'horaire, la loi, la situation). C'est l'équivalent du français « devoir / être obligé de / avoir à ».

\textbf{Exemple :} \esp{Tengo que llegar a tiempo al trabajo.} « Je dois arriver à l'heure au travail. »
\end{propriete}

Attention : \esp{tener} est un verbe \textbf{irrégulier} au présent (radical \esp{teng-} à la 1\textsuperscript{re} personne, diphtongaison \esp{e} → \esp{ie} aux trois personnes du singulier et à la 3\textsuperscript{e} personne du pluriel).

\begin{center}
\begin{tabular}{|l|l|l|}
\hline
\rowcolor{poporangeL} \textbf{Personne} & \textbf{Tener que} & \textbf{Exemple} \\
\hline
yo & tengo que & Tengo que llegar a tiempo. \\
\hline
tú & tienes que & Tienes que atender bien a los clientes. \\
\hline
él / ella / usted & tiene que & Ella tiene que abrir la tienda a las ocho. \\
\hline
nosotros / nosotras & tenemos que & Tenemos que ahorrar. \\
\hline
vosotros / vosotras & tenéis que & Tenéis que comprar productos locales. \\
\hline
ellos / ellas / ustedes & tienen que & Tienen que pagar sus impuestos. \\
\hline
\end{tabular}
\end{center}

\begin{exemplebox}[Tener que + infinitif : exemples traduits]
\begin{itemize}
\item \esp{Tenemos que ahorrar energía en casa.} « Nous devons économiser l'énergie à la maison. »
\item \esp{El vendedor tiene que dar un recibo al cliente.} « Le vendeur doit donner un reçu au client. »
\item \esp{¿Tienes que trabajar los sábados?} « Dois-tu travailler les samedis ? »
\item \esp{En Douala, los comerciantes tienen que levantarse muy temprano.} « À Douala, les commerçants doivent se lever très tôt. »
\end{itemize}
\end{exemplebox}

\courssub{Règle 3 : \esp{no se debe} + infinitif — l'interdiction générale}

\begin{propriete}[No se debe + infinitif]
\textbf{Forme :} \esp{no} + \esp{se} + \esp{debe} (toujours à la 3\textsuperscript{e} personne du singulier) + infinitif.

\textbf{Emploi :} cette tournure \cle{impersonnelle} exprime l'interdiction ou le conseil général, valable pour tout le monde : « on ne doit pas », « il ne faut pas ». On la trouve sur les panneaux, dans les règlements et dans les textes de morale ou d'éthique.

\textbf{Exemple :} \esp{No se debe mentir.} « On ne doit pas mentir. »
\end{propriete}

\begin{exemplebox}[No se debe + infinitif : exemples traduits]
\begin{itemize}
\item \esp{No se debe tirar la comida.} « On ne doit pas jeter la nourriture. »
\item \esp{No se debe engañar a los consumidores.} « On ne doit pas tromper les consommateurs. »
\item \esp{No se debe llegar tarde al trabajo.} « On ne doit pas arriver en retard au travail. »
\item \esp{No se debe malgastar el agua.} « On ne doit pas gaspiller l'eau. »
\end{itemize}
\end{exemplebox}

Pour renforcer l'interdiction (panneaux, règlements stricts), l'espagnol utilise aussi \esp{estar prohibido} + infinitif : \esp{Está prohibido fumar.} « Il est interdit de fumer. »

\courssub{Complément utile : \esp{hay que} + infinitif — « il faut »}

\begin{propriete}[Hay que + infinitif]
Le français « il faut » se traduit le plus souvent par \esp{hay que} + infinitif (tournure impersonnelle, invariable) ou par \esp{se debe} + infinitif.

\textbf{Exemple :} \esp{Hay que respetar al cliente.} « Il faut respecter le client. »

\esp{Hay que} exprime la \cle{nécessité générale}, sans désigner personne en particulier. C'est une tournure très fréquente à l'oral et très utile au Bac.
\end{propriete}

\begin{center}
\begin{tabularx}{\linewidth}{|X|X|X|}
\hline
\rowcolor{popblueL} \textbf{Français} & \textbf{Español} & \textbf{Remarque} \\
\hline
Je dois (devoir moral) & debo + inf. & sans \esp{que} \\
\hline
Je dois (obligation concrète) & tengo que + inf. & avec \esp{que} obligatoire \\
\hline
Il faut & hay que + inf. & impersonnel, invariable \\
\hline
On doit / il faut & se debe + inf. & impersonnel, 3\textsuperscript{e} pers. sing. \\
\hline
On ne doit pas / il ne faut pas & no se debe + inf. & interdiction générale \\
\hline
Il est interdit de & está prohibido + inf. & interdiction forte (règlement) \\
\hline
\end{tabularx}
\end{center}

\courssub{Nuances : quelle obligation choisir ?}

Les trois tournures ne sont pas interchangeables : elles forment une véritable échelle, du simple conseil à l'interdiction.

\begin{itemize}
\item \esp{se debe} : conseil général, morale collective (« on doit »).
\item \esp{deber} : devoir moral, obligation intérieure, éthique (« je dois être honnête »).
\item \esp{tener que} : obligation forte, imposée par les circonstances (« je dois arriver à 8 h, sinon je perds mon emploi »).
\item \esp{no se debe} / \esp{está prohibido} : interdiction.
\end{itemize}

\begin{popfigure}
\begin{center}
\begin{tikzpicture}[scale=1]
\draw[-{Stealth[length=3mm]}, line width=1.2pt, popdark] (0,0) -- (12.5,0);
\foreach \x/\c/\t/\s in {
  1.2/popgreen/Conseil/{se debe + inf.},
  4.4/popblue/Obligation morale/{deber + inf.},
  7.9/poporange/Obligation forte/{tener que + inf.},
  11.2/poppink/Interdiction/{no se debe / está prohibido}}{
  \fill[\c] (\x,0) circle (0.12);
  \node[above, align=center, font=\small\bfseries, text=\c] at (\x,0.25) {\t};
  \node[below, align=center, font=\footnotesize, text=popdark] at (\x,-0.25) {\s};
}
\node[font=\footnotesize\itshape, text=popmuted] at (6.25,-1.1) {Force croissante de la contrainte};
\end{tikzpicture}
\end{center}
\legende{L'échelle de l'obligation en espagnol : du conseil à l'interdiction.}
\end{popfigure}

Comparons sur un même exemple :

\begin{itemize}
\item \esp{Se debe ahorrar agua.} « On doit économiser l'eau » (principe général).
\item \esp{Debemos ahorrar agua.} « Nous devons économiser l'eau » (devoir que nous ressentons).
\item \esp{Tenemos que ahorrar agua porque hay sequía.} « Nous devons économiser l'eau parce qu'il y a la sécheresse » (nécessité concrète).
\item \esp{No se debe malgastar el agua.} « On ne doit pas gaspiller l'eau » (interdiction).
\end{itemize}

\courssub{Tableau de conjugaison récapitulatif}

\begin{popfigure}
\begin{center}
\begin{tikzpicture}[font=\small]
\node[draw=popblue, fill=popblueL, rounded corners, align=center, minimum width=5.6cm, minimum height=0.9cm] (t1) at (0,0) {\textbf{DEBER} (régulier en -er)};
\node[draw=popblue, rounded corners, align=left, text width=5.4cm, below=0.15cm of t1] (c1) {debo \quad debes \quad debe\\debemos \quad debéis \quad deben};
\node[draw=popdark, fill=popcream, rounded corners, align=center, minimum width=5.6cm, below=0.15cm of c1] (s1) {+ \textbf{infinitif} (sin \emph{que})};

\node[draw=poporange, fill=poporangeL, rounded corners, align=center, minimum width=5.6cm, minimum height=0.9cm] (t2) at (7.2,0) {\textbf{TENER QUE} (irrégulier)};
\node[draw=poporange, rounded corners, align=left, text width=5.4cm, below=0.15cm of t2] (c2) {tengo que \quad tienes que \quad tiene que\\tenemos que \quad tenéis que \quad tienen que};
\node[draw=popdark, fill=popcream, rounded corners, align=center, minimum width=5.6cm, below=0.15cm of c2] (s2) {+ \textbf{infinitif} (con \emph{que})};
\end{tikzpicture}
\end{center}
\legende{Les deux conjugaisons à connaître par cœur : la structure « verbe conjugué + infinitif » est identique, mais \esp{deber} refuse \esp{que} et \esp{tener} l'exige.}
\end{popfigure}

\courssub{Mise en contexte : mini-texte d'application}

\begin{textebox}[La ética de un buen comerciante]
En el mercado de Yaoundé, don Manuel es un comerciante respetado por todos. ¿Cuál es su secreto? Él lo explica con sencillez: « Un profesional debe cumplir sus promesas. Yo debo ser honesto con los clientes: nunca engaño con la balanza y siempre doy un precio justo. Además, tengo que llegar a tiempo cada mañana, porque los clientes me esperan a las siete. También tenemos que respetar a los demás vendedores y mantener limpio el mercado. Hay que recordar algunas reglas básicas: no se debe mentir sobre la calidad de los productos, no se debe tirar la comida en mal estado delante de los clientes y no se debe subir los precios sin razón. Si un cliente no está contento, tengo que escucharlo y debo cambiar el producto. El comercio justo no es solamente una palabra: es una manera de trabajar cada día. Los jóvenes que quieren ser comerciantes deben aprender que la honestidad es más importante que el dinero. Con la publicidad se puede vender mucho, pero sin ética se pierde la confianza de la gente. Y sin confianza, no hay clientes. »

\textbf{Glossaire :} \esp{la balanza} = la balance ; \esp{engañar} = tromper ; \esp{subir los precios} = augmenter les prix ; \esp{la confianza} = la confiance ; \esp{sencillez} = simplicité ; \esp{en mal estado} = en mauvais état ; \esp{a tiempo} = à l'heure / ponctuel.
\end{textebox}

Questions de repérage grammatical : (1) Relevez toutes les tournures d'obligation et d'interdiction du texte. (2) Pour chacune, dites s'il s'agit d'un devoir moral, d'une obligation pratique ou d'une interdiction. (3) Pourquoi don Manuel utilise-t-il \esp{tengo que llegar} et non \esp{debo llegar} ?

\courssub{Exercice résolu}

\begin{exoresolu}[Obligation ou interdiction ?]
Complétez avec la bonne forme : \esp{deber}, \esp{tener que} ou \esp{no se debe} (au présent).

1. Los consumidores \dots\ comprar productos locales. (devoir moral)
2. Yo \dots\ pagar al vendedor antes de irme. (obligation concrète)
3. \dots\ derrochar el dinero. (interdiction générale)
4. Nosotros \dots\ llegar puntuales a la cita. (nécessité pratique)
5. Un empleado \dots\ ser respetuoso con los clientes. (éthique professionnelle)
\tcblower
\textbf{Solution détaillée :}

1. \esp{Los consumidores \textbf{deben} comprar productos locales.} — Il s'agit d'un devoir moral, d'un principe de consommation responsable : \esp{deber} + infinitif, 3\textsuperscript{e} personne du pluriel (\esp{deben}), sans \esp{que}.

2. \esp{Yo \textbf{tengo que} pagar al vendedor antes de irme.} — Obligation concrète, situation pratique : \esp{tener que} + infinitif ; 1\textsuperscript{re} personne du singulier du verbe irrégulier : \esp{tengo}, avec \esp{que}.

3. \esp{\textbf{No se debe} derrochar el dinero.} — Interdiction générale valable pour tous : tournure impersonnelle \esp{no se debe} + infinitif.

4. \esp{Nosotros \textbf{tenemos que} llegar puntuales a la cita.} — Nécessité pratique (un rendez-vous, un horaire) : \esp{tenemos que} + infinitif ; pas de diphtongaison à \esp{nosotros}.

5. \esp{Un empleado \textbf{debe} ser respetuoso con los clientes.} — Règle d'éthique professionnelle, devoir moral : \esp{debe} + infinitif, sans \esp{que}.
\end{exoresolu}

\begin{attention}[Récapitulatif des pièges]
\begin{itemize}
\item \esp{debo que trabajar} : FAUX. \esp{deber} se construit \textbf{sans} \esp{que} : \esp{debo trabajar}.
\item \esp{tengo trabajar} : FAUX. \esp{tener} exige \esp{que} : \esp{tengo que trabajar}.
\item Après \esp{deber}, \esp{tener que}, \esp{hay que} et \esp{no se debe}, le verbe est toujours à l'\textbf{infinitif}, jamais conjugué : \esp{debemos consumir}, et non \esp{debemos consumimos}.
\item Ne confondez pas \esp{hay que} (il faut) et \esp{hay} (il y a) : \esp{Hay que ahorrar} = il faut économiser ; \esp{Hay muchos productos} = il y a beaucoup de produits.
\item Dans \esp{no se debe}, le verbe reste au singulier même si l'on parle de plusieurs personnes : \esp{no se debe tirar la comida}.
\item \textbf{Attention au sens :} \esp{deber de} + infinitif exprime la \emph{probabilité} ou la supposition (\esp{Debe de ser las ocho} = « Il doit être huit heures »), et non l'obligation. Ne pas confondre avec \esp{deber} + infinitif (obligation).
\end{itemize}
\end{attention}

\begin{savaistu}[« Deber », le verbe de l'argent]
Le verbe \esp{deber} ne signifie pas seulement « devoir » au sens moral : il veut dire aussi « devoir » au sens d'\emph{être redevable d'une somme} ! \esp{Le debo 5000 francos.} « Je lui dois 5000 francs. » \esp{¿Cuánto le debo?} « Combien vous dois-je ? » — phrase très utile au marché de Yaoundé ou de Douala ! D'ailleurs, le nom \esp{la deuda} (la dette) vient de la même racine latine \emph{debere}. Dans le champ lexical de l'argent et du commerce, \esp{deber} est donc un mot central : celui qui doit de l'argent est \esp{el deudor} (le débiteur). Et le participe \esp{debido} donne l'expression \esp{debido a} = « à cause de, en raison de » : \esp{El precio sube debido a la crisis.} « Le prix augmente à cause de la crise. »
\end{savaistu}

\begin{aretenir}[L'essentiel de la grammaire]
\begin{itemize}
\item \esp{deber} + infinitif (sans \esp{que}) = devoir moral : \esp{Debo ser honesto.}
\item \esp{tener que} + infinitif (avec \esp{que}) = obligation pratique : \esp{Tengo que llegar a tiempo.}
\item \esp{hay que} + infinitif = il faut (nécessité générale) : \esp{Hay que respetar al cliente.}
\item \esp{no se debe} + infinitif = on ne doit pas (interdiction générale) : \esp{No se debe mentir.}
\item Interdiction forte : \esp{está prohibido} + infinitif.
\item L'infinitif suit toujours directement ces tournures.
\end{itemize}
\end{aretenir}

Ces outils grammaticaux maîtrisés, nous sommes maintenant prêts pour la section suivante : la méthode du commentaire de texte, où il faudra repérer et analyser ces tournures d'obligation dans un document authentique.

\coursec{Méthode du commentaire de texte}

Dans la section précédente, nous avons appris à exprimer l'obligation et l'interdiction avec \esp{deber}, \esp{tener que} et \esp{no se debe}. Ces tournures ne servent pas seulement à parler : elles sont aussi des \cle{indices précieux} pour analyser un texte. Quand un auteur écrit \esp{debemos comprar menos}, il ne raconte pas une histoire : il \emph{convainc}, il \emph{argumente}. Savoir repérer et commenter ces procédés, c'est exactement ce qu'on te demandera au baccalauréat dans l'épreuve de commentaire de texte. Cette section te donne une méthode complète, étape par étape, avec un modèle rédigé et les expressions toutes prêtes.

\courssub{Qu'est-ce qu'un commentaire de texte ?}

\begin{definition}[Le commentaire de texte]
Le \cle{commentaire de texte} (\esp{comentario de texto}) est un exercice écrit qui consiste à \textbf{présenter} un texte, à \textbf{analyser} ses idées et ses procédés, puis à \textbf{donner son avis personnel argumenté}. Il ne s'agit ni de recopier le texte, ni de le résumer mot à mot : il s'agit de montrer que tu l'as \emph{compris} et que tu sais \emph{réagir}.
\end{definition}

\begin{attention}[Résumé ou commentaire ?]
Erreur classique des candidats francophones : confondre le \textbf{résumé} et le \textbf{commentaire}.
\begin{itemize}
\item Le \textbf{résumé} (\esp{resumen}) \emph{répète} le texte en plus court, sans avis personnel.
\item Le \textbf{commentaire} (\esp{comentario}) \emph{analyse} le texte (idées, procédés, intention de l'auteur) et se termine par une \emph{opinion personnelle justifiée}.
\end{itemize}
Si ta copie ne contient que ce que dit le texte, sans analyse ni avis, tu as fait un résumé : tu perds la moiti\% des points.
\end{attention}

\courssub{Les cinq étapes du commentaire}

\begin{methode}[Réussir le commentaire de texte au Bac]
\begin{enumerate}
\item \textbf{Présenter le texte} (\esp{presentación}) : indique le \emph{type de texte} (argumentativo, narrativo, descriptivo, periodístico), le \emph{thème}, l'\emph{auteur ou le narrateur} et la \emph{situation} (qui parle ? à qui ? où ?).
\item \textbf{Dégager l'idée principale} (\esp{idea principal}) : formule en une ou deux phrases ce que l'auteur veut démontrer. Demande-toi : « Que veut me faire penser l'auteur ? »
\item \textbf{Analyser les idées secondaires} (\esp{análisis}) : repère le \emph{problème} posé (ici : \esp{el derroche}, \esp{la publicidad}) et la \emph{solution} proposée (ici : \esp{el consumo responsable}, \esp{el comercio justo}). Justifie chaque idée par une citation.
\item \textbf{Relever les procédés} (\esp{procedimientos}) : opposition (\esp{en cambio}, \esp{sin embargo}), énumération, impératifs (\esp{compra}, \esp{elige}), tournures d'obligation (\esp{debemos}, \esp{no se debe}). Explique à quoi ils servent : convaincre, opposer, énumérer, obliger.
\item \textbf{Donner ton avis argumenté} (\esp{opinión personal}) : réponds aux questions \esp{¿Estás de acuerdo? ¿Por qué?} en deux ou trois phrases, avec un exemple concret (le Cameroun, ton quartier, ton école).
\end{enumerate}
\textbf{Règle d'or} : cite toujours le texte \emph{en espagnol, entre guillemets}, pour justifier chaque idée, puis traduis ou explique la citation. Une idée sans citation n'est pas prouvée ; une citation sans explication n'est pas analysée.
\end{methode}

\begin{popfigure}
\begin{center}
\begin{tikzpicture}[
node distance=0.55cm and 0.7cm,
etape/.style={rectangle, rounded corners=4pt, draw=popblue, fill=popblueL, text width=4.2cm, align=center, font=\small\bfseries, minimum height=0.9cm},
desc/.style={rectangle, draw=popline, fill=popcream, text width=5.2cm, align=center, font=\footnotesize, minimum height=0.9cm},
fleche/.style={-{Stealth[length=3mm]}, thick, poporange}
]
\node[etape] (e1) {1. Presentación};
\node[desc, right=of e1] (d1) {Tipo de texto, tema, narrador, situación};
\node[etape, below=of e1] (e2) {2. Idea principal};
\node[desc, right=of e2] (d2) {¿Qué quiere demostrar el autor?};
\node[etape, below=of e2] (e3) {3. Análisis};
\node[desc, right=of e3] (d3) {Problema y solución + citas del texto};
\node[etape, below=of e3] (e4) {4. Procedimientos};
\node[desc, right=of e4] (d4) {Oposición, enumeración, obligación};
\node[etape, below=of e4] (e5) {5. Opinión personal};
\node[desc, right=of e5] (d5) {¿Estás de acuerdo? ¿Por qué?};
\draw[fleche] (e1) -- (e2);
\draw[fleche] (e2) -- (e3);
\draw[fleche] (e3) -- (e4);
\draw[fleche] (e4) -- (e5);
\end{tikzpicture}
\end{center}
\legende{Organigramme du commentaire de texte : cinq blocs, dans l'ordre, de la présentation à l'opinion personnelle.}
\end{popfigure}

\courssub{Étape 1 — Présenter le texte}

La présentation tient en deux ou trois phrases. Elle répond à quatre questions : \emph{quoi} (type de texte), \emph{de quoi} (thème), \emph{qui} (auteur ou narrateur), \emph{dans quelles circonstances} (situation).

\begin{exemplebox}[Présenter le texte « Comprar menos, comprar mejor »]
\esp{Este texto es un texto argumentativo sobre el consumo responsable. El narrador es un joven que observa a su familia: su madre compra productos locales en el mercado, mientras que su hermano se deja llevar por la publicidad.}

« Ce texte est un texte argumentatif sur la consommation responsable. Le narrateur est un jeune qui observe sa famille : sa mère achète des produits locaux au marché, tandis que son frère se laisse influencer par la publicité. »
\end{exemplebox}

Vocabulaire utile pour le type de texte :

\begin{center}
\begin{tabularx}{\linewidth}{|X|X|X|}
\hline
\rowcolor{popblueL} \textbf{Español} & \textbf{Français} & \textbf{Remarque} \\
\hline
\esp{un texto argumentativo} & un texte argumentatif & l'auteur défend une thèse \\
\hline
\esp{un texto narrativo} & un texte narratif & on raconte des événements \\
\hline
\esp{un texto descriptivo} & un texte descriptif & on décrit une personne, un lieu \\
\hline
\esp{un artículo de periódico} & un article de journal & presse, actualité \\
\hline
\esp{un diálogo} & un dialogue & deux personnages ou plus parlent \\
\hline
\end{tabularx}
\end{center}

\courssub{Étape 2 — Dégager l'idée principale}

L'idée principale est souvent résumée par le \textbf{titre} du texte ou par sa dernière phrase. Pour notre texte, elle tient en une phrase : \esp{Debemos « comprar menos, comprar mejor »} (« Nous devons "acheter moins, acheter mieux" »).

Formule-la avec les tournures : \esp{La idea principal es que...} / \esp{El autor quiere demostrar que...} / \esp{El texto defiende la idea de que...}

\begin{attention}[Attention au verbe \esp{tratar de}]
\esp{El texto trata de...} signifie « Le texte traite de... / parle de... ». Ne traduis jamais \esp{tratar de} par « traiter » au sens de « soigner » ou « négocier » : c'est ici un semi-faux ami. Et n'oublie pas la préposition \esp{de} : on écrit \esp{trata de la importancia}, jamais \esp{trata la importancia}.
\end{attention}

\courssub{Étape 3 — Analyser les idées secondaires : problème et solution}

Un texte argumentatif s'organise presque toujours autour d'un \textbf{problème} et d'une \textbf{solution}. C'est la colonne vertébrale de ton analyse.

\begin{center}
\begin{tabularx}{\linewidth}{|X|X|}
\hline
\rowcolor{popblueL} \textbf{El problema} & \textbf{La solución} \\
\hline
\esp{el derroche} (le gaspillage) & \esp{el ahorro} (l'économie, l'épargne) \\
\hline
\esp{la publicidad} (la publicité) & \esp{comprar menos, comprar mejor} \\
\hline
\esp{los productos importados} & \esp{los productos locales} \\
\hline
\esp{el consumismo} (le consumérisme) & \esp{el comercio justo} (le commerce équitable) \\
\hline
\end{tabularx}
\end{center}

Dans la rédaction, chaque idée suit le schéma : \textbf{idée + citation + explication}.

\begin{exemplebox}[Analyser avec une citation]
\esp{El narrador denuncia el derroche: « Mi hermano compra cosas que no necesita ». Esta frase muestra que la publicidad empuja a los jóvenes a gastar su dinero sin pensar.}

« Le narrateur dénonce le gaspillage : "Mon frère achète des choses dont il n'a pas besoin". Cette phrase montre que la publicité pousse les jeunes à dépenser leur argent sans réfléchir. »
\end{exemplebox}

\courssub{Étape 4 — Relever les procédés de l'auteur}

Le correcteur attend que tu montres \emph{comment} l'auteur convainc. Trois procédés à repérer dans notre texte :

\begin{itemize}
\item \textbf{L'opposition} : \esp{en cambio} (« en revanche »), \esp{sin embargo} (« cependant »), \esp{mientras que} (« tandis que »). Elle oppose les deux comportements : la mère économe contre le frère dépensier.
\item \textbf{L'énumération} : les listes (\esp{productos locales, comercio justo, ahorro}) donnent du rythme et accumulent les arguments.
\item \textbf{Les tournures d'obligation et d'interdiction} : \esp{debemos}, \esp{tenemos que}, \esp{no se debe}, et les impératifs (\esp{compra}, \esp{elige}). Elles transforment le lecteur en acteur : on ne lui raconte pas, on lui \emph{prescrit}.
\end{itemize}

\begin{exemplebox}[Commenter un procédé]
\esp{El autor utiliza la oposición « en cambio » para comparar a la madre y al hermano: así, el lector ve claramente la diferencia entre el consumo responsable y el derroche.}

« L'auteur utilise l'opposition "en revanche" pour comparer la mère et le frère : ainsi, le lecteur voit clairement la différence entre la consommation responsable et le gaspillage. »
\end{exemplebox}

\courssub{Étape 5 — Donner son avis argumenté}

Termine toujours par toi : \esp{¿Estás de acuerdo? ¿Por qué?} Deux ou trois phrases suffisent, mais l'argument doit être \emph{concret}. Le contexte camerounais est un excellent appui : marchés locaux, produits importés, publicité à la télévision et sur les téléphones.

\begin{exemplebox}[Opinions traduites]
\esp{El texto trata de la importancia del consumo responsable.} « Le texte traite de l'importance de la consommation responsable. »

\esp{En mi opinión, el derroche es un problema grave.} « À mon avis, le gaspillage est un problème grave. »

\esp{Estoy de acuerdo porque comprar productos locales ayuda a los campesinos de mi región.} « Je suis d'accord parce qu'acheter des produits locaux aide les paysans de ma région. »

\esp{No estoy de acuerdo del todo: los productos importados a veces son más baratos para las familias pobres.} « Je ne suis pas tout à fait d'accord : les produits importés sont parfois moins chers pour les familles pauvres. »
\end{exemplebox}

\courssub{Boîte à expressions du commentaire}

\begin{aretenir}[Les expressions à connaître par cœur]
\begin{center}
\begin{tabularx}{\linewidth}{|X|X|}
\hline
\rowcolor{popblueL} \textbf{Español} & \textbf{Français} \\
\hline
\esp{El texto trata de...} & Le texte traite de... \\
\hline
\esp{Según el narrador...} & Selon le narrateur... \\
\hline
\esp{La idea principal es que...} & L'idée principale est que... \\
\hline
\esp{El autor utiliza... para convencer al lector} & L'auteur utilise... pour convaincre le lecteur \\
\hline
\esp{Esta frase muestra que...} & Cette phrase montre que... \\
\hline
\esp{En mi opinión...} & À mon avis... \\
\hline
\esp{Estoy de acuerdo porque...} & Je suis d'accord parce que... \\
\hline
\esp{No estoy de acuerdo del todo} & Je ne suis pas tout à fait d'accord \\
\hline
\end{tabularx}
\end{center}
\end{aretenir}

\courssub{Un commentaire modèle}

\begin{textebox}[Comentario modelo — « Comprar menos, comprar mejor »]
Este texto es un texto argumentativo sobre el consumo responsable. El narrador compara a su madre, que compra productos locales, con su hermano, que se deja llevar por la publicidad. La idea principal es que debemos « comprar menos, comprar mejor ». El autor utiliza la oposición (« en cambio ») y las obligaciones (« debemos », « no se debe ») para convencer al lector. En mi opinión, este texto es muy actual: en Camerún, como en España, el consumo responsable ayuda a la economía local y protege el medio ambiente.
\end{textebox}

\textbf{Traduction} : « Ce texte est un texte argumentatif sur la consommation responsable. Le narrateur compare sa mère, qui achète des produits locaux, à son frère, qui se laisse influencer par la publicité. L'idée principale est que nous devons "acheter moins, acheter mieux". L'auteur utilise l'opposition ("en revanche") et les tournures d'obligation ("nous devons", "on ne doit pas") pour convaincre le lecteur. À mon avis, ce texte est très actuel : au Cameroun comme en Espagne, la consommation responsable aide l'économie locale et protège l'environnement. »

Observe la structure du modèle : \textbf{phrase 1} = présentation ; \textbf{phrase 2} = situation et personnages ; \textbf{phrase 3} = idée principale ; \textbf{phrase 4} = procédés ; \textbf{phrase 5} = opinion personnelle. Cinq phrases, trois parties : \esp{presentación → análisis → opinión personal}. C'est le plan type à reproduire.

\begin{exoresolu}[Rédiger le début d'un commentaire]
\textbf{Énoncé.} À partir du texte « Comprar menos, comprar mejor », rédige en espagnol : 1) la présentation du texte ; 2) l'idée principale ; 3) une analyse du problème avec une citation.
\tcblower
\textbf{Solution détaillée.}

1) \textbf{Présentation} : \esp{Este texto es un texto argumentativo sobre el consumo responsable. El narrador es un joven que habla de su familia.} (type + thème + narrateur : les trois éléments obligatoires sont présents).

2) \textbf{Idée principale} : \esp{La idea principal es que debemos comprar menos y comprar mejor.} (une seule phrase, introduite par \esp{la idea principal es que}).

3) \textbf{Analyse avec citation} : \esp{El narrador denuncia el problema del derroche: su hermano « se deja llevar por la publicidad ». Esta expresión significa que el hermano compra sin pensar, solo porque ve anuncios.} (idée + citation entre guillemets + explication : le schéma complet est respecté).

\textbf{À vérifier dans ta copie} : les guillemets autour de la citation, la préposition \esp{de} après \esp{trata}, et l'accent sur \esp{opinión}, \esp{análisis}, \esp{está}.
\end{exoresolu}

\begin{exercice}[À ton tour]
Rédige en espagnol un commentaire complet (cinq à huit phrases) du texte « Comprar menos, comprar mejor », en suivant le plan type : présentation, idée principale, analyse (problème, solution, un procédé cité), opinion personnelle avec un exemple camerounais.
\end{exercice}

\begin{corrige}[Exercice — éléments de correction]
Un bon commentaire contient : une phrase de présentation avec \esp{texto argumentativo} et \esp{trata de} ; l'idée principale (\esp{debemos comprar menos, comprar mejor}) ; une analyse qui cite le texte entre guillemets (par exemple \esp{« en cambio »} ou \esp{« no se debe »}) et explique le procédé ; une opinion personnelle introduite par \esp{En mi opinión} ou \esp{Estoy de acuerdo porque}, appuyée sur un exemple local (le marché de Yaoundé, les produits du village, la publicité à la télévision). Pénalise-toi pour : résumé sans avis, citation non expliquée, \esp{trata} sans \esp{de}, accents oubliés.
\end{corrige}

\begin{aretenir}[L'essentiel de la méthode]
Un commentaire = \textbf{présentation} (type, thème, narrateur) + \textbf{idée principale} (une phrase) + \textbf{analyse} (problème, solution, procédés, chaque idée prouvée par une citation espagnole expliquée) + \textbf{opinion personnelle argumentée}. Jamais de résumé pur, jamais d'idée sans citation, jamais de citation sans explication. Structure type : \esp{presentación → análisis → opinión personal}.
\end{aretenir}

\coursec{Production guidée}

Dans la section précédente, nous avons appris à \cle{commenter un texte} : identifier le thème, relever les idées principales, dégager l'intention de l'auteur. Il est temps maintenant de passer de la lecture à l'\cle{écriture} : c'est vous qui allez produire des textes en espagnol. Cette section vous guide pas à pas dans trois tâches concrètes : rédiger des \cle{règles d'éthique professionnelle}, écrire un \cle{paragraphe d'opinion} et créer une \cle{affiche de sensibilisation}. Vous allez réutiliser tout ce que vous avez appris dans la leçon : le lexique de la consommation (\esp{consumo, consumidor, derroche, ahorro, comercio justo, producto local, publicidad, honestidad}) et les structures d'\cle{obligation} et d'\cle{interdiction} (\esp{deber, tener que, no se debe}).

\courssub{Méthode : écrire un paragraphe en espagnol}

\begin{methode}[Écrire un paragraphe en trois étapes]
Un paragraphe d'opinion de 6 à 8 lignes se construit toujours de la même manière :
\begin{enumerate}
\item \textbf{Phrase d'introduction du sujet} : annoncez le thème en une phrase simple. Exemple : \esp{Hoy en día, el consumo responsable es muy importante.} (Aujourd'hui, la consommation responsable est très importante.)
\item \textbf{Deux ou trois arguments avec des connecteurs} : développez vos idées en les reliant avec \esp{además} (de plus), \esp{también} (aussi), \esp{sin embargo} (cependant), \esp{por eso} (c'est pourquoi), \esp{por ejemplo} (par exemple).
\item \textbf{Phrase de conclusion avec votre avis} : terminez par votre opinion personnelle avec \esp{en mi opinión} (à mon avis), \esp{yo pienso que} (je pense que), \esp{para mí} (pour moi).
\end{enumerate}
\end{methode}

\begin{exemplebox}[Un paragraphe modèle]
Observez ce paragraphe et repérez les trois étapes :

\esp{Hoy en día, muchas personas compran demasiado y derrochan. Sin embargo, yo intento ser un consumidor responsable. Siempre compro productos locales en el mercado de mi barrio porque son frescos y baratos. Además, evito el derroche de agua y de comida en casa. También prefiero el comercio justo cuando es posible. En mi opinión, el ahorro es una buena costumbre para la familia y para el país.}

Traduction : « Aujourd'hui, beaucoup de personnes achètent trop et gaspillent. Cependant, j'essaie d'être un consommateur responsable. J'achète toujours des produits locaux au marché de mon quartier parce qu'ils sont frais et bon marché. De plus, j'évite le gaspillage d'eau et de nourriture à la maison. Je préfère aussi le commerce équitable quand c'est possible. À mon avis, l'épargne est une bonne habitude pour la famille et pour le pays. »

Comptez les mots du lexique : \esp{derrochan, consumidor responsable, productos locales, derroche, comercio justo, ahorro} — six mots, la consigne en demandait cinq !
\end{exemplebox}

\courssub{Tâche 1 : Rédiger 5 règles d'éthique professionnelle}

\textbf{Consigne :} \esp{Redacta cinco reglas de ética profesional para un comerciante camerunés. Utiliza \emph{deber}, \emph{tener que} y \emph{no se debe}.} (Rédigez cinq règles d'éthique professionnelle pour un commerçant camerounais. Utilisez \emph{deber}, \emph{tener que} et \emph{no se debe}.)

\begin{propriete}[Rappel : les trois structures]
\begin{itemize}
\item \esp{deber + infinitivo} : devoir (obligation morale, règle) → \esp{El vendedor debe ser honesto.} (Le vendeur doit être honnête.)
\item \esp{tener que + infinitivo} : devoir (obligation pratique, nécessité) → \esp{El comerciante tiene que abrir la tienda a las ocho.} (Le commerçant doit ouvrir la boutique à huit heures.)
\item \esp{no se debe + infinitivo} : on ne doit pas (interdiction générale, impersonnelle) → \esp{No se debe mentir.} (On ne doit pas mentir.)
\end{itemize}
\end{propriete}

\begin{exemplebox}[Exemples traduits]
\begin{itemize}
\item \esp{Un buen vendedor debe ser honesto y puntual.} = Un bon vendeur doit être honnête et ponctuel.
\item \esp{No se debe engañar al cliente.} = On ne doit pas tromper le client.
\item \esp{El comerciante tiene que respetar a los clientes.} = Le commerçant doit respecter les clients.
\item \esp{Un buen profesional debe cumplir sus promesas.} = Un bon professionnel doit tenir ses promesses.
\item \esp{No se debe vender productos de mala calidad.} = On ne doit pas vendre des produits de mauvaise qualité.
\end{itemize}
\end{exemplebox}

\begin{attention}[Pièges des francophones]
\begin{itemize}
\item Après \esp{deber} et \esp{tener que}, le verbe est toujours à l'\textbf{infinitif} : on écrit \esp{debe ser} et non \esp{debe es}.
\item \esp{No se debe} est impersonnel : ne traduisez pas mot à mot « on » par \esp{uno}. « On ne doit pas tricher » = \esp{No se debe engañar}.
\item Attention au verbe \esp{engañar} : il signifie « tromper », pas « gagner » (gagner = \esp{ganar}). C'est un \cle{faux ami} classique !
\end{itemize}
\end{attention}

\begin{exoresolu}[Tâche 1 résolue : les cinq règles]
Énoncé : rédiger cinq règles d'éthique professionnelle pour un commerçant camerounais.
\tcblower
Solution détaillée :
\begin{enumerate}
\item \esp{Un buen comerciante debe ser honesto con sus clientes.} (Un bon commerçant doit être honnête avec ses clients.)
\item \esp{El comerciante tiene que abrir su tienda a tiempo todos los días.} (Le commerçant doit ouvrir sa boutique à l'heure tous les jours.)
\item \esp{Debe vender productos de buena calidad.} (Il doit vendre des produits de bonne qualité.)
\item \esp{No se debe engañar al cliente con el peso o el precio.} (On ne doit pas tromper le client sur le poids ou le prix.)
\item \esp{Un buen profesional tiene que respetar a todo el mundo.} (Un bon professionnel doit respecter tout le monde.)
\end{enumerate}
Vérification : les trois structures sont utilisées, les verbes sont à l'infinitif, le vocabulaire de la leçon est présent (\esp{honesto, cliente, respetar, engañar, productos, calidad}).
\end{exoresolu}

\courssub{Tâche 2 : « ¿Eres un consumidor responsable ? »}

\textbf{Consigne :} \esp{Escribe un párrafo de seis a ocho líneas. Utiliza al menos cinco palabras del léxico de la lección.} (Écrivez un paragraphe de six à huit lignes. Utilisez au moins cinq mots du lexique de la leçon.)

\begin{methode}[Grille d'aide : amorces de phrases]
Si vous bloquez, commencez par ces amorces et complétez-les :
\begin{itemize}
\item \esp{Un buen profesional debe...} (Un bon professionnel doit...)
\item \esp{No se debe...} (On ne doit pas...)
\item \esp{Yo siempre compro...} (J'achète toujours...)
\item \esp{En mi familia, evitamos el derroche de...} (Dans ma famille, nous évitons le gaspillage de...)
\item \esp{En mi opinión, el consumo responsable es...} (À mon avis, la consommation responsable est...)
\item \esp{Prefiero los productos locales porque...} (Je préfère les produits locaux parce que...)
\end{itemize}
\end{methode}

\begin{center}
\begin{tabularx}{\linewidth}{|X|X|X|}
\hline
\rowcolor{popblueL} Français & Español & Remarque \\
\hline
J'achète toujours & Yo siempre compro & l'adverbe se place avant le verbe \\
\hline
à mon avis & en mi opinión & ne pas dire \esp{en mi aviso} (faux ami : \esp{aviso} = avertissement) \\
\hline
parce que & porque & une seule orthographe pour la cause \\
\hline
cependant & sin embargo & connecteur de contraste \\
\hline
de plus & además & avec accent sur la dernière syllabe \\
\hline
le gaspillage & el derroche & ne pas confondre avec \esp{derrota} (défaite) \\
\hline
\end{tabularx}
\end{center}

\begin{exoresolu}[Tâche 2 résolue : un paragraphe type]
Énoncé : rédiger un paragraphe de 6 à 8 lignes sur « ¿Eres un consumidor responsable ? » avec au moins cinq mots du lexique.
\tcblower
Solution détaillée :

\esp{¿Soy un consumidor responsable? Yo pienso que sí. En mi familia, siempre compramos productos locales en el mercado de Yaoundé porque son frescos y ayudan a los campesinos. Además, evitamos el derroche: no tiramos la comida y apagamos la luz. También intento ahorrar un poco de dinero cada mes. Sin embargo, a veces compro cosas inútiles por la publicidad. En mi opinión, el ahorro y el comercio justo son importantes para el futuro de Camerún.}

Traduction : « Suis-je un consommateur responsable ? Je pense que oui. Dans ma famille, nous achetons toujours des produits locaux au marché de Yaoundé parce qu'ils sont frais et qu'ils aident les paysans. De plus, nous évitons le gaspillage : nous ne jetons pas la nourriture et nous éteignons la lumière. J'essaie aussi d'économiser un peu d'argent chaque mois. Cependant, j'achète parfois des choses inutiles à cause de la publicité. À mon avis, l'épargne et le commerce équitable sont importants pour l'avenir du Cameroun. »

Mots du lexique utilisés : \esp{consumidor responsable, productos locales, derroche, ahorrar, publicidad, ahorro, comercio justo} — sept mots, consigne respectée.
\end{exoresolu}

\courssub{Tâche 3 : Créer une affiche de sensibilisation}

\textbf{Consigne :} \esp{Crea un cartel de sensibilización con un eslogan y tres consejos.} (Créez une affiche de sensibilisation avec un slogan et trois conseils.)

\begin{propriete}[L'impératif dans les slogans]
Dans un slogan, on utilise souvent l'\textbf{impératif} pour s'adresser directement au public :
\begin{itemize}
\item Forme affirmative (avec \esp{tú}) : comme la 3\textsuperscript{e} personne du singulier du présent → \esp{¡Ahorra agua!} (Économise l'eau !), \esp{¡Compra local!} (Achète local !), \esp{¡Consume productos cameruneses!} (Consomme des produits camerounais !)
\item Forme négative : \esp{no + subjonctif} → \esp{¡No derroches!} (Ne gaspille pas !), \esp{¡No tires la comida!} (Ne jette pas la nourriture !)
\end{itemize}
Attention aux formes : \esp{ahorra} (impératif affirmatif) mais \esp{no derroches} (subjonctif après \esp{no}).
\end{propriete}

\begin{attention}[Erreur fréquente]
Ne dites jamais \esp{¡No derrocha!} : l'impératif négatif exige le subjonctif. Comparez :
\begin{itemize}
\item \esp{¡Derrocha menos!} → affirmatif : forme simple.
\item \esp{¡No derroches!} → négatif : subjonctif (\esp{derroches}).
\end{itemize}
Autre exemple : \esp{¡Compra local!} mais \esp{¡No compres demasiado!} (N'achète pas trop !).
\end{attention}

\begin{exemplebox}[Exemples de slogans]
\begin{itemize}
\item \esp{¡No derroches, ahorra!} = Ne gaspille pas, économise !
\item \esp{¡Compra local, vive mejor!} = Achète local, vis mieux !
\item \esp{¡Consume menos, vive más!} = Consomme moins, vis plus !
\item \esp{¡Respeta al cliente!} = Respecte le client !
\end{itemize}
\end{exemplebox}

\begin{popfigure}
\begin{tikzpicture}
\draw[fill=popcream, draw=popdark, thick, rounded corners=4pt] (0,0) rectangle (10,9);
\draw[fill=poporange, rounded corners=3pt] (0.5,7.4) rectangle (9.5,8.6);
\node[text=white, font=\bfseries\large, align=center] at (5,8) {¡No derroches, ahorra!};
\draw[fill=popblueL, rounded corners=3pt] (0.7,5.4) rectangle (9.3,6.6);
\node[align=center] at (5,6) {1. Apaga la luz y el grifo.};
\draw[fill=popgreenL, rounded corners=3pt] (0.7,4.0) rectangle (9.3,5.2);
\node[align=center] at (5,4.6) {2. Compra productos locales.};
\draw[fill=poppurpleL, rounded corners=3pt] (0.7,2.6) rectangle (9.3,3.8);
\node[align=center] at (5,3.2) {3. No tires la comida.};
\draw[fill=popgoldL, draw=popgold, thick] (3.5,0.5) rectangle (6.5,2.2);
\node[align=center, font=\small] at (5,1.35) {Dibuja aquí\\una tienda o\\un mercado};
\end{tikzpicture}
\legende{Structure d'une affiche de sensibilisation : le slogan en haut, les trois conseils au centre, le dessin en bas.}
\end{popfigure}

\courssub{Correction collective d'une production type}

\begin{textebox}[Una producción de alumno para corregir]
\esp{Un buen comerciante debe es honesto con los clientes. No se debe engañar a la gente. Yo siempre compro en el mercado porque los productos locales son buenos. En mi familia no derrochamos la comida. Además, yo ahorro dinero para mis estudios. En mi opinión, el consumo responsable es muy importante para Camerún.}

Glossaire : \esp{la gente} = les gens ; \esp{ahorrar} = économiser ; \esp{los estudios} = les études.

Questions de correction : 1) Repérez l'erreur de grammaire dans la première phrase. 2) Comptez les mots du lexique de la leçon. 3) La consigne « 6 à 8 lignes » est-elle respectée ?
\end{textebox}

\begin{corrige}[Correction collective]
1) Erreur : \esp{debe es honesto} → après \esp{deber}, il faut l'infinitif : \esp{debe ser honesto}. C'est l'erreur typique du francophone qui conjugue le deuxième verbe. 2) Mots du lexique : \esp{comerciante, honesto, clientes, engañar, productos locales, derrochamos, ahorro, consumo responsable} — huit mots, excellent. 3) Oui, le paragraphe fait six lignes. Note globale : très bonne production, une seule faute à corriger.
\end{corrige}

\courssub{Critères d'évaluation}

\begin{center}
\begin{tabularx}{\linewidth}{|X|X|}
\hline
\rowcolor{popblueL} Critère & Ce que le professeur regarde \\
\hline
Respect de la consigne & nombre de lignes, nombre de mots du lexique, structures demandées \\
\hline
Exactitude grammaticale & \esp{deber/tener que + infinitivo}, impératif correct, accords \\
\hline
Richesse du lexique & variété des mots de la leçon, pas de répétition \\
\hline
Cohérence & connecteurs logiques (\esp{además, sin embargo, por eso}), plan en trois étapes \\
\hline
\end{tabularx}
\end{center}

\begin{savaistu}[La production écrite au Bac]
Au Baccalauréat, la production écrite est notée sur quatre critères : la \textbf{tâche} (avez-vous répondu à la consigne ?), la \textbf{grammaire}, le \textbf{lexique} et la \textbf{cohérence}. Le secret des bons élèves : mieux vaut des phrases simples et correctes que des phrases longues et fautives. Une phrase comme \esp{Yo compro productos locales porque son frescos} rapporte plus de points qu'une phrase compliquée pleine d'erreurs. \cle{Écrivez simple, écrivez juste !}
\end{savaistu}

\begin{experience}[Atelier d'écriture en classe]
Par groupes de quatre : 1) chaque groupe rédige ses cinq règles d'éthique pour un métier choisi (taximan, coiffeur, vendeuse au marché) ; 2) chaque élève écrit individuellement son paragraphe « ¿Eres un consumidor responsable ? » ; 3) le groupe conçoit une affiche avec slogan et trois conseils, puis la présente à la classe. Les autres élèves évaluent avec la grille des critères.
\end{experience}

\begin{exercice}[À vous de jouer]
1) Rédigez cinq règles d'éthique professionnelle pour un chauffeur de bendskin à Douala, avec \esp{deber}, \esp{tener que} et \esp{no se debe}. 2) Écrivez un paragraphe de 6 à 8 lignes : « ¿Eres un consumidor responsable ? » avec au moins cinq mots du lexique. 3) Inventez un slogan avec un impératif affirmatif et un impératif négatif.
\end{exercice}

\newpage

\coursec{Activité d'intégration}

\courssub{Objectifs de l'activité}

À la fin de cette activité d'intégration, l'élève sera capable de :

\begin{itemize}
\item mobiliser le \cle{lexique} de la consommation responsable (\esp{consumo, consumidor, publicidad, derroche, ahorro, producto local, comercio justo}) et de l'éthique professionnelle (\esp{honestidad, responsabilidad, puntualidad, respeto al cliente}) dans une production authentique ;
\item employer correctement les structures d'\cle{obligation} (\esp{deber + infinitivo}, \esp{tener que + infinitivo}) et d'\cle{interdiction} (\esp{no se debe + infinitivo}) ;
\item rédiger en groupe un document professionnel simple : un \esp{código de ética} (code d'éthique) ;
\item présenter oralement une affiche, répondre à des questions et justifier ses choix ;
\item écouter les autres groupes, poser des questions pertinentes et voter en argumentant.
\end{itemize}

\courssub{Situation}

Le quartier Mvog-Ada à Yaoundé compte de nombreux petits commerces : boutiques, ateliers de couture, restaurants, cybercafés. Pour lutter contre le gaspillage et améliorer les relations avec les clients, l'association des jeunes commerçants du quartier organise un concours : chaque commerce doit afficher son \esp{código de ética} et ses conseils pour une consommation responsable. Le meilleur code sera publié dans le bulletin de l'association. Votre classe participe au concours !

Voici l'annonce officielle du concours, telle qu'elle a été affichée au marché :

\begin{textebox}[Anuncio del concurso — Asociación de Jóvenes Comerciantes de Mvog-Ada]
¡Atención, comerciantes del barrio!

Nuestra asociación organiza un concurso: « El código de ética de mi tienda ».

Cada comercio debe presentar un código de ética con seis reglas profesionales y tres consejos para los clientes. Las reglas deben hablar de honestidad, puntualidad, respeto y responsabilidad. Los consejos deben ayudar a los clientes a consumir mejor: comprar productos locales, evitar el derroche, ahorrar y apoyar el comercio justo.

El código ganador será publicado en nuestro boletín mensual y el comercio recibirá un diploma de « Comercio Ejemplar del Año ».

Fecha límite: el viernes a las 16 horas. ¡Buena suerte a todos!
\end{textebox}

\textbf{Traduction de l'annonce :} « Attention, commerçants du quartier ! Notre association organise un concours : "Le code d'éthique de ma boutique". Chaque commerce doit présenter un code d'éthique avec six règles professionnelles et trois conseils pour les clients. Les règles doivent parler d'honnêteté, de ponctualité, de respect et de responsabilité. Les conseils doivent aider les clients à mieux consommer : acheter des produits locaux, éviter le gaspillage, économiser et soutenir le commerce équitable. Le code gagnant sera publié dans notre bulletin mensuel et le commerce recevra un diplôme de "Commerce exemplaire de l'année". Date limite : vendredi à 16 heures. Bonne chance à tous ! »

\begin{definition}[Glossaire de l'annonce]
\begin{itemize}
\item \esp{el barrio} : le quartier ; \esp{el concurso} : le concours.
\item \esp{la regla} : la règle ; \esp{el consejo} : le conseil.
\item \esp{evitar} : éviter ; \esp{apoyar} : soutenir.
\item \esp{ganador / ganadora} : gagnant(e) ; \esp{el boletín mensual} : le bulletin mensuel.
\item \esp{la fecha límite} : la date limite ; \esp{¡Buena suerte!} : bonne chance !
\end{itemize}
\end{definition}

\courssub{Consignes}

Travail en groupes de 4 élèves. Durée totale : environ 1 heure (30 minutes de préparation, 30 minutes de présentations et de vote).

\begin{enumerate}
\item \textbf{Choisir le commerce.} En groupe, choisissez un petit commerce camerounais : une boutique d'alimentation, un atelier de couture, un restaurant, un cybercafé, un salon de coiffure. Donnez-lui un nom en espagnol (par exemple : \esp{Tienda « La Confianza »}).
\item \textbf{Relire l'annonce.} Soulignez dans le texte du concours les mots du lexique de la leçon : \esp{honestidad, puntualidad, respeto, responsabilidad, productos locales, derroche, ahorrar, comercio justo}. Vérifiez que vous comprenez chaque mot.
\item \textbf{Rédiger les six règles professionnelles.} Écrivez trois règles d'obligation avec \esp{deber + infinitivo} ou \esp{tener que + infinitivo} (exemple : \esp{Tenemos que ser puntuales.}) et trois règles d'interdiction avec \esp{no se debe + infinitivo} (exemple : \esp{No se debe engañar al cliente.}). Chaque règle doit utiliser au moins un mot du lexique de la leçon.
\item \textbf{Rédiger les trois conseils aux clients.} Écrivez trois conseils pour consommer de manière responsable, avec \esp{deber}, \esp{tener que} ou \esp{no se debe} (exemple : \esp{El consumidor debe comprar productos locales.}).
\item \textbf{Préparer l'affiche.} Présentez votre code sur une grande feuille : le nom du commerce, les six règles, les trois conseils. Soignez l'orthographe et les accents.
\item \textbf{Présenter à l'oral.} Chaque membre du groupe lit une partie du code. Les autres groupes posent au moins une question avec \esp{¿Por qué... ?} (exemple : \esp{¿Por qué no se deben vender productos de mala calidad?}). Répondez avec \esp{porque...}.
\item \textbf{Voter.} La classe vote pour le code le plus convaincant. Chaque électeur justifie son vote en une phrase : \esp{Voto por el código de... porque...}.
\end{enumerate}

\begin{methode}[Comment rédiger une règle d'éthique en quatre étapes]
\begin{enumerate}
\item \textbf{Choisir l'idée} : honnêteté, ponctualité, respect, responsabilité, lutte contre le gaspillage...
\item \textbf{Choisir la structure} : \esp{deber} pour un devoir moral, \esp{tener que} pour une obligation concrète (horaire, tâche), \esp{no se debe} pour une interdiction générale.
\item \textbf{Construire la phrase} : sujet + verbe conjugué + \textbf{infinitif} + complément. Exemple : \esp{Nosotros + debemos + respetar + a los clientes.}
\item \textbf{Vérifier} : l'infinitif après \esp{deber / tener que / no se debe}, les accents, la concordance de l'adjectif (\esp{puntuales} avec \esp{nosotros}).
\end{enumerate}
\end{methode}

\begin{attention}[Pièges à éviter]
\begin{itemize}
\item Après \esp{deber}, \esp{tener que} et \esp{no se debe}, le verbe est toujours à l'\cle{infinitif} : \esp{Debemos respetar} et non \esp{debemos respetamos}.
\item N'oubliez pas \esp{que} dans \esp{tener que} : \esp{Tenemos que abrir a las ocho}.
\item \esp{No se debe} est impersonnel quand il est suivi d'un infinitif seul : on dit \esp{No se debe engañar}. En revanche, si un nom pluriel suit, l'accord est possible : \esp{No se deben vender productos de mala calidad} (« il ne faut pas vendre des produits de mauvaise qualité »). Au niveau Première, la forme simple \esp{no se debe + infinitivo} suffit et est toujours correcte.
\item Attention au faux ami : \esp{asistir al cliente} signifie « servir / recevoir le client », pas « assister » au sens d'aider ; pour « aider », utilisez \esp{ayudar al cliente}.
\item Ne confondez pas \esp{ahorrar} (économiser de l'argent, de l'eau) et \esp{salvar} (sauver une personne).
\end{itemize}
\end{attention}

\courssub{Ressources à mobiliser}

\begin{aretenir}[Structures d'obligation et d'interdiction]
\begin{itemize}
\item \esp{deber + infinitivo} : devoir (règle morale, conseil) — \esp{El comerciante debe ser honesto.} « Le commerçant doit être honnête. »
\item \esp{tener que + infinitivo} : devoir (obligation concrète) — \esp{Tenemos que abrir la tienda a las ocho.} « Nous devons ouvrir la boutique à huit heures. »
\item \esp{no se debe + infinitivo} : il ne faut pas (interdiction générale) — \esp{No se debe derrochar la comida.} « Il ne faut pas gaspiller la nourriture. »
\end{itemize}
\end{aretenir}

\begin{center}
\begin{tabularx}{\linewidth}{|X|X|}
\hline
\rowcolor{popblueL} \textbf{La consommation} & \textbf{L'éthique professionnelle} \\
\hline
el consumo, el consumidor & la honestidad, ser honesto \\
\hline
la publicidad & la responsabilidad, ser responsable \\
\hline
el derroche, derrochar & la puntualidad, ser puntual \\
\hline
el ahorro, ahorrar & el respeto al cliente, respetar \\
\hline
el producto local & engañar (tromper), la calidad \\
\hline
el comercio justo & el precio justo, la confianza \\
\hline
\end{tabularx}
\end{center}

\begin{propriete}[Rappel : conjugaison de deber et tener au présent de l'indicatif]
\begin{center}
\begin{tabular}{|l|l|l|}
\hline
\rowcolor{popgreenL} \textbf{Personne} & \textbf{deber} & \textbf{tener} \\
\hline
yo & debo & tengo \\
\hline
tú & debes & tienes \\
\hline
él, ella, usted & debe & tiene \\
\hline
nosotros, nosotras & debemos & tenemos \\
\hline
vosotros, vosotras & debéis & tenéis \\
\hline
ellos, ellas, ustedes & deben & tienen \\
\hline
\end{tabular}
\end{center}
\esp{deber} est un verbe régulier en \esp{-er} ; \esp{tener} est irrégulier (\esp{tengo, tienes, tiene...}) : c'est le même radical que dans \esp{tener que + infinitivo}.
\end{propriete}

\begin{savaistu}[Le commerce équitable et le cacao camerounais]
Le \esp{comercio justo} (commerce équitable) est un système d'échange qui garantit un \esp{precio justo} aux producteurs, de bonnes conditions de travail et le respect de l'environnement. Le Cameroun est un grand producteur de \esp{cacao} : quand un consommateur achète un chocolat issu du commerce équitable, il aide les \esp{campesinos} (paysans) à vivre dignement de leur travail. En espagnol, on dit aussi \esp{consumo responsable} ou \esp{consumo ético} pour parler d'une consommation qui pense aux autres et à la planète.
\end{savaistu}

\courssub{Proposition de production et corrigé commenté}

\begin{exoresolu}[Modèle de production : le code d'éthique d'un restaurant]
Énoncé : rédigez le \esp{código de ética} complet d'un petit restaurant de quartier, avec six règles professionnelles (3 obligations, 3 interdictions) et trois conseils aux clients.
\tcblower
\textbf{Production modèle :}

\esp{Restaurante « El Buen Sabor » — Mvog-Ada, Yaoundé}

\esp{Código de ética}

\esp{Nuestras reglas profesionales:}

\esp{1. Debemos ser honestos con los clientes: los precios son claros y justos.}

\esp{2. Tenemos que ser puntuales: abrimos el restaurante todos los días a las once.}

\esp{3. Debemos respetar a cada cliente y escuchar sus opiniones.}

\esp{4. No se debe vender comida de mala calidad.}

\esp{5. No se debe engañar al consumidor con publicidad falsa.}

\esp{6. No se debe derrochar el agua ni la electricidad.}

\esp{Consejos para nuestros clientes:}

\esp{1. El consumidor debe comprar productos locales: ¡nuestro ndolé y nuestro pescado son de la región!}

\esp{2. No se debe pedir más comida de la necesaria: hay que evitar el derroche.}

\esp{3. Usted tiene que apoyar el comercio justo: así ayuda a los campesinos cameruneses.}

\textbf{Commentaire des choix de langue :}
\begin{itemize}
\item Les deux verbes d'obligation sont variés : \esp{debemos} (règles 1 et 3, devoir moral) et \esp{tenemos que} (règle 2, obligation concrète d'horaire). Cela montre la maîtrise des deux structures.
\item Les trois interdictions utilisent la forme impersonnelle \esp{no se debe} + infinitif (\esp{vender, engañar, derrochar}), correctement construite.
\item Le lexique de la leçon est bien mobilisé : \esp{honestos, puntuales, respetar, calidad, publicidad, derroche, consumidor, productos locales, comercio justo}.
\item Les conseils varient les personnes : \esp{el consumidor debe} (3\textsuperscript{e} personne), \esp{no se debe} (impersonnel), \esp{usted tiene que} (2\textsuperscript{e} personne de politesse) — ce qui rend le texte vivant.
\item Le contexte camerounais est intégré naturellement (\esp{ndolé, campesinos cameruneses, Mvog-Ada}), ce qui rend la production authentique.
\item Orthographe : attention aux accents (\esp{atención, región, electricidad}) et à la coordination négative correcte \esp{ni... ni} dans la règle 6.
\end{itemize}
\end{exoresolu}

\begin{experience}[Présentation et vote]
Chaque groupe dispose de trois minutes pour présenter son affiche. Le professeur veille à ce que chaque membre parle. Pendant les présentations, les auditeurs remplissent une grille simple : le code est-il clair ? les structures \esp{deber / tener que / no se debe} sont-elles correctes ? le lexique de la leçon est-il utilisé ? Après les questions (\esp{¿Por qué no se debe... ?}), la classe vote à main levée et le groupe gagnant reçoit le titre de \esp{« Comercio Ejemplar de la Clase »}.
\end{experience}

\begin{exercice}[Pour aller plus loin : thème]
Traduisez en espagnol :
\begin{enumerate}
\item Le commerçant doit respecter ses clients.
\item Nous devons éviter le gaspillage de l'eau.
\item Il ne faut pas tromper le consommateur avec une publicité fausse.
\item Vous (politesse) devez acheter des produits locaux.
\end{enumerate}
\end{exercice}

\begin{corrige}[Exercice « Pour aller plus loin : thème »]
\begin{enumerate}
\item \esp{El comerciante debe respetar a sus clientes.} — Attention à la \emph{a} personnelle devant \esp{sus clientes}.
\item \esp{Tenemos que evitar el derroche del agua.} — \esp{tener que} + infinitif ; contraction \esp{de + el = del}.
\item \esp{No se debe engañar al consumidor con publicidad falsa.} — forme impersonnelle ; \esp{a + el = al}.
\item \esp{Usted debe comprar productos locales.} (ou \esp{Usted tiene que comprar productos locales.}) — accord de l'adjectif : \esp{productos locales}.
\end{enumerate}
\end{corrige}

\courssub{Bilan de l'activité}

Cette activité d'intégration a permis de réunir toutes les compétences de la leçon dans une tâche finale authentique et motivante :

\begin{itemize}
\item \textbf{Compréhension de l'écrit :} lire et exploiter l'annonce du concours, identifier les mots-clés du lexique.
\item \textbf{Lexique :} réemployer dans un contexte nouveau les huit mots du champ de la consommation et les quatre notions d'éthique professionnelle.
\item \textbf{Grammaire :} produire correctement des phrases avec \esp{deber + infinitivo}, \esp{tener que + infinitivo} et \esp{no se debe + infinitivo}, en distinguant obligation morale, obligation concrète et interdiction générale.
\item \textbf{Expression écrite :} rédiger un document professionnel réel (un code d'éthique), avec un destinataire précis (les clients) et une fonction claire (informer et engager).
\item \textbf{Expression orale et interaction :} présenter, écouter, poser des questions avec \esp{¿Por qué... ?}, justifier avec \esp{porque...}, voter en argumentant.
\item \textbf{Dimension citoyenne :} réfléchir concrètement, dans le contexte camerounais, à l'honnêteté commerciale, au respect du client, à la lutte contre le gaspillage et au soutien aux produits locaux et au commerce équitable.
\end{itemize}

\begin{aretenir}[Ce que je dois retenir]
Pour exprimer une règle professionnelle en espagnol, j'utilise \esp{deber} ou \esp{tener que} + infinitif ; pour une interdiction générale, \esp{no se debe} + infinitif. Un bon code d'éthique parle d'\esp{honestidad}, de \esp{puntualidad}, de \esp{respeto} et de \esp{responsabilidad} ; un bon consommateur évite le \esp{derroche}, préfère les \esp{productos locales} et soutient le \esp{comercio justo}.
\end{aretenir}

\newpage

\coursec{Méthodes et exercices résolus}

\courssub{Méthode 1 : Lire et comprendre un texte sur la consommation responsable}

\begin{methode}[Comprendre un texte de presse sur le consumo responsable]
Pour réussir les questions de compréhension à l'examen, suis ces étapes dans l'ordre :
\begin{enumerate}
\item \textbf{Lecture globale silencieuse.} Lis le texte deux fois sans dictionnaire. Repère le thème (ici : \esp{comprar menos, comprar mejor}), le type de texte (article d'opinion, reportage) et la thèse de l'auteur.
\item \textbf{Repérage des mots-clés du champ lexical.} Souligne les termes de la leçon : \esp{consumo, consumidor, publicidad, derroche, ahorro, producto local, comercio justo}. Ils portent le sens du texte.
\item \textbf{Localisation des réponses.} Pour chaque question, repère le paragraphe concerné : les réponses se trouvent presque toujours dans l'ordre du texte.
\item \textbf{Réponse avec ses propres mots.} Ne recopie jamais une phrase entière du texte : reformule en espagnol simple, avec un verbe conjugué correctement.
\item \textbf{Justification.} Quand la consigne demande \esp{justifica tu respuesta}, cite un court extrait entre guillemets.
\item \textbf{Vérification.} Relis : accord du verbe avec le sujet, accents, ponctuation espagnole (¿...? ¡...!).
\end{enumerate}
Structures utiles : \esp{Según el texto...} (selon le texte), \esp{El autor afirma que...} (l'auteur affirme que), \esp{Se trata de...} (il s'agit de).
\end{methode}

\begin{exoresolu}[Compréhension de l'écrit : « Comprar menos, comprar mejor »]
\textbf{Énoncé.} Lis l'extrait suivant, puis réponds en espagnol avec des phrases complètes.

\begin{center}
\emph{« En los mercados de Yaoundé, cada vez más jóvenes prefieren comprar productos locales en lugar de productos importados. Comprar un aguacate del pueblo o tela de pagne del barrio es también una forma de consumo responsable: ayuda a las familias, reduce el derroche y respeta el medio ambiente. Sin embargo, la publicidad empuja a muchos consumidores a comprar lo que no necesitan. »}
\end{center}

\begin{enumerate}
\item ¿Qué prefieren comprar los jóvenes de Yaoundé?
\item ¿Por qué comprar productos locales es una forma de consumo responsable? Cita dos razones.
\item ¿Qué efecto negativo tiene la publicidad, según el texto?
\end{enumerate}
\tcblower
\textbf{Raisonnement.} Question 1 : la réponse est dans la première phrase (\esp{prefieren comprar productos locales}). Question 2 : le texte donne trois raisons (\esp{ayuda a las familias, reduce el derroche, respeta el medio ambiente}) ; on en choisit deux et on reformule. Question 3 : dernière phrase (\esp{empuja a... comprar lo que no necesitan}).

\textbf{Réponses rédigées.}
\begin{enumerate}
\item \esp{Los jóvenes de Yaoundé prefieren comprar productos locales en lugar de productos importados.} (Ils préfèrent acheter des produits locaux plutôt que des produits importés.)
\item \esp{Según el texto, comprar productos locales es responsable porque ayuda a las familias y reduce el derroche.} (Selon le texte, acheter des produits locaux est responsable parce que cela aide les familles et réduit le gaspillage.) — \esp{porque} introduit la cause ; on reformule sans recopier toute la phrase.
\item \esp{La publicidad empuja a muchos consumidores a comprar lo que no necesitan.} (La publicité pousse de nombreux consommateurs à acheter ce dont ils n'ont pas besoin.) — Attention à l'orthographe : \esp{publicidad} se termine en \emph{-dad}, sans accent.
\end{enumerate}
\textbf{Conclusion.} Trois réponses courtes, personnelles, avec des verbes conjugués correctement : c'est exactement ce que le correcteur attend.
\end{exoresolu}

\courssub{Méthode 2 : Distinguer consommation responsable et gaspillage}

\begin{methode}[Classer et argumenter : consumo responsable / derroche]
Pour distinguer les deux notions et les utiliser dans une production :
\begin{enumerate}
\item \textbf{Définis chaque notion en une phrase.} \esp{El consumo responsable} = acheter ce dont on a besoin, local, durable, solidaire. \esp{El derroche} = acheter trop, jeter, gaspiller l'argent et les ressources.
\item \textbf{Construis deux colonnes mentales.} À gauche les mots positifs : \esp{ahorro, producto local, comercio justo, reciclar, necesidad}. À droite les mots négatifs : \esp{derroche, gastar, tirar, publicidad engañosa, moda}.
\item \textbf{Utilise les connecteurs d'opposition.} \esp{sin embargo} (cependant), \esp{en cambio} (en revanche), \esp{al contrario} (au contraire), \esp{mientras que} (tandis que).
\item \textbf{Donne un exemple concret} du contexte camerounais ou hispanophone : marché, téléphone, nourriture, eau.
\item \textbf{Conclus par un jugement personnel} avec \esp{en mi opinión}, \esp{pienso que}, \esp{me parece que}.
\end{enumerate}
Structures types : \esp{El consumo responsable consiste en + infinitif}, \esp{Derrochar es + infinitif}.
\end{methode}

\begin{exoresolu}[Expression écrite : opposition consumo responsable / derroche]
\textbf{Énoncé.} Rédige un court paragraphe (4 à 6 phrases) en espagnol : \emph{« Distingue el consumo responsable del derroche y da un ejemplo de tu entorno. »}
\tcblower
\textbf{Raisonnement.} On applique la méthode : définition des deux notions, connecteur d'opposition, exemple camerounais (le marché, l'eau, la nourriture), opinion personnelle. Temps utilisé : présent de l'indicatif, car on parle d'habitudes et de vérités générales.

\textbf{Production modèle.}

\esp{El consumo responsable consiste en comprar solo lo que necesitamos y preferir los productos locales. El derroche, en cambio, es gastar el dinero en cosas inútiles y tirar lo que todavía sirve. Por ejemplo, en mi barrio de Douala, algunas familias tiran comida todos los días, mientras que otras compran en el mercado solo lo necesario. En mi opinión, ahorrar y respetar los productos locales es una forma de ayudar a la sociedad.}

Traduction : « La consommation responsable consiste à acheter seulement ce dont nous avons besoin et à préférer les produits locaux. Le gaspillage, en revanche, c'est dépenser l'argent en choses inutiles et jeter ce qui sert encore. Par exemple, dans mon quartier de Douala, certaines familles jettent de la nourriture tous les jours, tandis que d'autres n'achètent au marché que le nécessaire. À mon avis, économiser et respecter les produits locaux est une façon d'aider la société. »

\textbf{Points grammaticaux.} \esp{consistir en + infinitif} ; \esp{en cambio} et \esp{mientras que} pour opposer ; \esp{lo que necesitamos} (ce dont nous avons besoin) avec subordonnée à l'indicatif car c'est un fait général ; accord \esp{cosas inútiles} (féminin pluriel).

\textbf{Conclusion.} Un paragraphe structuré (définition — opposition — exemple — opinion) obtient la totalité des points de contenu.
\end{exoresolu}

\courssub{Méthode 3 : Exprimer l'obligation et l'interdiction (éthique professionnelle)}

\begin{methode}[Deber, tener que, no se debe : les trois degrés de la règle]
Pour formuler des règles d'éthique professionnelle (honnêteté, ponctualité, respect du client) :
\begin{enumerate}
\item \textbf{Obligation personnelle :} \esp{tener que + infinitif} (devoir, contrainte concrète). \esp{El vendedor tiene que ser honesto.} — conjugué selon la personne : \esp{tengo que, tienes que, tiene que, tenemos que, tenéis que, tienen que}.
\item \textbf{Devoir moral ou général :} \esp{deber + infinitif}. \esp{Un profesional debe respetar al cliente.} (Un professionnel doit respecter le client.)
\item \textbf{Règle impersonnelle :} \esp{se debe / se deben + infinitif}. \esp{Se debe llegar puntual.} (On doit arriver à l'heure.) Avec un nom pluriel, le verbe s'accorde : \esp{Se deben respetar las normas.} (On doit respecter les règles.)
\item \textbf{Interdiction :} \esp{no se debe + infinitif} ou \esp{estar prohibido + infinitif}. \esp{No se debe mentir al cliente.} (Il ne faut pas mentir au client.) \esp{Está prohibido fumar aquí.} (Il est interdit de fumer ici.)
\item \textbf{Nécessité générale :} \esp{hay que + infinitif} (il faut, impersonnel, invariable). \esp{Hay que ser puntual.} (Il faut être ponctuel.)
\end{enumerate}
Réflexe : après \esp{deber}, \esp{tener que}, \esp{hay que}, le verbe est TOUJOURS à l'infinitif.
\end{methode}

\begin{exoresolu}[Thème : traduire des règles d'éthique professionnelle]
\textbf{Énoncé.} Traduis en espagnol :
\begin{enumerate}
\item Un employé doit être honnête avec les clients.
\item Il faut arriver à l'heure au travail.
\item Il ne faut pas gaspiller l'argent de l'entreprise.
\item On doit respecter les règles de l'entreprise.
\end{enumerate}
\tcblower
\textbf{Raisonnement.} Phrase 1 : devoir moral + sujet précis $\rightarrow$ \esp{deber} conjugué. Phrase 2 : « il faut » impersonnel $\rightarrow$ \esp{hay que}. Phrase 3 : interdiction impersonnelle $\rightarrow$ \esp{no se debe}. Phrase 4 : règle générale avec nom pluriel $\rightarrow$ \esp{se deben} (accord avec \esp{las reglas}).

\textbf{Solutions.}
\begin{enumerate}
\item \esp{Un empleado debe ser honesto con los clientes.} — \esp{debe} (3\textsuperscript{e} personne du singulier de \esp{deber}) + infinitif \esp{ser} ; « honnête » se dit \esp{honesto} ou \esp{honrado} : le h s'écrit mais ne se prononce pas (« o-nés-to »).
\item \esp{Hay que llegar a la hora al trabajo.} — \esp{hay que} est impersonnel et invariable ; \esp{a la hora} = à l'heure ; on peut aussi dire \esp{llegar puntual al trabajo}.
\item \esp{No se debe derrochar el dinero de la empresa.} — la négation se place devant \esp{se debe} ; \esp{derrochar} = gaspiller (on peut aussi dire \esp{malgastar}).
\item \esp{Se deben respetar las reglas de la empresa.} — accord pluriel : \esp{se deben} car \esp{las reglas} est pluriel. Avec un nom singulier, on dirait \esp{se debe respetar la regla}.
\end{enumerate}
\textbf{Conclusion.} Choisir la bonne structure d'obligation selon le sujet (personnel, moral, impersonnel) est la clé de ce point de grammaire.
\end{exoresolu}

\courssub{Méthode 4 : La version (traduire de l'espagnol vers le français)}

\begin{methode}[Traduire un texte sur la consommation sans calquer]
\begin{enumerate}
\item \textbf{Comprends d'abord le sens global} de chaque phrase avant de traduire mot à mot.
\item \textbf{Repère les faux amis} : \esp{consumidor} = consommateur (correct), mais \esp{asistir} = assister à (être présent), pas « aider » ; \esp{actualmente} = en ce moment, de nos jours.
\item \textbf{Respecte les temps} : un présent espagnol se traduit par un présent français ; un \esp{pretérito imperfecto} par un imparfait.
\item \textbf{Restitue les structures d'obligation} : \esp{se debe} = « il faut / on doit » ; \esp{hay que} = « il faut ».
\item \textbf{Reformule en français naturel} : l'ordre des mots espagnol n'est pas toujours celui du français. \esp{productos locales} = « produits locaux » (ordre identique), mais \esp{el dinero de la empresa} = « l'argent de l'entreprise ».
\item \textbf{Relis ta version française} : elle doit être grammaticalement parfaite en français.
\end{enumerate}
\end{methode}

\begin{exoresolu}[Version : un texte sur le comercio justo]
\textbf{Énoncé.} Traduis en français :

\esp{El comercio justo permite a los campesinos vender sus productos a un precio digno. En Guinea Ecuatorial y en Camerún, muchos productores de cacao no ganan bastante dinero. Por eso, los consumidores responsables deben elegir productos del comercio justo. No se debe comprar siempre lo más barato.}
\tcblower
\textbf{Raisonnement.} Phrase 1 : \esp{permitir a alguien + infinitif} = « permettre à quelqu'un de » ; \esp{precio digno} = « prix décent, équitable ». Phrase 2 : \esp{productores de cacao} = « producteurs de cacao » ; \esp{no ganan bastante} = « ne gagnent pas assez ». Phrase 3 : \esp{por eso} = « c'est pourquoi » ; \esp{deben elegir} = « doivent choisir ». Phrase 4 : \esp{no se debe comprar} = « il ne faut pas acheter » ; \esp{lo más barato} = « le moins cher ».

\textbf{Traduction.}

« Le commerce équitable permet aux paysans de vendre leurs produits à un prix digne. En Guinée équatoriale et au Cameroun, beaucoup de producteurs de cacao ne gagnent pas assez d'argent. C'est pourquoi les consommateurs responsables doivent choisir des produits du commerce équitable. Il ne faut pas toujours acheter le moins cher. »

\textbf{Points de vigilance.} \esp{comercio justo} se traduit par l'expression consacrée « commerce équitable » (pas « commerce juste ») ; \esp{por eso} = « c'est pourquoi » ; \esp{Camerún} prend un accent en espagnol, mais « Cameroun » n'en prend pas en français.

\textbf{Conclusion.} Une bonne version est fidèle au sens ET naturelle en français : les deux critères comptent dans la notation.
\end{exoresolu}

\courssub{Méthode 5 : Rédiger un commentaire de texte sur la consommation}

\begin{methode}[Le commentaire de texte en trois parties]
Pour commenter un texte comme « Comprar menos, comprar mejor » :
\begin{enumerate}
\item \textbf{Introduction (2-3 phrases).} Présente le texte : type, thème, problématique. \esp{Este texto es un artículo de opinión que trata del consumo responsable. El autor plantea la siguiente pregunta: ¿cómo podemos consumir mejor?}
\item \textbf{Développement (2 paragraphes).} Paragraphe 1 : les idées principales du texte, reformulées et illustrées. Paragraphe 2 : ton avis argumenté, avec \esp{en mi opinión}, \esp{estoy de acuerdo porque...}, \esp{sin embargo...}
\item \textbf{Conclusion (1-2 phrases).} Bilan + ouverture. \esp{En conclusión, consumir de forma responsable es un deber de todos.}
\item \textbf{Connecteurs à mobiliser.} \esp{en primer lugar, además, sin embargo, por eso, en conclusión}.
\item \textbf{Langue.} Présent de l'indicatif pour les généralités ; phrases courtes ; vocabulaire de la leçon (\esp{ahorro, derroche, comercio justo, honestidad}).
\end{enumerate}
\end{methode}

\begin{exoresolu}[Commentaire guidé : introduction et conclusion]
\textbf{Énoncé.} Rédige en espagnol l'introduction et la conclusion d'un commentaire du texte « Comprar menos, comprar mejor », qui défend la consommation responsable et critique la publicité.
\tcblower
\textbf{Raisonnement.} Introduction : type de texte + thème + problématique. Conclusion : bilan + ouverture. On utilise le présent, les connecteurs, et on évite de recopier le titre mot pour mot dans la problématique.

\textbf{Introduction modèle.}

\esp{Este texto es un artículo de opinión titulado « Comprar menos, comprar mejor ». Trata del consumo responsable y critica el derroche provocado por la publicidad. El autor plantea una pregunta importante: ¿podemos cambiar nuestra manera de comprar para ayudar a la sociedad?}

(« Ce texte est un article d'opinion intitulé "Acheter moins, acheter mieux". Il traite de la consommation responsable et critique le gaspillage provoqué par la publicité. L'auteur pose une question importante : pouvons-nous changer notre façon d'acheter pour aider la société ? »)

\textbf{Conclusion modèle.}

\esp{En conclusión, el texto nos enseña que consumir de forma responsable es un deber de todos los ciudadanos. Comprar productos locales y evitar el derroche no solo ahorra dinero, sino que también construye un mundo más justo.}

(« En conclusion, le texte nous enseigne que consommer de façon responsable est un devoir de tous les citoyens. Acheter des produits locaux et éviter le gaspillage n'économise pas seulement de l'argent, mais construit aussi un monde plus juste. »)

\textbf{Points grammaticaux.} \esp{tratar de} = traiter de ; \esp{no solo... sino que también} = non seulement... mais aussi ; \esp{enseñar que + indicatif} car on présente un fait ; \esp{para ayudar} : \esp{para} + infinitif exprime le but.

\textbf{Conclusion.} Une introduction qui annonce clairement le thème et une conclusion qui ouvre sur une idée générale donnent immédiatement une bonne impression au correcteur.
\end{exoresolu}

\begin{attention}[Les 5 erreurs qui coûtent des points à l'examen]
\begin{enumerate}
\item \textbf{Fautes d'accent.} \esp{publicidad}, \esp{honestidad}, \esp{consumidor} s'écrivent SANS accent (mots terminant en \emph{-d} ou \emph{-r}, accentués sur la dernière syllabe : c'est la règle normale). En revanche : \esp{ética, económico, política, Camerún} prennent un accent. Un accent oublié ou mal placé est pénalisé à chaque occurrence.
\item \textbf{Faux amis.} \esp{asistir} = assister à (pas « aider » : \esp{ayudar}) ; \esp{actualmente} = en ce moment (pas « en fait ») ; \esp{una librería} = une librairie (pas une bibliothèque : \esp{biblioteca}) ; \esp{éxito} = succès (pas « sortie » : \esp{salida}).
\item \textbf{Calques du français.} « Il faut » se traduit par \esp{hay que + infinitif} ou \esp{se debe + infinitif}. « Être à l'heure » = \esp{llegar a la hora} ou \esp{ser puntual}, pas \emph{estar en la hora}.
\item \textbf{Accords.} Adjectif accordé avec le nom : \esp{productos locales}, \esp{las reglas de la empresa}, \esp{una profesional honesta}. Avec \esp{se debe(n)}, accorde le verbe avec le nom qui suit : \esp{se deben respetar las normas}.
\item \textbf{Choix du temps.} Pour des règles, des habitudes et des vérités générales (l'éthique, la consommation), utilise le \textbf{présent de l'indicatif}, pas le futur ni le conditionnel. Après \esp{deber / tener que / hay que}, toujours l'\textbf{infinitif}, jamais un verbe conjugué.
\end{enumerate}
\end{attention}

\newpage

\coursec{Exercices}

\courssub{Je vérifie mes connaissances}

\begin{exercice}[QCM : la consommation responsable]
Entoure la bonne réponse.

\begin{enumerate}
\item Un \esp{consumidor responsable} est quelqu'un qui :
\begin{enumerate}
\item compra muchos productos baratos.
\item compra pensando en el medio ambiente y en los demás.
\item no compra nunca nada.
\end{enumerate}
\item El \esp{comercio justo} garantiza :
\begin{enumerate}
\item precios muy bajos para el consumidor.
\item un salario digno para los productores.
\item productos importados de Europa.
\end{enumerate}
\item \esp{Derrochar} significa :
\begin{enumerate}
\item ahorrar dinero.
\item gastar sin necesidad.
\item comprar productos locales.
\end{enumerate}
\item \esp{No se debe tirar la comida} expresa :
\begin{enumerate}
\item un consejo.
\item una interdicción.
\item una descripción.
\end{enumerate}
\end{enumerate}
\end{exercice}

\begin{exercice}[Vrai ou faux]
Lis les phrases suivantes et indique si elles sont \esp{verdaderas} (V) ou \esp{falsas} (F) selon le texte « Comprar menos, comprar mejor » étudié en classe. Justifie chaque réponse par une courte citation du texte.

\begin{enumerate}
\item \esp{Según el texto, la publicidad ayuda a ahorrar.}
\item \esp{Comprar productos locales beneficia a los campesinos de la región.}
\item \esp{El texto dice que el derroche es un problema solamente de los ricos.}
\item \esp{El ahorro es una forma de consumo responsable.}
\end{enumerate}
\end{exercice}

\begin{exercice}[Vocabulaire : appariement]
Relie chaque mot espagnol à sa traduction française.

\begin{center}
\begin{tabularx}{\linewidth}{|X|X|}
\hline
\rowcolor{popblueL} \textbf{Español} & \textbf{Français} \\
\hline
1. consumo & a. l'épargne \\
\hline
2. derroche & b. l'honnêteté \\
\hline
3. ahorro & c. le gaspillage \\
\hline
4. honestidad & d. la consommation \\
\hline
5. puntualidad & e. le commerce équitable \\
\hline
6. comercio justo & f. la ponctualité \\
\hline
\end{tabularx}
\end{center}
\end{exercice}

\begin{exercice}[Conjugaison : deber et tener que]
Complète les phrases avec le verbe entre parenthèses, conjugué à la bonne personne.

\begin{enumerate}
\item Yo \ldots\ldots\ldots\ldots (deber) comprar productos locales.
\item Nosotros \ldots\ldots\ldots\ldots (tener que) respetar a los clientes.
\item ¿Tú \ldots\ldots\ldots\ldots (deber) apagar la luz al salir?
\item Los vendedores \ldots\ldots\ldots\ldots (tener que) ser honestos.
\item Ella \ldots\ldots\ldots\ldots (deber) llegar puntual al trabajo.
\end{enumerate}
\end{exercice}

\courssub{Je m'entraîne}

\begin{exercice}[Texte à trous : obligation et interdiction]
Complète les phrases avec \esp{deber}, \esp{tener que}, \esp{hay que} ou \esp{no se debe} (attention à la conjugaison et au sens).

\begin{enumerate}
\item Un profesional \ldots\ldots\ldots\ldots respetar a sus clientes.
\item \ldots\ldots\ldots\ldots gaspiller la nourriture : ¡hay mucha gente con hambre!
\item Para ahorrar, \ldots\ldots\ldots\ldots comparar los precios antes de comprar.
\item Los empleados \ldots\ldots\ldots\ldots llegar a tiempo todos los días.
\item \ldots\ldots\ldots\ldots mentir en la publicidad.
\item Yo \ldots\ldots\ldots\ldots estudiar esta noche porque mañana tengo examen.
\item \ldots\ldots\ldots\ldots comprar productos de temporada : son más baratos.
\item Un comerciante honesto \ldots\ldots\ldots\ldots dar siempre el cambio correcto.
\item \ldots\ldots\ldots\ldots tirar los plásticos al río.
\item Nosotros \ldots\ldots\ldots\ldots consumir menos agua durante la estación seca.
\end{enumerate}
\end{exercice}

\begin{exercice}[Thème : traduction en espagnol]
Traduis en espagnol les phrases suivantes.

\begin{enumerate}
\item Un professionnel doit respecter ses clients.
\item On ne doit pas gaspiller la nourriture.
\item Nous devons acheter des produits locaux au marché de Yaoundé.
\item Il faut économiser l'eau et l'électricité.
\item Les vendeurs doivent être honnêtes et ponctuels.
\end{enumerate}
\end{exercice}

\begin{exercice}[Transformation de phrases]
Transforme chaque phrase selon la consigne entre parenthèses.

\begin{enumerate}
\item \esp{Debo apagar la luz.} (mets à la forme impersonnelle avec \esp{hay que})
\item \esp{Tenemos que respetar a los clientes.} (mets à la forme négative avec \esp{no se debe} + infinitif, en gardant le sens : « ne pas manquer de respect »)
\item \esp{No se debe derrochar el agua.} (transforme en obligation affirmative avec \esp{hay que} : « économiser »)
\item \esp{El vendedor tiene que ser honesto.} (remplace \esp{tener que} par \esp{deber})
\item \esp{Hay que comprar productos locales.} (mets à la première personne du pluriel avec \esp{tener que})
\end{enumerate}
\end{exercice}

\begin{exercice}[Version : traduction en français]
Traduis en français le paragraphe suivant.

\begin{textebox}[La publicidad y el derroche]
La publicidad está en todas partes: en la televisión, en la radio, en Internet y hasta en las calles de Douala. Muchas veces nos empuja a comprar cosas que no necesitamos. Así, algunas familias gastan su dinero en productos inútiles y después no pueden pagar la escuela de sus hijos. Para evitar este derroche, hay que pensar antes de comprar y hacer una lista de lo necesario. El consumidor inteligente no compra más, compra mejor.
\end{textebox}
\end{exercice}

\courssub{Je me prépare au Bac}

\begin{exercice}[Compréhension de texte type Bac]
Lis le texte suivant, puis réponds aux questions \textbf{en espagnol}, avec des phrases complètes.

\begin{textebox}[El comercio justo en México]
En las montañas del estado de Chiapas, en el sur de México, miles de familias campesinas cultivan café desde hace generaciones. Durante muchos años, vendían su cosecha a intermediarios que pagaban precios muy bajos. Los campesinos trabajaban mucho pero vivían en la pobreza, y sus hijos no podían ir a la escuela.

Hace algunos años, estas familias decidieron organizarse en cooperativas de comercio justo. Ahora venden su café directamente a compradores de Europa y de América del Norte, sin intermediarios. Gracias a este sistema, reciben un precio estable y digno por su trabajo. Con el dinero ganado, las comunidades han construido escuelas y centros de salud.

El comercio justo también protege el medio ambiente: los campesinos no utilizan productos químicos peligrosos y respetan la tierra de sus abuelos. Además, los consumidores europeos saben exactamente de dónde viene el café que beben cada mañana.

Comprar productos del comercio justo cuesta un poco más, pero es una manera de consumir con responsabilidad y con honestidad. Como dice un campesino de Chiapas: « Cuando usted compra nuestro café, no compra solamente una bebida: compra la educación de nuestros niños. »
\end{textebox}

\textbf{Preguntas de comprensión} (5 points) :
\begin{enumerate}
\item ¿Dónde cultivan café las familias del texto?
\item ¿Por qué vivían en la pobreza antes de las cooperativas?
\item ¿Qué han construido las comunidades con el dinero del comercio justo?
\item ¿Cómo protege el comercio justo el medio ambiente?
\item Según el campesino citado al final, ¿qué compra realmente un consumidor de café del comercio justo?
\end{enumerate}

\textbf{Questions de langue} (5 points) :
\begin{enumerate}
\item Trouve dans le texte deux mots de la famille de \esp{comprar}.
\item Relève dans le texte une phrase qui exprime une obligation avec \esp{hay que} ou une expression équivalente, et traduis-la en français.
\item Conjugue le verbe \esp{vender} au présent de l'indicatif à toutes les personnes.
\item Traduis en français la phrase : \esp{« Comprar productos del comercio justo cuesta un poco más, pero es una manera de consumir con responsabilidad. »}
\item Donne le contraire de : \esp{precios muy bajos} ; \esp{la pobreza}.
\end{enumerate}
\end{exercice}

\begin{exercice}[Thème et version type Bac]

\textbf{A. Thème} (5 points) : traduis en espagnol.

\begin{enumerate}
\item Les consommateurs doivent acheter moins mais acheter mieux.
\item Il ne faut pas croire toute la publicité.
\item Un employé honnête doit arriver à l'heure et respecter ses clients.
\item Nous devons soutenir les producteurs locaux du Cameroun.
\item On ne doit pas jeter les plastiques dans la rue.
\end{enumerate}

\textbf{B. Version} (5 points) : traduis en français le texte suivant.

\begin{textebox}[El consumidor responsable]
Ser un consumidor responsable no es difícil. Primero, hay que comprar solamente lo necesario y evitar el derroche. Segundo, debemos preferir los productos locales, porque ayudan a los campesinos y a los comerciantes de nuestro país. Tercero, no se debe tirar la comida: millones de personas sufren hambre en el mundo. Finalmente, tenemos que reutilizar y reciclar para proteger la naturaleza. Con pequeños gestos cada día, todos podemos consumir mejor.
\end{textebox}
\end{exercice}

\begin{exercice}[Production écrite type Bac]
\textbf{Sujet} : \esp{Escribe un texto de 60 a 80 palabras sobre tu comportamiento como consumidor.}

Consignes :
\begin{itemize}
\item Tu texte doit être rédigé entièrement \textbf{en espagnol}.
\item Tu dois répondre aux questions suivantes : ¿Qué compras normalmente? ¿Compras productos locales? ¿Qué debes hacer para consumir mejor? ¿Qué no se debe hacer?
\item Utilise au moins \textbf{trois expressions d'obligation ou d'interdiction} : \esp{deber}, \esp{tener que}, \esp{hay que}, \esp{no se debe}.
\item Emploie au moins \textbf{quatre mots du lexique de la leçon} : \esp{consumo}, \esp{consumidor}, \esp{publicidad}, \esp{derroche}, \esp{ahorro}, \esp{producto local}, \esp{comercio justo}, \esp{honestidad}.
\item Soigne l'orthographe et les accents.
\end{itemize}
\end{exercice}

\newpage

\coursec{Corrigés des exercices}

\begin{corrige}[Exercice 1 — QCM : la consommation responsable]
\begin{enumerate}
\item \textbf{b.} \esp{Compra pensando en el medio ambiente y en los demás.} — Un consommateur responsable pense aux conséquences de ses achats (environnement, producteurs), il ne s'agit ni d'acheter beaucoup (a) ni de ne rien acheter (c).
\item \textbf{b.} \esp{Un salario digno para los productores.} — Le commerce équitable garantit un prix juste aux producteurs, pas forcément des prix bas pour le consommateur.
\item \textbf{b.} \esp{Gastar sin necesidad.} — \esp{Derrochar} = gaspiller, dépenser sans nécessité. C'est le contraire d'\esp{ahorrar} (économiser).
\item \textbf{b.} \esp{Una prohibición.} — La tournure impersonnelle \esp{no se debe} + infinitif exprime l'\cle{interdiction} : « on ne doit pas jeter la nourriture ». \emph{Attention : le mot espagnol est \textbf{prohibición} (et non « interdicción »)}.
\end{enumerate}
\end{corrige}

\begin{corrige}[Exercice 2 — Vrai ou faux]
\begin{enumerate}
\item \textbf{F.} Le texte dit le contraire : \esp{« la publicidad nos empuja a comprar cosas que no necesitamos »} (la publicité nous pousse à acheter, donc à dépenser, pas à économiser).
\item \textbf{V.} \esp{« Comprar productos locales beneficia a los campesinos de la región »} : acheter local fait vivre les producteurs du pays.
\item \textbf{F.} Le gaspillage concerne tout le monde : \esp{« el derroche es un problema de todos »} ; chacun, riche ou pauvre, peut gaspiller.
\item \textbf{V.} \esp{« El ahorro es una forma de consumo responsable »} : économiser, c'est consommer de manière responsable.
\end{enumerate}
\textbf{Méthode} : pour justifier, cite toujours un court fragment du texte entre guillemets, puis explique en une phrase.
\end{corrige}

\begin{corrige}[Exercice 3 — Vocabulaire : appariement]
\begin{center}
\begin{tabularx}{\linewidth}{|X|X|X|}
\hline
\rowcolor{popblueL} \textbf{Español} & \textbf{Français} & \textbf{Réponse} \\
\hline
1. consumo & d. la consommation & 1 — d \\
\hline
2. derroche & c. le gaspillage & 2 — c \\
\hline
3. ahorro & a. l'épargne & 3 — a \\
\hline
4. honestidad & b. l'honnêteté (probité) & 4 — b \\
\hline
5. puntualidad & f. la ponctualité & 5 — f \\
\hline
6. comercio justo & e. le commerce équitable & 6 — e \\
\hline
\end{tabularx}
\end{center}
\textbf{Astuce} : \esp{derroche} et \esp{ahorro} sont des contraires ; \esp{honestidad} désigne l'honnêteté morale (probité), à ne pas confondre avec la politesse (\esp{cortesía}).
\end{corrige}

\begin{corrige}[Exercice 4 — Conjugaison : \esp{deber} et \esp{tener que}]
\begin{enumerate}
\item \esp{Yo \textbf{debo} comprar productos locales.} — \esp{dever} est régulier en \cle{-er} : \esp{debo, debes, debe, debemos, debéis, deben}.
\item \esp{Nosotros \textbf{tenemos que} respetar a los clientes.} — \esp{tener} est irrégulier (\cle{e $\rightarrow$ ie}, 1\iere{} pers. \esp{tengo}) : \esp{tengo, tienes, tiene, tenemos, tenéis, tienen}.
\item \esp{¿Tú \textbf{debes} apagar la luz al salir?} — « Dois-tu éteindre la lumière en sortant ? »
\item \esp{Los vendedores \textbf{tienen que} ser honestos.} — 3\iere{} pers. du pluriel : \esp{tienen}.
\item \esp{Ella \textbf{debe} llegar puntual al trabajo.} — 3\iere{} pers. du singulier : \esp{debe}.
\end{enumerate}
\textbf{Rappel} : après \esp{deber} et \esp{tener que}, le verbe suivant est toujours à l'\textbf{infinitif}.
\end{corrige}

\begin{corrige}[Exercice 5 — Texte à trous : obligation et interdiction]
\begin{enumerate}
\item \esp{Un profesional \textbf{debe} respetar a sus clientes.} (ou \esp{tiene que})
\item \esp{\textbf{No se debe} desperdiciar la comida : ¡hay mucha gente con hambre!} — interdiction impersonnelle. (Note : « gaspiller la nourriture » = \esp{desperdiciar / tirar la comida}.)
\item \esp{Para ahorrar, \textbf{hay que} comparar los precios antes de comprar.} — obligation générale impersonnelle (\cle{hay que} + infinitif).
\item \esp{Los empleados \textbf{tienen que} llegar a tiempo todos los días.} (ou \esp{deben})
\item \esp{\textbf{No se debe} mentir en la publicidad.} — interdiction.
\item \esp{Yo \textbf{tengo que} estudiar esta noche porque mañana tengo examen.} (ou \esp{debo}) — nécessité personnelle : \esp{tener que} convient très bien.
\item \esp{\textbf{Hay que} comprar productos de temporada : son más baratos.} — conseil général.
\item \esp{Un comerciante honesto \textbf{debe} dar siempre el cambio correcto.} (ou \esp{tiene que})
\item \esp{\textbf{No se debe} tirar los plásticos al río.} — interdiction.
\item \esp{Nosotros \textbf{tenemos que} consumir menos agua durante la estación seca.} (ou \esp{debemos})
\end{enumerate}
\textbf{Points de vigilance} : \esp{hay que} et \esp{no se debe} sont \cle{impersonnels} (jamais de sujet de personne) ; \esp{deber} et \esp{tener que} se conjuguent selon le sujet.
\end{corrige}

\begin{corrige}[Exercice 6 — Thème : traduction en espagnol]
\begin{enumerate}
\item \esp{Un profesional debe respetar a sus clientes.} — attention à la \cle{préposition \esp{a}} devant le complément de personne (\esp{respetar a alguien}).
\item \esp{No se debe desperdiciar la comida.} (ou \esp{No se debe tirar la comida.})
\item \esp{Debemos comprar productos locales en el mercado de \textbf{Yaundé}.} (ou \esp{Tenemos que comprar...}) — \emph{Yaundé} s'écrit avec accent aigu sur le \emph{é} en espagnol.
\item \esp{Hay que ahorrar agua y electricidad.} — « il faut » impersonnel = \esp{hay que} ; pas de sujet.
\item \esp{Los vendedores deben ser honestos y puntuales.} — adjectifs au pluriel : \esp{honestos}, \esp{puntuales}.
\end{enumerate}
\textbf{Piège fréquent} : ne pas traduire « on ne doit pas » par \esp{no tenemos que} ; la tournure impersonnelle correcte est \esp{no se debe} + infinitif.
\end{corrige}

\begin{corrige}[Exercice 7 — Transformation de phrases]
\begin{enumerate}
\item \esp{Debo apagar la luz.} $\rightarrow$ \esp{\textbf{Hay que apagar la luz.}} — la forme impersonnelle efface la personne.
\item \esp{Tenemos que respetar a los clientes.} $\rightarrow$ \esp{\textbf{No se debe faltar al respeto a los clientes.}} — « ne pas manquer de respect » = \esp{no faltar al respeto}.
\item \esp{No se debe derrochar el agua.} $\rightarrow$ \esp{\textbf{Hay que ahorrar agua.}} — l'interdiction de gaspiller devient l'obligation d'économiser (contraire).
\item \esp{El vendedor tiene que ser honesto.} $\rightarrow$ \esp{\textbf{El vendedor debe ser honesto.}} — \esp{deber} remplace \esp{tener que} sans changer le sens (obligation).
\item \esp{Hay que comprar productos locales.} $\rightarrow$ \esp{\textbf{Tenemos que comprar productos locales.}} — 1\iere{} pers. du pluriel de \esp{tener} : \esp{tenemos}.
\end{enumerate}
\end{corrige}

\begin{corrige}[Exercice 8 — Version : traduction en français]
\textbf{Traduction proposée :}

« La publicité est partout : à la télévision, à la radio, sur Internet et jusque dans les rues de Douala. Souvent, elle nous pousse à acheter des choses dont nous n'avons pas besoin. Ainsi, certaines familles dépensent leur argent en produits inutiles et ensuite ne peuvent pas payer l'école de leurs enfants. Pour éviter ce gaspillage, il faut réfléchir avant d'acheter et faire une liste de ce qui est nécessaire. Le consommateur intelligent n'achète pas plus, il achète mieux. »

\textbf{Points de vigilance} :
\begin{itemize}
\item \esp{hasta} = « jusqu'à / jusque » (faux ami : ce n'est pas « hâte »).
\item \esp{hacer una lista de lo necesario} = « faire une liste du nécessaire » ; \esp{lo} + adjectif = « ce qui est... ».
\item \esp{no compra más, compra mejor} : rendre le contraste « n'achète pas plus, achète mieux ».
\end{itemize}
\end{corrige}

\begin{corrige}[Exercice 9 — Compréhension de texte type Bac]
\textbf{Preguntas de comprensión} (1 point par réponse, 5 points) :
\begin{enumerate}
\item \esp{Las familias cultivan café en las montañas del estado de Chiapas, en el sur de México.}
\item \esp{Vivían en la pobreza porque vendían su cosecha a intermediarios que pagaban precios muy bajos.}
\item \esp{Con el dinero del comercio justo, las comunidades han construido escuelas y centros de salud.}
\item \esp{El comercio justo protege el medio ambiente porque los campesinos no utilizan productos químicos peligrosos y respetan la tierra.}
\item \esp{Según el campesino, el consumidor no compra solamente una bebida: compra la educación de los niños.}
\end{enumerate}
\textbf{Méthode} : répondre par une phrase complète qui reprend les mots de la question ; ne jamais recopier tout un paragraphe.

\textbf{Questions de langue} (1 point par question, 5 points) :
\begin{enumerate}
\item Deux mots de la famille de \esp{comprar} : \esp{compradores} (acheteurs) et \esp{compra} (achète / achat). On accepte aussi \esp{comprar} lui-même répété dans le texte.
\item \textbf{Relever une phrase d'obligation dans le texte} (ex. \esp{hay que / debe / tienen que} + infinitif). Si le texte n'en contient pas explicitement, transformer une phrase en obligation : \esp{« Hay que comprar productos del comercio justo »} = « Il faut acheter des produits du commerce équitable ».
\item Conjugaison de \esp{vender} au présent : \esp{yo vendo, tú vendes, él/ella vende, nosotros vendemos, vosotros vendéis, ellos venden.} — verbe régulier en \cle{-er}.
\item Traduction : « Acheter des produits du commerce équitable coûte un peu plus cher, mais c'est une manière de consommer avec responsabilité. »
\item Contraires : \esp{precios muy bajos} $\rightarrow$ \esp{precios muy altos} (ou \esp{elevados}) ; \esp{la pobreza} $\rightarrow$ \esp{la riqueza}.
\end{enumerate}
\textbf{Barème indicatif} : compréhension 5 pts (1 pt par question : 0,5 pour le sens, 0,5 pour la langue) ; langue 5 pts (1 pt par question).
\end{corrige}

\begin{corrige}[Exercice 10 — Thème et version type Bac]
\textbf{A. Thème} (1 point par phrase, 5 points) :
\begin{enumerate}
\item \esp{Los consumidores deben comprar menos pero comprar mejor.}
\item \esp{No hay que creer toda la publicidad.} (ou \esp{No se debe creer...})
\item \esp{Un empleado honesto debe llegar \textbf{a tiempo} y respetar a sus clientes.} — « à l'heure / à l'heure exacte » = \esp{a tiempo} (meilleur que \esp{a la hora} pour la ponctualité habituelle).
\item \esp{Debemos apoyar a los productores locales de \textbf{Camerún}.} (ou \esp{Tenemos que apoyar...}) — \esp{apoyar a} + personnes ; \esp{Camerún} avec accent sur le \emph{u}.
\item \esp{No se debe tirar los plásticos a la calle.} (ou \esp{en la calle}) — \esp{tirar a la calle} = « jeter dans la rue ».
\end{enumerate}
\textbf{Points de vigilance} : \esp{apoyar a} + personnes ; \esp{Camerún} avec accent ; \esp{tirar a la calle} = « jeter dans la rue ».

\textbf{B. Version} (5 points) — traduction proposée :

« Être un consommateur responsable n'est pas difficile. D'abord, il faut acheter seulement le nécessaire et éviter le gaspillage. Ensuite, nous devons préférer les produits locaux, parce qu'ils aident les paysans et les commerçants de notre pays. Troisièmement, on ne doit pas jeter la nourriture : des millions de personnes souffrent de la faim dans le monde. Enfin, nous devons réutiliser et recycler pour protéger la nature. Avec de petits gestes chaque jour, nous pouvons tous mieux consommer. »

\textbf{Barème indicatif} : thème 5 pts (1 pt/phrase : vocabulaire 0,5, grammaire 0,5) ; version 5 pts (fidélité 3 pts, qualité du français 2 pts).
\end{corrige}

\begin{corrige}[Exercice 11 — Production écrite type Bac]
\textbf{Proposition de rédaction modèle} (environ 75 mots) :

\begin{textebox}[Mi comportamiento como consumidor]
Normalmente compro comida, ropa y material escolar. Intento comprar productos locales en el mercado, porque ayudan a los campesinos de mi región. Para consumir mejor, debo hacer una lista antes de ir de compras y tengo que comparar los precios. Además, hay que evitar el derroche: no se debe tirar la comida ni comprar cosas inútiles por culpa de la publicidad. Como consumidor responsable, también quiero apoyar el comercio justo y practicar el ahorro. Con honestidad y pequeños gestos, todos podemos consumir mejor.
\end{textebox}

\textbf{Traduction du modèle} : « Normalement, j'achète de la nourriture, des vêtements et du matériel scolaire. J'essaie d'acheter des produits locaux au marché, parce qu'ils aident les paysans de ma région. Pour mieux consommer, je dois faire une liste avant d'aller faire les courses et je dois comparer les prix. De plus, il faut éviter le gaspillage : on ne doit pas jeter la nourriture ni acheter des choses inutiles à cause de la publicité. Comme consommateur responsable, je veux aussi soutenir le commerce équitable et pratiquer l'épargne. Avec honnêteté et de petits gestes, nous pouvons tous mieux consommer. »

\textbf{Barème indicatif} (10 points) :
\begin{itemize}
\item Respect de la consigne et du nombre de mots (60-80) : 2 pts.
\item Contenu : réponse aux quatre questions (achats habituels, efforts, règles, engagement) : 2 pts.
\item Trois expressions d'obligation/interdiction correctement employées (\esp{debo}, \esp{tengo que}, \esp{hay que}, \esp{no se debe}) : 2 pts.
\item Quatre mots du lexique de la leçon (\esp{derroche}, \esp{ahorro}, \esp{comercio justo}, \esp{honestidad}, \esp{productos locales}, \esp{publicidad}...) : 2 pts.
\item Orthographe, accents, grammaire : 2 pts.
\end{itemize}

\textbf{Points de vigilance} :
\begin{itemize}
\item Vérifier la conjugaison de \esp{deber} (\esp{debo}, \esp{debe}) et de \esp{tener que} (\esp{tengo que}).
\item \esp{hay que} et \esp{no se debe} restent \cle{impersonnels} : jamais \esp{yo hay que}.
\item Accents obligatoires : \esp{publicidad}, \esp{región}, \esp{estación}, \esp{Yaundé}, \esp{Camerún} ; \esp{por culpa de} et non « porque de ».
\item Ne pas écrire \esp{consumar} : le verbe est \esp{consumir}, le nom \esp{el consumo}.
\end{itemize}
\end{corrige}

\newpage

\coursec{Fiche bilan}

\begin{popfigure}
\begin{tikzpicture}[mindmap, grow cyclic, every node/.style={concept, text=white, font=\small},
  root concept/.append style={concept color=popblue, font=\bfseries\small, minimum size=3.2cm},
  level 1 concept/.append style={sibling angle=60, level distance=4.2cm, minimum size=2.6cm}]
\node[root concept]{Consumo responsable y ética profesional}
  child[concept color=poporange]{ node{Texto: Comprar menos, comprar mejor} }
  child[concept color=popgreen]{ node{Lexique: consumo, ahorro, derroche} }
  child[concept color=poppurple]{ node{Comercio justo y producto local} }
  child[concept color=poppink]{ node{Ética: honestidad, puntualidad, respeto} }
  child[concept color=popgold]{ node{Obligation: deber, tener que} }
  child[concept color=popdark]{ node{Interdiction: no se debe, estar prohibido} };
\end{tikzpicture}
\legende{Carte mentale de la leçon : consommation responsable et éthique professionnelle.}
\end{popfigure}

\begin{bilan}
\textbf{1. Lexique clé à connaître par cœur}

\begin{center}
\begin{tabularx}{\linewidth}{|X|X|}
\hline
\rowcolor{popblueL} \textbf{Español} & \textbf{Français} \\
\hline
el consumo / el consumidor & la consommation / le consommateur \\
\hline
la publicidad & la publicité \\
\hline
el derroche & le gaspillage \\
\hline
el ahorro & l'épargne, les économies \\
\hline
el producto local & le produit local \\
\hline
el comercio justo & le commerce équitable \\
\hline
la honestidad & l'honnêteté \\
\hline
la puntualidad & la ponctualité \\
\hline
el respeto al cliente & le respect du client \\
\hline
la responsabilidad & la responsabilité \\
\hline
\end{tabularx}
\end{center}

\textbf{2. Grammaire : obligation et interdiction}

\begin{center}
\begin{tabularx}{\linewidth}{|X|X|X|}
\hline
\rowcolor{popgreenL} \textbf{Structure} & \textbf{Emploi} & \textbf{Exemple} \\
\hline
\esp{deber + infinitivo} & obligation morale, conseil & \esp{Debemos ahorrar.} « Nous devons économiser. » \\
\hline
\esp{tener que + infinitivo} & obligation extérieure, nécessité & \esp{Tengo que llegar a tiempo.} « Je dois arriver à l'heure. » \\
\hline
\esp{hay que + infinitivo} & obligation impersonnelle & \esp{Hay que respetar al cliente.} « Il faut respecter le client. » \\
\hline
\esp{no se debe + infinitivo} & interdiction, déconseil & \esp{No se debe derrochar.} « On ne doit pas gaspiller. » \\
\hline
\esp{estar prohibido + infinitivo} & interdiction formelle & \esp{Está prohibido mentir.} « Il est interdit de mentir. » \\
\hline
\end{tabularx}
\end{center}

\textbf{3. Idées directrices du texte}
\begin{itemize}
\item Consommation responsable : acheter moins, acheter mieux, préférer les produits locaux et le commerce équitable.
\item Gaspillage : acheter sous l'influence de la publicité, jeter, surconsommer.
\item Éthique professionnelle : honnêteté, responsabilité, ponctualité, respect du client.
\end{itemize}

\textbf{4. Méthodes express}
\begin{itemize}
\item \textbf{Compréhension} : lire deux fois, repérer le thème, souligner les mots-clés, dégager l'idée principale de chaque paragraphe.
\item \textbf{Commentaire} : introduction (auteur, thème, problématique) ; développement (idées + exemples du texte) ; conclusion (bilan + avis personnel).
\item \textbf{Traduction} : repérer d'abord la structure d'obligation française (« il faut », « on doit ») avant de choisir \esp{hay que}, \esp{deber} ou \esp{tener que}.
\end{itemize}
\end{bilan}

\begin{aretenir}[Checklist avant le Bac]
\begin{itemize}
\item[$\square$] Je connais le lexique : \esp{consumo, consumidor, publicidad, derroche, ahorro}.
\item[$\square$] Je sais traduire \esp{producto local} et \esp{comercio justo} sans faux amis.
\item[$\square$] Je cite les quatre valeurs de l'éthique professionnelle : \esp{honestidad, responsabilidad, puntualidad, respeto al cliente}.
\item[$\square$] Je distingue \esp{deber} (devoir moral) et \esp{tener que} (obligation concrète).
\item[$\square$] J'emploie \esp{hay que} pour une obligation impersonnelle.
\item[$\square$] Je forme une interdiction avec \esp{no se debe + infinitivo}.
\item[$\square$] Je conjugue \esp{tener} correctement : \esp{tengo, tienes, tiene, tenemos, tenéis, tienen}.
\item[$\square$] Je sais que l'infinitif suit toujours \esp{deber}, \esp{tener que}, \esp{hay que}.
\item[$\square$] Je peux résumer le texte en trois idées principales.
\item[$\square$] Je sais structurer un commentaire : introduction, développement, conclusion.
\item[$\square$] Je peux donner mon avis : \esp{En mi opinión... / Pienso que...}.
\item[$\square$] Je vérifie les accents : \esp{honestidad, publicidad, responsabilidad, ética}.
\end{itemize}
\end{aretenir}

\findecours

\end{document}
